\documentclass[numbers=enddot,12pt,final,onecolumn,notitlepage]{scrartcl}%
\usepackage[headsepline,footsepline,manualmark]{scrlayer-scrpage}
\usepackage{amsfonts}
\usepackage{amssymb}
\usepackage{amsmath}
\usepackage{amsthm}
\usepackage{framed}
\usepackage{comment}
\usepackage{color}
\usepackage[breaklinks=true]{hyperref}
\usepackage{lmodern} % Rendering font compatibility: unavailable native Palatino metrics.
\usepackage[T1]{fontenc}
\usepackage{needspace}
\usepackage{ytableau}
\usepackage{hyperxmp}
%TCIDATA{OutputFilter=latex2.dll}
%TCIDATA{Version=5.50.0.2960}
%TCIDATA{LastRevised=Friday, April 17, 2026 18:23:50}
%TCIDATA{SuppressPackageManagement}
%TCIDATA{<META NAME="GraphicsSave" CONTENT="32">}
%TCIDATA{<META NAME="SaveForMode" CONTENT="1">}
%TCIDATA{BibliographyScheme=Manual}
%TCIDATA{Language=American English}
%BeginMSIPreambleData
\providecommand{\U}[1]{\protect\rule{.1in}{.1in}}
%EndMSIPreambleData
\theoremstyle{definition}
\newtheorem{theo}{Theorem}[section]
\newenvironment{theorem}[1][]
{\begin{theo}[#1]\begin{leftbar}}
{\end{leftbar}\end{theo}}
\newtheorem{lem}[theo]{Lemma}
\newenvironment{lemma}[1][]
{\begin{lem}[#1]\begin{leftbar}}
{\end{leftbar}\end{lem}}
\newtheorem{prop}[theo]{Proposition}
\newenvironment{proposition}[1][]
{\begin{prop}[#1]\begin{leftbar}}
{\end{leftbar}\end{prop}}
\newtheorem{defi}[theo]{Definition}
\newenvironment{definition}[1][]
{\begin{defi}[#1]\begin{leftbar}}
{\end{leftbar}\end{defi}}
\newtheorem{remk}[theo]{Remark}
\newenvironment{remark}[1][]
{\begin{remk}[#1]\begin{leftbar}}
{\end{leftbar}\end{remk}}
\newtheorem{coro}[theo]{Corollary}
\newenvironment{corollary}[1][]
{\begin{coro}[#1]\begin{leftbar}}
{\end{leftbar}\end{coro}}
\newtheorem{conv}[theo]{Convention}
\newenvironment{convention}[1][]
{\begin{conv}[#1]\begin{leftbar}}
{\end{leftbar}\end{conv}}
\newtheorem{quest}[theo]{Question}
\newenvironment{question}[1][]
{\begin{quest}[#1]\begin{leftbar}}
{\end{leftbar}\end{quest}}
\newtheorem{warn}[theo]{Warning}
\newenvironment{conclusion}[1][]
{\begin{warn}[#1]\begin{leftbar}}
{\end{leftbar}\end{warn}}
\newtheorem{conj}[theo]{Conjecture}
\newenvironment{conjecture}[1][]
{\begin{conj}[#1]\begin{leftbar}}
{\end{leftbar}\end{conj}}
\newtheorem{exam}[theo]{Example}
\newenvironment{example}[1][]
{\begin{exam}[#1]\begin{leftbar}}
{\end{leftbar}\end{exam}}
\newcommand{\silentsection}{\section}
\newenvironment{verlong}{}{}
\newenvironment{vershort}{}{}
\newenvironment{noncompile}{}{}
\includecomment{verlong}
\excludecomment{vershort}
\excludecomment{noncompile}
\newenvironment{statement}{\begin{quote}}{\end{quote}}
\newenvironment{fineprint}{\begin{small}}{\end{small}}
\iffalse
\newenvironment{proof}[1][Proof]{\noindent\textbf{#1.} }{\ \rule{0.5em}{0.5em}}
\newenvironment{question}[1][Question]{\noindent\textbf{#1.} }{\ \rule{0.5em}{0.5em}}
\NOEXPAND{\silentsection}{\section}
\fi
\newcommand{\id}{\operatorname{id}}
\newcommand{\set}[1]{\left\{ #1 \right\}}
\newcommand{\abs}[1]{\left| #1 \right|}
\newcommand{\tup}[1]{\left( #1 \right)}
\newcommand{\ive}[1]{\left[ #1 \right]}
\newcommand{\kk}{\mathbf{k}}
\let\sumnonlimits\sum
\let\prodnonlimits\prod
\renewcommand{\sum}{\sumnonlimits\limits}
\renewcommand{\prod}{\prodnonlimits\limits}
\newcommand\arxiv[1]{\href{http://www.arxiv.org/abs/#1}{\texttt{arXiv:#1}}}
\voffset=0cm
\hoffset=-0.7cm
\setlength\textheight{22.5cm}
\setlength\textwidth{15.5cm}
\ihead{Representations of somewhere-to-below shuffles, version 2026-04-17}
\ohead{page \thepage}
\cfoot{}
\title{Specht-module spectra of arbitrary one-sided cycle shuffles}
\author{Darij Grinberg}
\date{arxiv:2508.00752v3; CC0 1.0}
\begin{document}
\maketitle
\noindent\textit{Editorial note.} Complete original mathematical document169--3994 retained, including all original supporting arguments. Native long/short proof controls[112, 113] are switched only to include verlong/exclude vershort, selecting author-provided detailed content; raw TeX/archive and logical line count unchanged. This is a version choice, not agent verification.\par\medskip
\noindent\textit{Editorial note.} Original CC0 1.0 license preserved exactly. Intro restatements, standard/special-case forms, consequences and applications are not extra counts.\par\medskip
\noindent\textit{Editorial note.} The selected theorem retains its original any-field statement, multiplicity-collision convention and generic diagonalizability condition. Characteristic-zero filtration statements used in the original proof and all omitted supporting proofs remain intact; no field hypothesis is strengthened, weakened or inferred.\par\medskip
\noindent\textit{Editorial note.} Provided original119-line s2b3.bbl is inserted verbatim once; original biblio.bib preserved. The absent custom halpha-abbrv style is not reconstructed or needed with the supplied author bbl.\par\medskip
\noindent\textit{Editorial note.} Author retains responsibility for mathematics. No proof repair, new assumptions, independent proof review, symbolic/numerical replay, counterexample search or certification.\par\medskip
\noindent\textit{Editorial note.} Actual r001 unavailable Palatino font metric replaced only in rendering preamble by existing Latin Modern. Original mathematical text, detailed-proof mode, raw source and failed r001 preserved; line locators unchanged.\par\medskip
\noindent\textit{Allowed standard prerequisites (editorial).}\begin{itemize}
\item Full original finite-field/Frobenius, noncommutative-ring, or Specht/shuffle framework as appropriate. All native definitions, coefficient conventions, field/ring hypotheses and parameter restrictions retained verbatim with cited prior facts.. Complete original detailed proof and supporting lemmas supplied, without reconstructing any external cited proof or independently certifying the arguments.
\end{itemize}\medskip
\section{Introduction}

This paper is a continuation of \cite{s2b1} with other means. Specifically,
our goal here is to answer some natural representation-theoretical questions
around the somewhere-to-below shuffles in the symmetric group algebra
(including \cite[Question 16.12]{s2b1}).

We recall that the \emph{somewhere-to-below shuffles} are $n$ special elements
$t_{1},t_{2},\ldots,t_{n}$ of the group algebra $\mathbf{k}\left[
S_{n}\right]  $ of a symmetric group $S_{n}$ over an arbitrary commutative
ring $\mathbf{k}$; they are defined by
\[
t_{\ell}:=\operatorname*{cyc}\nolimits_{\ell}+\operatorname*{cyc}%
\nolimits_{\ell,\ell+1}+\operatorname*{cyc}\nolimits_{\ell,\ell+1,\ell
+2}+\cdots+\operatorname*{cyc}\nolimits_{\ell,\ell+1,\ldots,n}\in
\mathbf{k}\left[  S_{n}\right]  ,
\]
where $\operatorname*{cyc}\nolimits_{\ell,\ell+1,\ldots,k}$ denotes the cycle
that sends $\ell\mapsto\ell+1\mapsto\ell+2\mapsto\cdots\mapsto k\mapsto\ell$
(and leaves all remaining elements of $\left[  n\right]  =\left\{
1,2,\ldots,n\right\}  $ unchanged). Together with their linear combinations
(called the \emph{one-sided cycle shuffles}), they have been introduced and
studied in the paper \cite{s2b1} (published with abridgements as
\cite{s2b1-pub}\footnote{The numbering of results in \cite{s2b1} and in
\cite{s2b1-pub} is identical except for Section 9, so the reader can consult
either version.}) by Lafreni\`{e}re and the present author. One of the main
results is \cite[Theorem 11.1]{s2b1}, which constructs a basis $\left(
a_{w}\right)  _{w\in S_{n}}$ of $\mathbf{k}\left[  S_{n}\right]  $ on which
each of the shuffles $t_{1},t_{2},\ldots,t_{n}$ acts (by right multiplication)
triangularly -- i.e., which satisfies%
\[
a_{w}t_{\ell}\in\operatorname*{span}\left\{  a_{v}\ \mid\ v\leq w\right\}
\ \ \ \ \ \ \ \ \ \ \text{for all }w\in S_{n}\text{ and }\ell\in\left\{
1,2,\ldots,n\right\}
\]
(for an appropriate total order $<$ on $S_{n}$). This entails that the
shuffles $t_{1},t_{2},\ldots,t_{n}$ and their linear combinations have integer
eigenvalues; these eigenvalues have indeed been found (\cite[\S 12]{s2b1})
along with their multiplicities (\cite[\S 13]{s2b1}). As a further
consequence, the $\mathbf{k}$-subalgebra of $\mathbf{k}\left[  S_{n}\right]  $
generated by $t_{1},t_{2},\ldots,t_{n}$ is isomorphic to an algebra of
upper-triangular matrices, and the commutators $\left[  t_{i},t_{j}\right]
:=t_{i}t_{j}-t_{j}t_{i}$ are nilpotent; a followup work \cite{s2b2} proves
even stronger claims.

However, like any elements of the group algebra $\mathbf{k}\left[
S_{n}\right]  $, the shuffles $t_{1},t_{2},\ldots,t_{n}$ act not just on the
whole algebra $\mathbf{k}\left[  S_{n}\right]  $, but on any of its modules,
i.e., on any representation of $S_{n}$. Thus, the question about eigenvalues
can be asked for each representation of $S_{n}$, in particular for the
\emph{Specht modules} (which are the irreducible representations of $S_{n}$,
at least in characteristic $0$).

The main goal of this paper is to answer this latter question. Let us give a
quick outline of the answer (which was announced in \cite[\S 11]%
{fps2024sn})\footnote{The formulation in \cite[\S 11]{fps2024sn} uses the
Frobenius characteristic map, but this has turned out to be a red herring.}:

We shall use some basic notions from the representation theory of $S_{n}$ and
from symmetric functions; the reader can find all prerequisites in
\cite[Chapters 6 and 7]{Fulton97}. For any partition $\lambda$ of $n$, a
Specht module $\mathcal{S}^{\lambda}$ is defined, which is a representation of
$S_{n}$ with a basis indexed by standard tableaux of shape $\lambda$. (In
\cite{Fulton97}, it is called $S^{\lambda}$.) This $S_{n}$-module
$\mathcal{S}^{\lambda}$ is irreducible when $\mathbf{k}$ has characteristic
$0$. Each $u\in\mathbf{k}\left[  S_{n}\right]  $ acts (on the left) on this
Specht module $\mathcal{S}^{\lambda}$; we let $L_{\lambda}\left(  u\right)  $
denote this action (viewed as a $\mathbf{k}$-module endomorphism of
$\mathcal{S}^{\lambda}$).

\begin{noncompile}
We let $\mathcal{R}$ denote the (inductive) representation ring of the
symmetric groups (called $R$ in \cite[\S 7.3]{Fulton97}), and $\Lambda$ denote
the ring of symmetric functions over $\mathbb{Z}$ (defined in \cite[\S 6.2]%
{Fulton97}). The \emph{Frobenius characteristic map} is the ring isomorphism
$\varphi:\Lambda\rightarrow\mathcal{R}$ defined in \cite[\S 7.3]{Fulton97},
and the famous \emph{Schur function} $s_{\lambda}\in\Lambda$ corresponding to
a partition $\lambda$ is the preimage of the Specht module $\mathcal{S}%
^{\lambda}$ (defined over a field of characteristic $0$) under this
isomorphism $\varphi$.
\end{noncompile}

We let $\Lambda$ denote the ring of symmetric functions over $\mathbb{Z}$
(defined in \cite[\S 6.2]{Fulton97}). We recall that it has a basis $\left(
s_{\lambda}\right)  _{\lambda\text{ is a partition}}$ of \emph{Schur
functions} $s_{\lambda}$.

For each $m\in\mathbb{N}$, we let $h_{m}\in\Lambda$ denote the $m$-th complete
homogeneous symmetric function. For each $m>1$, we let $z_{m}\in\Lambda$
denote the Schur function
\[
z_{m}:=s_{\left(  m-1,1\right)  }=h_{m-1}h_{1}-h_{m}\in\Lambda.
\]


A set of integers is called \emph{lacunar} if it contains no two consecutive
integers. For each lacunar subset $I$ of $\left[  n-1\right]  $, we define a
symmetric function%
\[
z_{I}:=h_{i_{1}-1}\prod_{j=2}^{m}z_{i_{j}-i_{j-1}}\in\Lambda,
\]
where $i_{1},i_{2},\ldots,i_{m}$ are the elements of $I\cup\left\{
n+1\right\}  $ in increasing order (so that $i_{m}=n+1$ and $I=\left\{
i_{1}<i_{2}<\cdots<i_{m-1}\right\}  $). When this symmetric function $z_{I}$
is expanded in the basis $\left(  s_{\lambda}\right)  _{\lambda\text{ is a
partition}}$ of $\Lambda$, the coefficient of a given Schur function
$s_{\lambda}$ shall be called $c_{I}^{\lambda}$. This coefficient
$c_{I}^{\lambda}$ is actually a Littlewood--Richardson coefficient (since
$z_{I}$ is a skew Schur function), hence a nonnegative integer.

We now claim the following:

\begin{theorem}
[part of Theorem \ref{thm.eigs}]\label{thm.intro.eigs}Let $\lambda$ be a
partition. Let $\omega_{1},\omega_{2},\ldots,\omega_{n}\in\mathbf{k}$. Then,
the eigenvalues of the operator $L_{\lambda}\left(  \omega_{1}t_{1}+\omega
_{2}t_{2}+\cdots+\omega_{n}t_{n}\right)  $ on the Specht module $\mathcal{S}%
^{\lambda}$ are the linear combinations%
\[
\omega_{1}m_{I,1}+\omega_{2}m_{I,2}+\cdots+\omega_{n}m_{I,n}%
\ \ \ \ \ \ \ \ \ \ \text{for }I\subseteq\left[  n-1\right]  \text{ lacunar
satisfying }c_{I}^{\lambda}\neq0,
\]
where the $m_{I,k}$ are certain nonnegative integers defined combinatorially
(namely, $m_{I,k}$ is the distance between $k$ and the smallest element of
$I\cup\left\{  n+1\right\}  $ that is $\geq k$). The algebraic multiplicities
of these eigenvalues are the $c_{I}^{\lambda}$ in the generic case (i.e., if
no two $I$'s produce the same linear combination; otherwise the multiplicities
of colliding eigenvalues should be added together). Moreover, if all these
linear combinations are distinct, then $L_{\lambda}\left(  \omega_{1}%
t_{1}+\omega_{2}t_{2}+\cdots+\omega_{n}t_{n}\right)  $ is diagonalizable.
\end{theorem}

The proof of this theorem will rely on the filtration $0=F_{0}\subseteq
F_{1}\subseteq F_{2}\subseteq\cdots\subseteq F_{f_{n+1}}=\mathbf{k}\left[
S_{n}\right]  $ of $\mathbf{k}\left[  S_{n}\right]  $ introduced in
\cite[\S 8.1]{s2b1}. We call this the \emph{Fibonacci filtration} of
$\mathbf{k}\left[  S_{n}\right]  $, as its length $f_{n+1}$ is the $\left(
n+1\right)  $-st Fibonacci number. We note that this filtration is not
completely canonical, as it depends on the choice of a listing $Q_{1}%
,Q_{2},\ldots,Q_{f_{n+1}}$ of all lacunar subsets of $\left[  n-1\right]  $ in
the order of increasing sum of elements (the ties can be broken arbitrarily,
whence the non-canonicity). Much about this filtration was already understood
in \cite{s2b1}, but we will need some additional information about the action
of $S_{n}$ on its subquotients $F_{i}/F_{i-1}$:

Let $\mathcal{A}$ be the $\mathbf{k}$-algebra $\mathbf{k}\left[  S_{n}\right]
$, and let $\mathcal{T}$ be its $\mathbf{k}$-subalgebra generated by
$t_{1},t_{2},\ldots,t_{n}$. Then, each $F_{i}$ is a left ideal of
$\mathcal{A}$ but is also fixed under right multiplication by each $t_{\ell}$;
therefore, each $F_{i}$ is an $\left(  \mathcal{A},\mathcal{T}\right)
$-subbimodule of $\mathcal{A}$. Thus, each subquotient $F_{i}/F_{i-1}$ of the
Fibonacci filtration is an $\left(  \mathcal{A},\mathcal{T}\right)
$-bimodule. As a right $\mathcal{T}$-module, it is \emph{scalar} (meaning that
each $t_{\ell}$ acts on it by a scalar, which is in fact the integer
$m_{Q_{i},\ell}$ from \cite[Theorem 8.1 \textbf{(c)}]{s2b1}). As a left
$\mathcal{A}$-module (i.e., as a representation of $S_{n}$), we describe it
explicitly here:

\begin{theorem}
[part of Theorem \ref{thm.Fi/Fi-1.as-ind}]\label{thm.intro.Fi/Fi-1}Let
$i\in\left[  f_{n+1}\right]  $.

\begin{noncompile}
Then, the isomorphism class of the $S_{n}$-representation $F_{i}/F_{i-1}$ is
$\varphi\left(  z_{Q_{i}}\right)  $. More explicitly, and true in arbitrary
characteristic, we can describe it as follows:
\end{noncompile}

Consider the lacunar subset $Q_{i}$ of $\left[  n-1\right]  $ (from the above
listing $Q_{1},Q_{2},\ldots,Q_{f_{n+1}}$). Write the set $Q_{i}\cup\left\{
n+1\right\}  $ as $\left\{  i_{1}<i_{2}<\cdots<i_{m}\right\}  $, so that
$i_{m}=n+1$. Furthermore, set $i_{0}:=1$. Set $j_{k}:=i_{k}-i_{k-1}$ for each
$k\in\left[  m\right]  $. Note that $j_{1}\geq0$ and $j_{2},j_{3},\ldots
,j_{m}>1$ and $j_{1}+j_{2}+\cdots+j_{m}=i_{m}-i_{0}=n$.

For each $p\in\mathbb{N}$, we let $\mathcal{H}_{p}$ denote the trivial
$1$-dimensional representation of $S_{p}$ (that is, the $\mathbf{k}$-module
$\mathbf{k}$ on which $S_{p}$ acts trivially), and we let $\mathcal{Z}_{p}$
denote the reflection quotient representation of $S_{p}$ (that is, the free
$\mathbf{k}$-module $\mathbf{k}^{p}$ on which $S_{p}$ acts by permuting the
coordinates, divided by the submodule consisting of all vectors of the form
$\left(  a,a,\ldots,a\right)  \in\mathbf{k}^{p}$). Then,%
\[
F_{i}/F_{i-1}\cong\operatorname*{Ind}\nolimits_{S_{j_{1}}\times S_{j_{2}%
}\times\cdots\times S_{j_{m}}}^{S_{n}}\underbrace{\left(  \mathcal{H}_{j_{1}%
}\otimes\mathcal{Z}_{j_{2}}\otimes\mathcal{Z}_{j_{3}}\otimes\cdots
\otimes\mathcal{Z}_{j_{m}}\right)  }_{\substack{\text{the first tensorand is
an }\mathcal{H}\text{,}\\\text{while all others are }\mathcal{Z}\text{'s}}}
\]
as $S_{n}$-representations. Here, we embed $S_{j_{1}}\times S_{j_{2}}%
\times\cdots\times S_{j_{m}}$ into $S_{n}$ by the usual parabolic embedding
(since $j_{1}+j_{2}+\cdots+j_{m}=n$).
\end{theorem}

This theorem will be a crucial stepping stone on our way to Theorem
\ref{thm.intro.eigs}.

We note that neither of our two main results requires any assumption about the
characteristic of $\mathbf{k}$. However, in positive characteristic, care must
be taken to distinguish between the reflection quotient representation
$\mathcal{Z}_{p}$ in Theorem \ref{thm.intro.Fi/Fi-1} and the reflection
subrepresentation $\mathcal{R}_{p}$ (which consists of the zero-sum vectors in
$\mathbf{k}^{p}$); the two representations have the same dimension $p-1$ (for
$p\geq1$), but are not isomorphic unless $\operatorname*{char}\mathbf{k}\neq
p$ or $p\leq2$.

\begin{noncompile}
It is worth noticing that the diagonalizability claim in Theorem
\ref{thm.intro.eigs} is an improvement upon the diagonalizability claim in
\cite[Theorem 12.3]{s2b1}. Indeed, if $\mathbf{k}$ is a field of
characteristic $0$, and if each partition $\lambda$ of $n$ has the property
that the scalars%
\[
\omega_{1}m_{I,1}+\omega_{2}m_{I,2}+\cdots+\omega_{n}m_{I,n}%
\ \ \ \ \ \ \ \ \ \ \text{for }I\subseteq\left[  n-1\right]  \text{ lacunar
satisfying }c_{I}^{\lambda}\neq0
\]
are distinct, then it follows that the action of $\omega_{1}t_{1}+\omega
_{2}t_{2}+\cdots+\omega_{n}t_{n}$ on each Specht module $\mathcal{S}^{\lambda
}$ is diagonalizable, and therefore so is the action of $\omega_{1}%
t_{1}+\omega_{2}t_{2}+\cdots+\omega_{n}t_{n}$ on $\mathcal{A}$ itself (since
$\mathcal{A}$ can be decomposed into a direct sum of Specht modules). In
\cite[Theorem 12.3]{s2b1}, this is proved only under the stronger condition
that the scalars%
\[
\omega_{1}m_{I,1}+\omega_{2}m_{I,2}+\cdots+\omega_{n}m_{I,n}%
\ \ \ \ \ \ \ \ \ \ \text{for }I\subseteq\left[  n-1\right]  \text{ lacunar}%
\]
are all distinct. However, I don't know of any specific cases in which this
helps (e.g., it does not help with the one-sided cycle shuffle $\sum_{k=1}%
^{n}\dfrac{1}{n+1-k}t_{k}\in\mathbb{Q}\left[  S_{n}\right]  $ for $n=12$).
\end{noncompile}

We suspect that our results can be generalized (\textquotedblleft%
$q$-deformed\textquotedblright) from the symmetric group algebra to the Hecke
algebra $\mathcal{H}_{q}\left(  S_{n}\right)  $. Most results from \cite{s2b1}
can definitely be generalized this way, as will be detailed in forthcoming work.

\paragraph{Acknowledgements}

The author would like to thank Sarah Brauner, Nadia Lafreni\`{e}re, Martin
Lorenz, Victor Reiner, Sheila Sundaram and Mark Wildon for inspiring discussions,
as well as two referees for helpful commentary
(in particular the suggestion of Example~\ref{exa.Fi/Fi-1.L3.e1}).

\begin{noncompile}
The SageMath CAS \cite{sagemath} was indispensable at every stage of the
research presented here.
\end{noncompile}


\section{Definitions and notations}

\subsection{Basics}

We recall some notations from \cite{s2b1}.

Let $\mathbf{k}$ be any commutative ring. (We don't require that $\mathbf{k}$
is a field or a $\mathbb{Q}$-algebra, but the reader can think of
$\mathbf{k}=\mathbb{Q}$ as a standing example.)

Let $\mathbb{N}:=\left\{  0,1,2,\ldots\right\}  $ be the set of all
nonnegative integers.

For any integers $a$ and $b$, we set
\[
\left[  a,b\right]  :=\left\{  k\in\mathbb{Z}\ \mid\ a\leq k\leq b\right\}
=\left\{  a,a+1,\ldots,b\right\}  .
\]
This is an empty set if $a>b$. In general, $\left[  a,b\right]  $ is called an
\emph{integer interval}.

For each $n\in\mathbb{Z}$, let $\left[  n\right]  :=\left[  1,n\right]
=\left\{  1,2,\ldots,n\right\}  $.

Fix an integer $n\in\mathbb{N}$. Let $S_{n}$ be the $n$-th symmetric group,
i.e., the group of all permutations of $\left[  n\right]  $. We multiply
permutations in the \textquotedblleft continental\textquotedblright\ way: that
is, $\left(  \pi\sigma\right)  \left(  i\right)  =\pi\left(  \sigma\left(
i\right)  \right)  $ for all $\pi,\sigma\in S_{n}$ and $i\in\left[  n\right]
$.

For any $k$ distinct elements $i_{1},i_{2},\ldots,i_{k}$ of $\left[  n\right]
$, we let $\operatorname*{cyc}\nolimits_{i_{1},i_{2},\ldots,i_{k}}$ be the
permutation in $S_{n}$ that sends $i_{1},i_{2},\ldots,i_{k-1},i_{k}$ to
$i_{2},i_{3},\ldots,i_{k},i_{1}$, respectively while leaving all remaining
elements of $\left[  n\right]  $ unchanged. This permutation is known as a
\emph{cycle}. Note that $\operatorname*{cyc}\nolimits_{i}=\operatorname*{id}$
for any single $i\in\left[  n\right]  $.

For any $i\in\left[  n-1\right]  $, we denote the cycle $\operatorname*{cyc}%
\nolimits_{i,i+1}$ by $s_{i}$ and call it a \emph{simple transposition}.

\subsection{Somewhere-to-below shuffles, $\mathcal{A}$ and $\mathcal{T}$}

Let $\mathcal{A}$ be the group algebra $\mathbf{k}\left[  S_{n}\right]  $. In
this algebra, define $n$ elements $t_{1},t_{2},\ldots,t_{n}$ by
setting\footnote{We view $S_{n}$ as a subset of $\mathbf{k}\left[
S_{n}\right]  $ in the obvious way.}%
\[
t_{\ell}:=\operatorname*{cyc}\nolimits_{\ell}+\operatorname*{cyc}%
\nolimits_{\ell,\ell+1}+\operatorname*{cyc}\nolimits_{\ell,\ell+1,\ell
+2}+\cdots+\operatorname*{cyc}\nolimits_{\ell,\ell+1,\ldots,n}\in
\mathbf{k}\left[  S_{n}\right]
\]
for each $\ell\in\left[  n\right]  $. Thus, in particular, $t_{n}%
=\operatorname*{cyc}\nolimits_{n}=\operatorname*{id}=1$ (where $1$ means the
unity of $\mathbf{k}\left[  S_{n}\right]  $). The $n$ elements $t_{1}%
,t_{2},\ldots,t_{n}$ are known as the \emph{somewhere-to-below shuffles}.

We let $\mathcal{T}$ be the $\mathbf{k}$-subalgebra of $\mathcal{A}$ generated
by these $n$ somewhere-to-below shuffles $t_{1},t_{2},\ldots,t_{n}$. Clearly,
$\mathcal{A}$ is an $\left(  \mathcal{A},\mathcal{T}\right)  $-bimodule (with
$\mathcal{A}$ acting from the left by multiplication, and $\mathcal{T}$ acting
from the right by multiplication).

\subsection{Some $S_{n}$-representation theory}

We recall that the representations of the symmetric group $S_{n}$ (over
$\mathbf{k}$) are precisely the left $\mathbf{k}\left[  S_{n}\right]
$-modules, i.e., the left $\mathcal{A}$-modules. We will use the following
four classes of $S_{n}$-representations in particular:

\begin{enumerate}
\item \emph{The Specht modules }$\mathcal{S}^{\lambda}$\emph{:} If $\lambda$
is any partition of $n$, then the \emph{Specht module} $\mathcal{S}^{\lambda}$
is a representation of $S_{n}$ constructed using the Young diagram of shape
$\lambda$. For its definition, see \cite[Definition 5.4.1 \textbf{(b)}]{sga}
(where it is called $\mathcal{S}^{Y\left(  \lambda\right)  }$) or
\cite[\S 7.2]{Fulton97} (where it is called $S^{\lambda}$). If $\mathbf{k}$ is
a field of characteristic $0$, then the Specht module $\mathcal{S}^{\lambda}$
is irreducible.

\item \emph{The trivial representation }$\mathcal{H}_{n}$\emph{:} We let
$\mathcal{H}_{n}$ denote the $\mathbf{k}$-module $\mathbf{k}$, equipped with a
trivial $S_{n}$-action (that is, $\sigma\cdot v=v$ for all $\sigma\in S_{n}$
and $v\in\mathbf{k}$). This is called the \emph{trivial representation} of
$S_{n}$. It is isomorphic to the Specht module $\mathcal{S}^{\left(  n\right)
}$.

\item \emph{The natural representation }$\mathcal{N}_{n}$\emph{:} We let
$\mathcal{N}_{n}$ denote the free $\mathbf{k}$-module $\mathbf{k}^{n}=\left\{
\left(  v_{1},v_{2},\ldots,v_{n}\right)  \ \mid\ \text{all }v_{i}\in
\mathbf{k}\right\}  $, on which $S_{n}$ acts by permuting the coordinates:%
\[
\sigma\cdot\left(  v_{1},v_{2},\ldots,v_{n}\right)  =\left(  v_{\sigma
^{-1}\left(  1\right)  },v_{\sigma^{-1}\left(  2\right)  },\ldots
,v_{\sigma^{-1}\left(  n\right)  }\right)  \ \ \ \ \ \ \ \ \ \ \text{for all
}\sigma\in S_{n}.
\]
This is called the \emph{natural representation} of $S_{n}$.

\item \emph{The reflection quotient representation }$\mathcal{Z}_{n}$\emph{:}
If $n>0$, then the natural representation $\mathcal{N}_{n}$ has a
$1$-dimensional subrepresentation
\[
\mathcal{D}_{n}:=\left\{  \left(  v_{1},v_{2},\ldots,v_{n}\right)
\ \mid\ \text{all }v_{i}\text{ are equal}\right\}  =\left\{  \left(
a,a,\ldots,a\right)  \ \mid\ a\in\mathbf{k}\right\}  .
\]
The quotient
\[
\mathcal{Z}_{n}:=\mathcal{N}_{n}/\mathcal{D}_{n}%
\]
is thus another representation of $S_{n}$. This $\mathcal{Z}_{n}$ is called
the \emph{reflection quotient representation} of $S_{n}$. As a $\mathbf{k}%
$-module, it is free of rank $n-1$ (with basis $\left(  \overline{e_{1}%
},\overline{e_{2}},\ldots,\overline{e_{n-1}}\right)  $, where $e_{1}%
,e_{2},\ldots,e_{n}$ are the standard basis vectors of $\mathbf{k}^{n}$). Here
and in the following, the notation $\overline{v}$ denotes the residue class of
a vector $v$ modulo some submodule (the submodule is to be inferred from the context).

If $\mathbf{k}$ is a field of characteristic $0$ (or, more generally, if $n$
is invertible in $\mathbf{k}$), then this representation $\mathcal{Z}_{n}$ is
isomorphic to the Specht module $\mathcal{S}^{\left(  n-1,1\right)  }$.
Without any such assumptions, $\mathcal{Z}_{n}$ is isomorphic to the dual
$\left(  \mathcal{S}^{\left(  n-1,1\right)  }\right)  ^{\ast}$ of this Specht
module. (See Proposition \ref{prop.Z-iso} below.)
\end{enumerate}

If $V$ is any $\mathbf{k}$-module, then $V^{\ast}$ shall denote its dual
$\mathbf{k}$-module $\operatorname*{Hom}\nolimits_{\mathbf{k}}\left(
V,\mathbf{k}\right)  $. If $V$ is an $S_{n}$-representation, then its dual
$V^{\ast}$ becomes an $S_{n}$-representation as well (see \cite[\S 5.19.3]%
{sga}).

\begin{proposition}
\label{prop.Z-iso}Let $n>1$ be an integer. Then:

\begin{enumerate}
\item[\textbf{(a)}] The reflection quotient representation $\mathcal{Z}_{n}$
is isomorphic (as an $S_{n}$-representation) to the dual $\left(
\mathcal{S}^{\left(  n-1,1\right)  }\right)  ^{\ast}$ of the Specht module
$\mathcal{S}^{\left(  n-1,1\right)  }$.

\item[\textbf{(b)}] If $n$ is invertible in $\mathbf{k}$, then $\mathcal{Z}%
_{n}$ is isomorphic (as an $S_{n}$-representation) to the Specht module
$\mathcal{S}^{\left(  n-1,1\right)  }$.
\end{enumerate}
\end{proposition}

This proposition is clearly part of the folklore, but we outline a proof in
the Appendix (Subsection \ref{subsec.apx.pf.prop.Z-iso}) for the sake of completeness.

\subsection{Tensor products, induction and induction products}

We shall now discuss certain ways to produce new representations from old.

The symbol \textquotedblleft$\otimes$\textquotedblright\ shall always mean a
tensor product over $\mathbf{k}$, unless a different base ring is provided as
a subscript.

It is well-known that if $A$ and $B$ are two $\mathbf{k}$-algebras, then the
tensor product $U\otimes V$ of any left $A$-module $U$ and any left $B$-module
$V$ is canonically a left $A\otimes B$-module. An analogous construction
exists for tensor products of $k$ left modules. Thus, if $U$ is a
representation of a group $G$, and if $V$ is a representation of a group $H$,
then $U\otimes V$ is a representation of $G\times H$, and a similar fact holds
for tensor products of $k$ representations.

We recall the notion of an induced representation: If $G$ is a group, and if
$H$ is a subgroup of $G$, then any $H$-representation $V$ gives rise to a
$G$-representation $\operatorname*{Ind}\nolimits_{H}^{G}V$ defined by%
\begin{equation}
\operatorname*{Ind}\nolimits_{H}^{G}V=\mathbf{k}\left[  G\right]
\otimes_{\mathbf{k}\left[  H\right]  }V, \label{eq.IndGH.def}%
\end{equation}
where we view $\mathbf{k}\left[  G\right]  $ as a $\left(  \mathbf{k}\left[
G\right]  ,\mathbf{k}\left[  H\right]  \right)  $-bimodule while viewing $V$
as a left $\mathbf{k}\left[  H\right]  $-module (so that the tensor product
over $\mathbf{k}\left[  H\right]  $ becomes a left $\mathbf{k}\left[
G\right]  $-module). This $G$-representation $\operatorname*{Ind}%
\nolimits_{H}^{G}V$ is called the \emph{induced representation} of $V$ to $G$.

We furthermore recall the notion of an induction product (\cite[\S 7.3]%
{Fulton97}):

\begin{definition}
Let $n$ and $m$ be two nonnegative integers. Then, the direct product
$S_{n}\times S_{m}$ can be canonically embedded as a subgroup into $S_{n+m}$,
by the group morphism that sends each pair $\left(  \sigma,\tau\right)  \in
S_{n}\times S_{m}$ to the permutation $\sigma\ast\tau\in S_{n+m}$ that applies
$\sigma$ to the first $n$ elements while applying $\tau$ (appropriately
shifted) to the last $m$ elements of $\left[  n+m\right]  $. (To be fully
precise: $\sigma\ast\tau$ is the permutation of $\left[  n+m\right]  $ that
sends $1,2,\ldots,n$ to $\sigma\left(  1\right)  ,\sigma\left(  2\right)
,\ldots,\sigma\left(  n\right)  $ while sending $n+1,n+2,\ldots,n+m$ to
$n+\tau\left(  1\right)  ,n+\tau\left(  2\right)  ,\ldots,n+\tau\left(
m\right)  $.) This is called the \emph{parabolic embedding} of $S_{n}\times
S_{m}$ into $S_{n+m}$.

Now, if $U$ is an $S_{n}$-representation and if $V$ is an $S_{m}%
$-representation, then the tensor product $U\otimes V$ is an $S_{n}\times
S_{m}$-representation, and thus (by the embedding of $S_{n}\times S_{m}$ into
$S_{n+m}$ we just explained) we can construct the induced representation%
\[
U\ast V:=\operatorname*{Ind}\nolimits_{S_{n}\times S_{m}}^{S_{n+m}}\left(
U\otimes V\right)
\]
of $S_{n+m}$. This induced representation $U\ast V$ is called the
\emph{induction product} of $U$ and $V$.

More generally, if $n_{1},n_{2},\ldots,n_{k}$ are any $k$ nonnegative
integers, and if $U_{i}$ is an $S_{n_{i}}$-representation for each
$i\in\left[  k\right]  $, then the \emph{induction product }$U_{1}\ast
U_{2}\ast\cdots\ast U_{k}$ is defined to be the $S_{n_{1}+n_{2}+\cdots+n_{k}}%
$-representation%
\[
\operatorname*{Ind}\nolimits_{S_{n_{1}}\times S_{n_{2}}\times\cdots\times
S_{n_{k}}}^{S_{n_{1}+n_{2}+\cdots+n_{k}}}\left(  U_{1}\otimes U_{2}%
\otimes\cdots\otimes U_{k}\right)  ,
\]
where we embed $S_{n_{1}}\times S_{n_{2}}\times\cdots\times S_{n_{k}}$ into
$S_{n_{1}+n_{2}+\cdots+n_{k}}$ in the obvious way (having each $S_{n_{i}}$ act
on an appropriate interval\footnotemark). The latter embedding is again called
the \emph{parabolic embedding} of $S_{n_{1}}\times S_{n_{2}}\times\cdots\times
S_{n_{k}}$ into $S_{n_{1}+n_{2}+\cdots+n_{k}}$.
\end{definition}

\footnotetext{To make this precise: Let $m_{i}:=n_{1}+n_{2}+\cdots+n_{i}$ for
each $i\in\left[  0,k\right]  $. Then, the integer interval $\left[
n_{1}+n_{2}+\cdots+n_{k}\right]  $ is partitioned into the intervals $\left[
m_{i-1}+1,\ m_{i}\right]  $ for all $i\in\left[  k\right]  $. The embedding of
$S_{n_{1}}\times S_{n_{2}}\times\cdots\times S_{n_{k}}$ into $S_{n_{1}%
+n_{2}+\cdots+n_{k}}$ sends each $k$-tuple $\left(  \sigma_{1},\sigma
_{2},\ldots,\sigma_{k}\right)  \in S_{n_{1}}\times S_{n_{2}}\times\cdots\times
S_{n_{k}}$ to the permutation $\sigma_{1}\ast\sigma_{2}\ast\cdots\ast
\sigma_{k}\in S_{n_{1}+n_{2}+\cdots+n_{k}}$ defined by%
\[
\left(  \sigma_{1}\ast\sigma_{2}\ast\cdots\ast\sigma_{k}\right)  \left(
m_{i-1}+x\right)  :=m_{i-1}+\sigma_{i}\left(  x\right)
\ \ \ \ \ \ \ \ \ \ \text{for each }i\in\left[  k\right]  \text{ and each
}x\in\left[  n_{i}\right]  .
\]
} These induction products satisfy associativity up to isomorphism: e.g., we
have isomorphisms
$\left(  U\ast V\right)  \ast W\cong U\ast V\ast W\cong U\ast\left(
V\ast W\right)  $ for all $U,V,W$. More generally:

\begin{proposition}
\label{prop.indprod.ass}Let $n_{1},n_{2},\ldots,n_{k}$ be any $k$ nonnegative
integers, and let $U_{i}$ be an $S_{n_{i}}$-representation for each
$i\in\left[  k\right]  $. Let $i\in\left[  0,k\right]  $. Then,%
\[
U_{1}\ast U_{2}\ast\cdots\ast U_{k}\cong\left(  U_{1}\ast U_{2}\ast\cdots\ast
U_{i}\right)  \ast\left(  U_{i+1}\ast U_{i+2}\ast\cdots\ast U_{k}\right)  .
\]

\end{proposition}

This is again a folklore result, but we sketch a proof in the Appendix
(Subsection \ref{subsec.apx.pf.prop.indprod.ass}) to fill a little gap in the literature.

\subsection{Lacunar sets and the submodules $F\left(  I\right)  $}

Next, we recall some more concepts from \cite{s2b1}.

If $I$ is a finite set of integers, then we let $\operatorname*{sum}I$ denote
the sum of all elements of $I$. For instance, $\operatorname*{sum}\left\{
3,7\right\}  =3+7=10$.

Let $\left(  f_{0},f_{1},f_{2},\ldots\right)  $ be the \emph{Fibonacci
sequence}. This is the sequence of integers defined recursively by%
\[
f_{0}=0,\ \ \ \ \ \ \ \ \ \ f_{1}=1,\ \ \ \ \ \ \ \ \ \ \text{and}%
\ \ \ \ \ \ \ \ \ \ f_{m}=f_{m-1}+f_{m-2}\text{ for all }m\geq2.
\]


We shall say that a set $I\subseteq\mathbb{Z}$ is \emph{lacunar} if it
contains no two consecutive integers (i.e., there exists no $i\in I$ such that
$i+1\in I$). For instance, the set $\left\{  1,4,6\right\}  $ is lacunar,
while the set $\left\{  1,4,5\right\}  $ is not.

The number of lacunar subsets of $\left[  n-1\right]  $ is the Fibonacci
number $f_{n+1}$. Let $Q_{1},Q_{2},\ldots,Q_{f_{n+1}}$ be all these $f_{n+1}$
lacunar subsets of $\left[  n-1\right]  $, listed in an order that satisfies%
\begin{equation}
\operatorname*{sum}\left(  Q_{1}\right)  \leq\operatorname*{sum}\left(
Q_{2}\right)  \leq\cdots\leq\operatorname*{sum}\left(  Q_{f_{n+1}}\right)  .
\label{pf.thm.t-simultri.sum-order}%
\end{equation}
We fix this order once and for all\footnote{For $n\leq3$, this order is
uniquely defined. For $n>3$, we need to make a choice.}. Many of our
constructions will formally (though rather shallowly) depend on this order.

For any subset $I$ of $\left[  n\right]  $, we define the following:

\begin{itemize}
\item We let $I-1$ denote the set $\left\{  i-1\ \mid\ i\in I\right\}
=\left\{  j\in\mathbb{Z}\ \mid\ j+1\in I\right\}  $. For instance, $\left\{
2,4,5\right\}  -1=\left\{  1,3,4\right\}  $. Note that $I$ is lacunar if and
only if $I\cap\left(  I-1\right)  =\varnothing$.

\item We let $I^{\prime}$ be the set $\left[  n-1\right]  \setminus\left(
I\cup\left(  I-1\right)  \right)  $. This is the set of all $i\in\left[
n-1\right]  $ satisfying $i\notin I$ and $i+1\notin I$. We shall refer to
$I^{\prime}$ as the \emph{non-shadow} of $I$.

For example, if $n=5$, then $\left\{  2,3\right\}  ^{\prime}=\left[  4\right]
\setminus\left\{  1,2,3\right\}  =\left\{  4\right\}  $.

\item We let
\[
F\left(  I\right)  :=\left\{  q\in\mathbf{k}\left[  S_{n}\right]
\ \mid\ qs_{i}=q\text{ for all }i\in I^{\prime}\right\}  .
\]
We can rewrite this equality as%
\begin{align}
F\left(  I\right)   &  =\left\{  t\in\mathbf{k}\left[  S_{n}\right]
\ \mid\ ts_{j}=t\text{ for all }j\in I^{\prime}\right\} \nonumber\\
&  =\left\{  t\in\mathcal{A}\ \mid\ ts_{j}=t\text{ for all }j\in I^{\prime
}\right\}  \label{eq.FI=.2}%
\end{align}
(since $\mathbf{k}\left[  S_{n}\right]  =\mathcal{A}$).
\end{itemize}

\section{The first main theorem: the Fibonacci filtration}

\subsection{The theorem}

For each $i\in\left[  0,f_{n+1}\right]  $, we define a $\mathbf{k}$-submodule%
\[
F_{i}:=F\left(  Q_{1}\right)  +F\left(  Q_{2}\right)  +\cdots+F\left(
Q_{i}\right)  \ \ \ \ \ \ \ \ \ \ \text{of }\mathbf{k}\left[  S_{n}\right]
\]
(so that $F_{0}=0$). In \cite[Theorem 8.1]{s2b1}, the following is shown:

\begin{theorem}
\label{thm.t-simultri}\ 

\begin{enumerate}
\item[\textbf{(a)}] We have%
\[
0=F_{0}\subseteq F_{1}\subseteq F_{2}\subseteq\cdots\subseteq F_{f_{n+1}%
}=\mathbf{k}\left[  S_{n}\right]  .
\]
In other words, the $\mathbf{k}$-submodules $F_{0},F_{1},\ldots,F_{f_{n+1}}$
form a $\mathbf{k}$-module filtration of $\mathbf{k}\left[  S_{n}\right]  $.

\item[\textbf{(b)}] We have $F_{i}\cdot t_{\ell}\subseteq F_{i}$ for each
$i\in\left[  0,f_{n+1}\right]  $ and $\ell\in\left[  n\right]  $.

\item[\textbf{(c)}] For each $i\in\left[  f_{n+1}\right]  $ and $\ell
\in\left[  n\right]  $, we have%
\[
F_{i}\cdot\left(  t_{\ell}-m_{Q_{i},\ell}\right)  \subseteq F_{i-1}.
\]
Here, $m_{Q_{i},\ell}$ is a certain integer whose definition we will give in
Subsection \ref{subsec.specht-spec.thm} (as we will not use it until then).
\end{enumerate}
\end{theorem}

The filtration $0=F_{0}\subseteq F_{1}\subseteq F_{2}\subseteq\cdots\subseteq
F_{f_{n+1}}=\mathbf{k}\left[  S_{n}\right]  $ will be called the
\emph{Fibonacci filtration} of $\mathcal{A}$. We can easily see that it is a
filtration of $\left(  \mathcal{A},\mathcal{T}\right)  $-bimodules:

\begin{proposition}
\label{prop.fibfilt.A}Let $i\in\left[  0,f_{n+1}\right]  $. Then, $F_{i}$ is
an $\left(  \mathcal{A},\mathcal{T}\right)  $-subbimodule of $\mathcal{A}$.
\end{proposition}

\begin{proof}
For any $I\subseteq\left[  n\right]  $, the set $F\left(  I\right)  $ is
closed under addition and left action of $\mathcal{A}$ (by its very
definition), hence is a left $\mathcal{A}$-submodule of $\mathcal{A}$. Thus,
$F_{i}$ (being defined as a sum of such sets $F\left(  I\right)  $) is also a
left $\mathcal{A}$-submodule of $\mathcal{A}$. Moreover, $F_{i}$ is also
closed under right multiplication by each $t_{\ell}$ (by Theorem
\ref{thm.t-simultri} \textbf{(b)}), and hence under the right action of
$\mathcal{T}$ (since $\mathcal{T}$ is the subalgebra generated by $t_{1}%
,t_{2},\ldots,t_{n}$). Thus, $F_{i}$ is also a right $\mathcal{T}$-submodule
of $\mathcal{A}$. Altogether, we conclude that $F_{i}$ is an $\left(
\mathcal{A},\mathcal{T}\right)  $-subbimodule of $\mathcal{A}$.
\end{proof}

Proposition \ref{prop.fibfilt.A} shows that the subquotients $F_{i}/F_{i-1}$
are $\left(  \mathcal{A},\mathcal{T}\right)  $-bimodules as well. In
particular, they are therefore left $\mathcal{A}$-modules, i.e.,
representations of $S_{n}$. Our second main theorem characterizes these representations:

\begin{theorem}
\label{thm.Fi/Fi-1.as-ind}Let $i\in\left[  f_{n+1}\right]  $. Consider the
lacunar subset $Q_{i}$ of $\left[  n-1\right]  $. Write the set $Q_{i}%
\cup\left\{  n+1\right\}  $ as $\left\{  i_{1}<i_{2}<\cdots<i_{m}\right\}  $,
so that $i_{m}=n+1$. Furthermore, set $i_{0}:=1$. Set $j_{k}:=i_{k}-i_{k-1}$
for each $k\in\left[  m\right]  $. (Note that $j_{1}\geq0$ and $j_{2}%
,j_{3},\ldots,j_{m}>1$ and $j_{1}+j_{2}+\cdots+j_{m}=n$; this follows from
Lemma \ref{lem.Fi/Fi-1.sum-js} below (applied to $I=Q_{i}$).) Then,%
\[
F_{i}/F_{i-1}\cong\underbrace{\mathcal{H}_{j_{1}}\ast\mathcal{Z}_{j_{2}}%
\ast\mathcal{Z}_{j_{3}}\ast\cdots\ast\mathcal{Z}_{j_{m}}}_{\substack{\text{the
first factor is an }\mathcal{H}\text{,}\\\text{while all others are
}\mathcal{Z}\text{'s}}}
\]
as $S_{n}$-representations.\footnotemark
\end{theorem}

\footnotetext{Note that the factor $\mathcal{H}_{j_{1}}$ can be omitted when
$j_{1}=0$, since $\mathcal{H}_{0}\cong\mathbf{k}$ with the trivial $S_{0}%
$-action.}We will spend the rest of this section proving this theorem, then
restating it (in the characteristic-$0$ case) using Littlewood--Richardson coefficients.

\subsection{Lemmas on $F\left(  Q_{i}\right)  $}

First, let us show some lemmas about lacunar sets $I$ and the corresponding
$\mathbf{k}$-modules $F\left(  I\right)  $:

\begin{lemma}
\label{lem.Fi/Fi-1.sum-js}Let $I$ be a lacunar subset of $\left[  n-1\right]
$. Write the set $I\cup\left\{  n+1\right\}  $ as $\left\{  i_{1}<i_{2}%
<\cdots<i_{m}\right\}  $, so that $i_{m}=n+1$. Furthermore, set $i_{0}:=1$.
Set $j_{k}:=i_{k}-i_{k-1}$ for each $k\in\left[  m\right]  $. Then, $j_{1}%
\geq0$ and $j_{2},j_{3},\ldots,j_{m}>1$ and $j_{1}+j_{2}+\cdots+j_{m}=n$.
\end{lemma}

\begin{proof}
By definition, we have $j_{1}=i_{1}-i_{0}\geq0$, since $i_{1}\geq1=i_{0}$.

Next, we recall that the set $I$ is lacunar. This lacunarity is preserved even
when we insert the new element $n+1$ into this set, since all existing
elements of $I$ are $\leq n-1$ (since $I\subseteq\left[  n-1\right]  $) and
thus cannot be consecutive with $n+1$. That is, the set $I\cup\left\{
n+1\right\}  $ is again lacunar. Since we have written this set as $\left\{
i_{1}<i_{2}<\cdots<i_{m}\right\}  $, this yields that any $k\in\left[
2,m\right]  $ satisfies $i_{k}-i_{k-1}>1$. In other words, any $k\in\left[
2,m\right]  $ satisfies $j_{k}>1$ (since $j_{k}=i_{k}-i_{k-1}$). In other
words, $j_{2},j_{3},\ldots,j_{m}>1$.

It remains to prove that $j_{1}+j_{2}+\cdots+j_{m}=n$. But recall that
$j_{k}=i_{k}-i_{k-1}$ for each $k\in\left[  m\right]  $. Hence,%
\begin{align*}
\sum_{k=1}^{m}j_{k}  &  =\sum_{k=1}^{m}\left(  i_{k}-i_{k-1}\right)
=\underbrace{i_{m}}_{=n+1}-\underbrace{i_{0}}_{=1}\ \ \ \ \ \ \ \ \ \ \left(
\text{by the telescope principle}\right) \\
&  =n+1-1=n.
\end{align*}
In other words, $j_{1}+j_{2}+\cdots+j_{m}=n$. Thus, Lemma
\ref{lem.Fi/Fi-1.sum-js} is fully proved.
\end{proof}

\begin{lemma}
\label{lem.Fi/Fi-1.dim}Let $i\in\left[  f_{n+1}\right]  $. Consider the
lacunar subset $Q_{i}$ of $\left[  n-1\right]  $. Write the set $Q_{i}%
\cup\left\{  n+1\right\}  $ as $\left\{  i_{1}<i_{2}<\cdots<i_{m}\right\}  $.
Furthermore, set $i_{0}:=1$. Set $j_{k}:=i_{k}-i_{k-1}$ for each $k\in\left[
m\right]  $. Then, the $\mathbf{k}$-module $F_{i}/F_{i-1}$ is free of rank%
\[
\dfrac{n!}{j_{1}!j_{2}!\cdots j_{m}!}\cdot\prod_{k=2}^{m}\left(
j_{k}-1\right)  .
\]

\end{lemma}

\begin{proof}
We have $Q_{i}\subseteq\left[  n-1\right]  $. Hence, $n+1$ is the largest
element of $Q_{i}\cup\left\{  n+1\right\}  $. Thus, from $Q_{i}\cup\left\{
n+1\right\}  =\left\{  i_{1}<i_{2}<\cdots<i_{m}\right\}  $, we obtain
\[
i_{m}=n+1\ \ \ \ \ \ \ \ \ \ \text{and}\ \ \ \ \ \ \ \ \ \ Q_{i}=\left\{
i_{1}<i_{2}<\cdots<i_{m-1}\right\}  .
\]


Lemma \ref{lem.Fi/Fi-1.sum-js} (applied to $I=Q_{i}$) shows that $j_{1}%
+j_{2}+\cdots+j_{m}=n$. Let $\dbinom{n}{j_{1},j_{2},\ldots,j_{m}}$ denote the
multinomial coefficient $\dfrac{n!}{j_{1}!j_{2}!\cdots j_{m}!}$. We know from
\cite[Theorem 13.1 \textbf{(a)} and \textbf{(c)}]{s2b1} (applied to $p=m-1$)
that the $\mathbf{k}$-module $F_{i}/F_{i-1}$ is free of rank%
\[
\delta_{i}=\dbinom{n}{j_{1},j_{2},\ldots,j_{m}}\cdot\prod_{k=2}^{m}\left(
j_{k}-1\right)  .
\]
In view of $\dbinom{n}{j_{1},j_{2},\ldots,j_{m}}=\dfrac{n!}{j_{1}!j_{2}!\cdots
j_{m}!}$, this is precisely the claim of Lemma \ref{lem.Fi/Fi-1.dim}.
\end{proof}

\begin{lemma}
\label{lem.Fi/Fi-1.L2a}Let $I$ be a subset of $\left[  n\right]  $. Let $j\in
I$. Then, there exists a lacunar subset $J$ of $\left[  n-1\right]  $ such
that $\operatorname*{sum}J<\operatorname*{sum}I$ and $J^{\prime}\subseteq
I^{\prime}\cup\left\{  j\right\}  $.
\end{lemma}

\begin{proof}
Set $K:=\left(  I\setminus\left\{  j\right\}  \right)  \cup\left\{
j-1\right\}  $ if $j>1$, and otherwise set $K:=I\setminus\left\{  j\right\}
$. Then, $K$ is a subset of $\left[  n\right]  $ and satisfies
$\operatorname*{sum}K<\operatorname*{sum}I$ (since $K$ is obtained from $I$ by
removing the element $j$ and possibly inserting the smaller element $j-1$).
Furthermore, \cite[Proposition 8.6 \textbf{(a)}]{s2b1} says that $K^{\prime
}\subseteq I^{\prime}\cup\left\{  j\right\}  $.

Now, \cite[Corollary 8.8]{s2b1} (applied to $K$ instead of $I$) shows that
there exists a lacunar subset $J$ of $\left[  n-1\right]  $ such that
$\operatorname*{sum}J\leq\operatorname*{sum}K$ and $J^{\prime}\subseteq
K^{\prime}$. Consider this $J$.

The set $J$ is a lacunar subset of $\left[  n-1\right]  $ and satisfies
$\operatorname*{sum}J<\operatorname*{sum}I$ (since $\operatorname*{sum}%
J\leq\operatorname*{sum}K<\operatorname*{sum}I$) and $J^{\prime}\subseteq
I^{\prime}\cup\left\{  j\right\}  $ (since $J^{\prime}\subseteq K^{\prime
}\subseteq I^{\prime}\cup\left\{  j\right\}  $). Hence, such a $J$ exists.
This proves Lemma \ref{lem.Fi/Fi-1.L2a}.
\end{proof}

\begin{lemma}
\label{lem.Fi/Fi-1.L2}Let $i\in\left[  f_{n+1}\right]  $. Consider the lacunar
subset $Q_{i}$ of $\left[  n-1\right]  $. Write the set $Q_{i}\cup\left\{
n+1\right\}  $ as $\left\{  i_{1}<i_{2}<\cdots<i_{m}\right\}  $. Let
$k\in\left[  m-1\right]  $. Then,%
\[
\left\{  t\in F\left(  Q_{i}\right)  \ \mid\ ts_{i_{k}}=t\right\}  \subseteq
F_{i-1}.
\]

\end{lemma}

\begin{proof}
As in the proof of Lemma \ref{lem.Fi/Fi-1.dim}, we find $Q_{i}=\left\{
i_{1}<i_{2}<\cdots<i_{m-1}\right\}  $. Thus, $i_{k}\in Q_{i}$ (since
$k\in\left[  m-1\right]  $). Hence, Lemma \ref{lem.Fi/Fi-1.L2a} (applied to
$I=Q_{i}$ and $j=i_{k}$) shows that there exists a lacunar subset $J$ of
$\left[  n-1\right]  $ such that $\operatorname*{sum}J<\operatorname*{sum}%
\left(  Q_{i}\right)  $ and $J^{\prime}\subseteq Q_{i}^{\prime}\cup\left\{
i_{k}\right\}  $. Consider this $J$. Since $J$ is lacunar, we have $J=Q_{s}$
for some $s\in\left[  f_{n+1}\right]  $. Consider this $s$. Thus, $Q_{s}=J$,
so that $\operatorname*{sum}\left(  Q_{s}\right)  =\operatorname*{sum}%
J<\operatorname*{sum}\left(  Q_{i}\right)  $ and therefore $s<i$ (by
(\ref{pf.thm.t-simultri.sum-order})). Hence, $s\leq i-1$, so that $F\left(
Q_{s}\right)  \subseteq F_{i-1}$ (since the definition of $F_{i-1}$ says that
$F_{i-1}=F\left(  Q_{1}\right)  +F\left(  Q_{2}\right)  +\cdots+F\left(
Q_{i-1}\right)  $). In view of $Q_{s}=J$, we can rewrite this as%
\[
F\left(  J\right)  \subseteq F_{i-1}.
\]


Now, (\ref{eq.FI=.2}) (applied to $Q_{i}$ instead of $I$) shows that%
\[
F\left(  Q_{i}\right)  =\left\{  t\in\mathcal{A}\ \mid\ ts_{j}=t\text{ for all
}j\in Q_{i}^{\prime}\right\}  ,
\]
so that%
\begin{align*}
&  \left\{  t\in F\left(  Q_{i}\right)  \ \mid\ ts_{i_{k}}=t\right\} \\
&  =\left\{  t\in\mathcal{A}\ \mid\ ts_{j}=t\text{ for all }j\in Q_{i}%
^{\prime}\text{, and also }ts_{i_{k}}=t\right\} \\
&  =\left\{  t\in\mathcal{A}\ \mid\ ts_{j}=t\text{ for all }j\in Q_{i}%
^{\prime}\cup\left\{  i_{k}\right\}  \right\} \\
&  \subseteq\left\{  t\in\mathcal{A}\ \mid\ ts_{j}=t\text{ for all }j\in
J^{\prime}\right\}  \ \ \ \ \ \ \ \ \ \ \left(  \text{since }J^{\prime
}\subseteq Q_{i}^{\prime}\cup\left\{  i_{k}\right\}  \right) \\
&  =F\left(  J\right)  \ \ \ \ \ \ \ \ \ \ \left(  \text{by (\ref{eq.FI=.2}),
applied to }J\text{ instead of }I\right) \\
&  \subseteq F_{i-1}.
\end{align*}
Thus, Lemma \ref{lem.Fi/Fi-1.L2} follows.
\end{proof}

\subsection{The elements $\nabla_{\mathbf{p}}$}

\begin{lemma}
\label{lem.Fi/Fi-1.L3}Let $i\in\left[  f_{n+1}\right]  $. Consider the lacunar
subset $Q_{i}$ of $\left[  n-1\right]  $. Write the set $Q_{i}\cup\left\{
n+1\right\}  $ as $\left\{  i_{1}<i_{2}<\cdots<i_{m}\right\}  $. Furthermore,
set $i_{0}:=1$.

For each $k\in\left[  m\right]  $, let $J_{k}$ denote the integer interval
$\left[  i_{k-1},\ i_{k}-1\right]  $. Note that the intervals $J_{1}%
,J_{2},\ldots,J_{m}$ are disjoint and -- except possibly for $J_{1}$ --
nonempty ($J_{1}$ is empty if and only if $1\in Q_{i}$), and their union is
$\left[  n\right]  $. Thus, we can view the direct product $S_{J_{1}}\times
S_{J_{2}}\times\cdots\times S_{J_{m}}$ as a subgroup of $S_{n}$ in the obvious
way (each factor $S_{J_{k}}$ acts on the elements of $J_{k}$ while leaving all
remaining elements of $\left[  n\right]  $ unchanged).

For each $\left(  m-1\right)  $-tuple $\mathbf{p}=\left(  p_{2},p_{3}%
,\ldots,p_{m}\right)  \in J_{2}\times J_{3}\times\cdots\times J_{m}$ (that is,
with $p_{k}\in J_{k}$ for each $k\in\left[  2,m\right]  $), we define an
element%
\[
\nabla_{\mathbf{p}}:=\sum_{\substack{\sigma\in S_{n};\\\sigma\left(
J_{k}\right)  =J_{k}\text{ for each }k\in\left[  m\right]  ;\\\sigma\left(
i_{k-1}\right)  =p_{k}\text{ for each }k\in\left[  2,m\right]  }}\sigma
\in\mathcal{A}.
\]
(Note that this also depends on $i$, not just on $\mathbf{p}$.)

Then:

\begin{enumerate}
\item[\textbf{(a)}] For any $\tau=\left(  \tau_{1},\tau_{2},\ldots,\tau
_{m}\right)  \in S_{J_{1}}\times S_{J_{2}}\times\cdots\times S_{J_{m}}$ and
$\mathbf{p}=\left(  p_{2},p_{3},\ldots,p_{m}\right)  \in J_{2}\times
J_{3}\times\cdots\times J_{m}$, we have%
\[
\tau\nabla_{\mathbf{p}}=\nabla_{\tau\mathbf{p}},
\]
where%
\[
\tau\mathbf{p}:=\left(  \tau_{2}\left(  p_{2}\right)  ,\tau_{3}\left(
p_{3}\right)  ,\ldots,\tau_{m}\left(  p_{m}\right)  \right)  =\left(
\tau\left(  p_{2}\right)  ,\tau\left(  p_{3}\right)  ,\ldots,\tau\left(
p_{m}\right)  \right)  .
\]


\item[\textbf{(b)}] The left $\mathcal{A}$-module $F\left(  Q_{i}\right)  $ is
generated (as a left $\mathcal{A}$-module) by any single element of the form
$\nabla_{\mathbf{p}}$ (with $\mathbf{p}\in J_{2}\times J_{3}\times\cdots\times
J_{m}$).

\item[\textbf{(c)}] Let $\ell\in\left[  2,m\right]  $. For each $k\in\left[
2,m\right]  \setminus\left\{  \ell\right\}  $, let $p_{k}\in J_{k}$ be an
element. Then,%
\[
\sum_{p_{\ell}\in J_{\ell}}\nabla_{\left(  p_{2},p_{3},\ldots,p_{m}\right)
}\in F_{i-1}.
\]
(Note that the elements $p_{2},p_{3},\ldots,p_{\ell-1},p_{\ell+1},p_{\ell
+2},\ldots,p_{m}$ in this sum are fixed, whereas $p_{\ell}$ runs through the
set $J_{\ell}$.)
\end{enumerate}
\end{lemma}

\begin{example}
\label{exa.Fi/Fi-1.L3.e1}Let $n=7$ and $Q_{i}=\left\{  3,6\right\}  $ (clearly
a lacunar subset of $\left[  n-1\right]  =\left[  6\right]  $). Then,
following the notations of Lemma \ref{lem.Fi/Fi-1.L3}, we have $\left\{
i_{1}<i_{2}<\cdots<i_{m}\right\}  =Q_{i}\cup\left\{  n+1\right\}  =\left\{
3,6,8\right\}  $, so that $i_{1}=3$ and $i_{2}=6$ and $i_{3}=8$ and $i_{0}=1$.
Thus, the integer intervals $J_{k}=\left[  i_{k-1},\ i_{k}-1\right]  $ are
\[
J_{1}=\left[  1,2\right]  ,\ \ \ \ \ \ \ \ \ \ J_{2}=\left[  3,5\right]
,\ \ \ \ \ \ \ \ \ \ J_{3}=\left[  6,7\right]  .
\]
Taking $\mathbf{p}$ to be the $2$-tuple $\left(  p_{2},p_{3}\right)  =\left(
4,7\right)  \in J_{2}\times J_{3}$, we then have%
\begin{align*}
\nabla_{\mathbf{p}}  & =\sum_{\substack{\sigma\in S_{n};\\\sigma\left(
J_{k}\right)  =J_{k}\text{ for each }k\in\left[  m\right]  ;\\\sigma\left(
i_{k-1}\right)  =p_{k}\text{ for each }k\in\left[  2,m\right]  }}\sigma
=\sum_{\substack{\sigma\in S_{n};\\\sigma\left(  \left[  1,2\right]  \right)
=\left[  1,2\right]  ;\ \sigma\left(  \left[  3,5\right]  \right)  =\left[
3,5\right]  ;\ \sigma\left(  \left[  6,7\right]  \right)  =\left[  6,7\right]
;\\\sigma\left(  3\right)  =4;\ \sigma\left(  6\right)  =7}}\sigma\\
& =\left[  1243576\right]  +\left[  1245376\right]  +\left[  2143576\right]
+\left[  2145376\right]
\end{align*}
(writing permutations in one-line notation). Lemma \ref{lem.Fi/Fi-1.L3}
\textbf{(b)} says that this element $\nabla_{\mathbf{p}}$ generates the left
$\mathcal{A}$-module $F\left(  Q_{i}\right)  $. Applying Lemma
\ref{lem.Fi/Fi-1.L3} \textbf{(c)} to $\ell=2$ and $p_{3}=7$, we obtain%
\[
\sum_{p_{2}\in J_{2}}\nabla_{\left(  p_{2},7\right)  }\in F_{i-1}%
,\ \ \ \ \ \ \ \ \ \ \text{that is,}\ \ \ \ \ \ \ \ \ \ \nabla_{\left(
3,7\right)  }+\nabla_{\left(  4,7\right)  }+\nabla_{\left(  5,7\right)  }\in
F_{i-1}.
\]

\end{example}

\begin{proof}
[Proof of Lemma \ref{lem.Fi/Fi-1.L3}.] As in the proof of Lemma
\ref{lem.Fi/Fi-1.dim}, we find $i_{m}=n+1$ and $Q_{i}=\left\{  i_{1}%
<i_{2}<\cdots<i_{m-1}\right\}  $. \medskip

\textbf{(a)} Let $\tau=\left(  \tau_{1},\tau_{2},\ldots,\tau_{m}\right)  \in
S_{J_{1}}\times S_{J_{2}}\times\cdots\times S_{J_{m}}$ and $\mathbf{p}=\left(
p_{2},p_{3},\ldots,p_{m}\right)  \in J_{2}\times J_{3}\times\cdots\times
J_{m}$. Recall that
\[
\nabla_{\mathbf{p}}=\sum_{\substack{\sigma\in S_{n};\\\sigma\left(
J_{k}\right)  =J_{k}\text{ for each }k\in\left[  m\right]  ;\\\sigma\left(
i_{k-1}\right)  =p_{k}\text{ for each }k\in\left[  2,m\right]  }}\sigma.
\]
Multiplying this equality by $\tau$ from the left, we obtain%
\begin{align*}
\tau\nabla_{\mathbf{p}} &  =\sum_{\substack{\sigma\in S_{n};\\\sigma\left(
J_{k}\right)  =J_{k}\text{ for each }k\in\left[  m\right]  ;\\\sigma\left(
i_{k-1}\right)  =p_{k}\text{ for each }k\in\left[  2,m\right]  }}\tau
\sigma=\sum_{\substack{\sigma\in S_{n};\\\tau\left(  \sigma\left(
J_{k}\right)  \right)  =\tau\left(  J_{k}\right)  \text{ for each }k\in\left[
m\right]  ;\\\tau\left(  \sigma\left(  i_{k-1}\right)  \right)  =\tau\left(
p_{k}\right)  \text{ for each }k\in\left[  2,m\right]  }}\tau\sigma\\
&  \ \ \ \ \ \ \ \ \ \ \ \ \ \ \ \ \ \ \ \ \left(
\begin{array}
[c]{c}%
\text{here, we have replaced the conditions \textquotedblleft}\sigma\left(
J_{k}\right)  =J_{k}\text{\textquotedblright}\\
\text{and \textquotedblleft}\sigma\left(  i_{k-1}\right)  =p_{k}%
\text{\textquotedblright\ under the summation sign}\\
\text{by the conditions \textquotedblleft}\tau\left(  \sigma\left(
J_{k}\right)  \right)  =\tau\left(  J_{k}\right)  \text{\textquotedblright}\\
\text{and \textquotedblleft}\tau\left(  \sigma\left(  i_{k-1}\right)  \right)
=\tau\left(  p_{k}\right)  \text{\textquotedblright\ (which are equivalent
to}\\
\text{the former two conditions because }\tau\text{ is injective)}%
\end{array}
\right)  \\
&  =\sum_{\substack{\sigma\in S_{n};\\\left(  \tau\sigma\right)  \left(
J_{k}\right)  =J_{k}\text{ for each }k\in\left[  m\right]  ;\\\left(
\tau\sigma\right)  \left(  i_{k-1}\right)  =\tau\left(  p_{k}\right)  \text{
for each }k\in\left[  2,m\right]  }}\tau\sigma\\
&  \ \ \ \ \ \ \ \ \ \ \ \ \ \ \ \ \ \ \ \ \left(
\begin{array}
[c]{c}%
\text{since each }k\in\left[  m\right]  \text{ satisfies }\tau\left(
\sigma\left(  J_{k}\right)  \right)  =\left(  \tau\sigma\right)  \left(
J_{k}\right)  \\
\text{and }\tau\left(  \sigma\left(  i_{k-1}\right)  \right)  =\left(
\tau\sigma\right)  \left(  i_{k-1}\right)  \text{ (if }k>1\text{)}\\
\text{and }\tau\left(  J_{k}\right)  =J_{k}\text{ (since }\tau\in S_{J_{1}%
}\times S_{J_{2}}\times\cdots\times S_{J_{m}}\text{)}%
\end{array}
\right)  \\
&  =\sum_{\substack{\sigma\in S_{n};\\\sigma\left(  J_{k}\right)  =J_{k}\text{
for each }k\in\left[  m\right]  ;\\\sigma\left(  i_{k-1}\right)  =\tau\left(
p_{k}\right)  \text{ for each }k\in\left[  2,m\right]  }}\sigma\\
&  \ \ \ \ \ \ \ \ \ \ \ \ \ \ \ \ \ \ \ \ \left(
\begin{array}
[c]{c}%
\text{here, we have substituted }\sigma\text{ for }\tau\sigma\text{ in the
sum,}\\
\text{since the map }S_{n}\rightarrow S_{n},\ \sigma\mapsto\tau\sigma\text{ is
a bijection}%
\end{array}
\right)  \\
&  =\nabla_{\tau\mathbf{p}}%
\end{align*}
(by the definition of $\nabla_{\tau\mathbf{p}}$, since $\tau\mathbf{p}=\left(
\tau\left(  p_{2}\right)  ,\tau\left(  p_{3}\right)  ,\ldots,\tau\left(
p_{m}\right)  \right)  $). This proves Lemma \ref{lem.Fi/Fi-1.L3}
\textbf{(a)}. \medskip

\textbf{(b)} The definition of the non-shadow $Q_{i}^{\prime}$ yields%
\begin{align*}
Q_{i}^{\prime} &  =\left[  n-1\right]  \setminus\left(  Q_{i}\cup\left(
Q_{i}-1\right)  \right)  \\
&  =\left[  n-1\right]  \setminus\left(  \left\{  i_{1}<i_{2}<\cdots
<i_{m-1}\right\}  \cup\left\{  i_{1}-1<i_{2}-1<\cdots<i_{m-1}-1\right\}
\right)
\end{align*}
(since $Q_{i}=\left\{  i_{1}<i_{2}<\cdots<i_{m-1}\right\}  $). In other words,
$Q_{i}^{\prime}$ consists of all elements of $\left[  n-1\right]  $ except for
those of the forms $i_{k}-1$ and $i_{k}$ for $k\in\left[  m-1\right]  $.

Let $\Gamma$ be the subgroup of $S_{n}$ generated by the simple transpositions
$s_{j}$ with $j\in Q_{i}^{\prime}$. Thus, $\Gamma$ is generated by all simple
transpositions $s_{1},s_{2},\ldots,s_{n-1}$ except for those of the forms
$s_{i_{k}-1}$ and $s_{i_{k}}$ for $k\in\left[  m-1\right]  $ (by the
description of $Q_{i}^{\prime}$ in the previous paragraph). Hence, every
permutation $\omega\in\Gamma$ preserves the intervals $J_{1},J_{2}%
,\ldots,J_{m}$ as well as the elements $i_{1},i_{2},\ldots,i_{m-1}$.

Conversely, if some permutation $\omega\in S_{n}$ preserves the intervals
$J_{1},J_{2},\ldots,J_{m}$ as well as the elements $i_{1},i_{2},\ldots
,i_{m-1}$, then $\omega$ must belong to $\Gamma$ (because such a permutation
$\omega$ must preserve the intervals $J_{1}=\left[  i_{0},\ i_{1}-1\right]  $
as well as $J_{k}\setminus\left\{  i_{k-1}\right\}  =\left[  i_{k-1}%
+1,\ i_{k}-1\right]  $ for all $k\in\left[  2,m\right]  $ (since it preserves
both $J_{k}$ and $i_{k-1}$) as well as the length-$1$ intervals $\left\{
i_{k-1}\right\}  $ for all $k\in\left[  2,m\right]  $, and thus must be a
composition of permutations of these intervals; but any such permutation
belongs to $\Gamma$ (since any permutation of an integer interval $\left[
a,b\right]  $ can be written as a product of simple transpositions $s_{j}$
with $j\in\left[  a,b-1\right]  $)).

The subgroup $\Gamma$ of $S_{n}$ acts from the right on $S_{n}$ (simply by
right multiplication), and thus also acts $\mathbf{k}$-linearly from the right
on $\mathcal{A}=\mathbf{k}\left[  S_{n}\right]  $ (by linear extension),
making $\mathcal{A}$ into a permutation module\footnote{Recall the definition
of a permutation module:
\par
Let $G$ be a finite group. Let $X$ be a right $G$-set. Let $\mathbf{k}%
^{\left(  X\right)  }$ be the free $\mathbf{k}$-module with basis $X$. Then,
$\mathbf{k}^{\left(  X\right)  }$ becomes a right $\mathbf{k}\left[  G\right]
$-module, where the action of $\mathbf{k}\left[  G\right]  $ on $\mathbf{k}%
^{\left(  X\right)  }$ is given by bilinearly extending the action of $G$ on
$X$ (that is, by the rule $\left(  \sum_{g\in G}\alpha_{g}g\right)  \left(
\sum_{x\in X}\beta_{x}x\right)  :=\sum_{g\in G}\ \ \sum_{x\in X}\alpha
_{g}\beta_{x}gx$). This is called the \emph{permutation module} corresponding
to the right $G$-set $X$.
\par
In our present setup, we apply this construction to $G=\Gamma$ and $X=S_{n}$.}
of $\Gamma$. Applying (\ref{eq.FI=.2}) to $I=Q_{i}$, we see that%
\begin{align*}
F\left(  Q_{i}\right)   &  =\left\{  t\in\mathcal{A}\ \mid\ ts_{j}=t\text{ for
all }j\in Q_{i}^{\prime}\right\}  \\
&  =\left\{  t\in\mathcal{A}\ \mid\ t\omega=t\text{ for all }\omega\in
\Gamma\right\}
\end{align*}
(since $\Gamma$ is the group generated by the $s_{j}$ with $j\in Q_{i}%
^{\prime}$, and therefore the condition \textquotedblleft$ts_{j}=t$ for all
$j\in Q_{i}^{\prime}$\textquotedblright\ is equivalent to \textquotedblleft%
$t\omega=t$ for all $\omega\in\Gamma$\textquotedblright). Thus, $F\left(
Q_{i}\right)  $ is the space of fixed points\footnote{Recall the definition of
a space of fixed points: If a $\mathbf{k}$-module $V$ is equipped with a
linear right action of a group $G$ (that is, if $V$ is a right $\mathbf{k}%
\left[  G\right]  $-module), then its \emph{space of fixed points} is defined
to be the set $\left\{  a\in V\ \mid\ ag=a\text{ for all }g\in G\right\}  $.
This is a $\mathbf{k}$-submodule of $V$.} of the right $\Gamma$-action on
$\mathcal{A}$.

However, we know from the basic theory of group actions (see, e.g.,
\cite[\S 3.3.1, \textquotedblleft Invariants of Permutation
Representations\textquotedblright]{Lorenz18} or \cite[Proposition
A.2]{GriPar25}) that when a finite group $G$ acts on a set $X$, the space of
fixed points of the corresponding permutation module is spanned by the orbit
sums\footnote{In more details:
\par
Let $G$ be a finite group. Let $X$ be a right $G$-set. Consider the
corresponding permutation module $\mathbf{k}^{\left(  X\right)  }$, with its
right $G$-action.
\par
For each $G$-orbit $\mathcal{O}$ on $X$, we define the \emph{orbit sum}
$z_{\mathcal{O}}:=\sum_{x\in\mathcal{O}}x\in\mathbf{k}^{\left(  X\right)  }$.
Now, the known fact that we are citing here is saying that these orbit sums
$z_{\mathcal{O}}$ (as $\mathcal{O}$ ranges over all $G$-orbits on $X$) form a
basis of the space of fixed points of $\mathbf{k}^{\left(  X\right)  }$ (as a
$\mathbf{k}$-module).
\par
In \cite[\S 3.3.1, \textquotedblleft Invariants of Permutation
Representations\textquotedblright]{Lorenz18} and in \cite[Proposition
A.2]{GriPar25}, this is stated for left $G$-actions, but the case of right
$G$-actions is analogous.}. Hence, the $\mathbf{k}$-module $F\left(
Q_{i}\right)  $ is spanned by the orbit sums of the right $\Gamma$-action on
$S_{n}$ (since $F\left(  Q_{i}\right)  $ is the set of fixed points of the
right $\Gamma$-action on $\mathcal{A}$, which is the permutation module
corresponding to the right $\Gamma$-action on $S_{n}$). In other words,
$F\left(  Q_{i}\right)  $ is spanned by the orbit sums $\sum_{\sigma\in
\tau\Gamma}\sigma$ for $\tau\in S_{n}$ (since each orbit of the right $\Gamma
$-action on $S_{n}$ has the form $\tau\Gamma$ for some $\tau\in S_{n}$). As a
left $\mathcal{A}$-module, $F\left(  Q_{i}\right)  $ is therefore generated by
any \textbf{one} of these orbit sums (since any two orbit sums $\sum
_{\sigma\in\tau_{1}\Gamma}\sigma$ and $\sum_{\sigma\in\tau_{2}\Gamma}\sigma$
can be transformed into each other by left multiplication by $\tau_{1}\tau
_{2}^{-1}\in S_{n}\subseteq\mathcal{A}$, and therefore each of them generates
the other).

Now, let $\mathbf{p}=\left(  p_{2},p_{3},\ldots,p_{m}\right)  \in J_{2}\times
J_{3}\times\cdots\times J_{m}$. We shall now show that $\nabla_{\mathbf{p}}$
is one of these orbit sums we just mentioned. Indeed, let $\Omega_{\mathbf{p}%
}$ be the set of all permutations $\sigma\in S_{n}$ that satisfy
\textquotedblleft$\sigma\left(  J_{k}\right)  =J_{k}$ for each $k\in\left[
m\right]  $\textquotedblright\ and \textquotedblleft$\sigma\left(
i_{k-1}\right)  =p_{k}$ for each $k\in\left[  2,m\right]  $\textquotedblright.
Then, the definition of $\nabla_{\mathbf{p}}$ can be rewritten as%
\begin{equation}
\nabla_{\mathbf{p}}=\sum_{\sigma\in\Omega_{\mathbf{p}}}\sigma
.\label{pf.lem.Fi/Fi-1.L3.nap=}%
\end{equation}
We shall now show that $\Omega_{\mathbf{p}}$ is an orbit of the right $\Gamma
$-action on $S_{n}$ (that is, a left coset of $\Gamma$ in $S_{n}$).

First, we show that the set $\Omega_{\mathbf{p}}$ is nonempty. Indeed, it is
easy to construct some permutation $\tau\in\Omega_{\mathbf{p}}$: Namely, we
pick a permutation $\tau_{1}\in S_{J_{1}}$ arbitrarily. Furthermore, for each
$k\in\left[  2,m\right]  $, we pick a permutation $\tau_{k}\in S_{J_{k}}$ that
sends $i_{k-1}\in J_{k}$ to $p_{k}\in J_{k}$. The $m$-tuple $\left(  \tau
_{1},\tau_{2},\ldots,\tau_{m}\right)  $ then belongs to $S_{J_{1}}\times
S_{J_{2}}\times\cdots\times S_{J_{m}}$ and -- viewed as an element of $S_{n}$
via the embedding $S_{J_{1}}\times S_{J_{2}}\times\cdots\times S_{J_{m}%
}\rightarrow S_{n}$ -- belongs to $\Omega_{\mathbf{p}}$.

Hence, $\Omega_{\mathbf{p}}$ is nonempty. Pick any $\tau\in\Omega_{\mathbf{p}%
}$. Then, $\tau\left(  J_{k}\right)  =J_{k}$ for each $k\in\left[  m\right]
$, and $\tau\left(  i_{k-1}\right)  =p_{k}$ for each $k\in\left[  2,m\right]
$. Moreover, these equalities remain valid if we replace $\tau$ by $\tau
\omega$ for any $\omega\in\Gamma$ (because every permutation $\omega\in\Gamma$
preserves the sets $J_{1},J_{2},\ldots,J_{m}$ as well as the elements
$i_{1},i_{2},\ldots,i_{m-1}$). Thus, for each $\omega\in\Gamma$, we have
$\tau\omega\in\Omega_{\mathbf{p}}$ as well. In other words, $\tau
\Gamma\subseteq\Omega_{\mathbf{p}}$.

Conversely, we claim that $\Omega_{\mathbf{p}}\subseteq\tau\Gamma$. Indeed,
let $\sigma\in\Omega_{\mathbf{p}}$ be arbitrary. Then, each $k\in\left[
m\right]  $ satisfies $\sigma\left(  J_{k}\right)  =J_{k}=\tau\left(
J_{k}\right)  $, whereas each $k\in\left[  2,m\right]  $ satisfies
$\sigma\left(  i_{k-1}\right)  =p_{k}=\tau\left(  i_{k-1}\right)  $. Set
$\omega=\tau^{-1}\sigma\in S_{n}$; thus, each $k\in\left[  m\right]  $
satisfies $\omega\left(  J_{k}\right)  =\tau^{-1}\left(  \sigma\left(
J_{k}\right)  \right)  =J_{k}$ (since we just saw that $\sigma\left(
J_{k}\right)  =\tau\left(  J_{k}\right)  $), and each $k\in\left[  2,m\right]
$ satisfies $\omega\left(  i_{k-1}\right)  =\tau^{-1}\left(  \sigma\left(
i_{k-1}\right)  \right)  =i_{k-1}$ (since we just saw that $\sigma\left(
i_{k-1}\right)  =\tau\left(  i_{k-1}\right)  $). Thus, the permutation
$\omega\in S_{n}$ preserves the intervals $J_{1},J_{2},\ldots,J_{m}$ as well
as the elements $i_{1},i_{2},\ldots,i_{m-1}$. Hence, $\omega\in\Gamma$
(because if some permutation $\omega\in S_{n}$ preserves the intervals
$J_{1},J_{2},\ldots,J_{m}$ as well as the elements $i_{1},i_{2},\ldots
,i_{m-1}$, then $\omega$ must belong to $\Gamma$). Now, from $\omega=\tau
^{-1}\sigma$, we obtain $\sigma=\tau\omega\in\tau\Gamma$ (since $\omega
\in\Gamma$). Forget that we fixed $\sigma$. We thus have proved that
$\sigma\in\tau\Gamma$ for each $\sigma\in\Omega_{\mathbf{p}}$. In other words,
$\Omega_{\mathbf{p}}\subseteq\tau\Gamma$.

Combining this with $\tau\Gamma\subseteq\Omega_{\mathbf{p}}$, we obtain
$\Omega_{\mathbf{p}}=\tau\Gamma$. Hence, $\Omega_{\mathbf{p}}$ is an orbit of
the right $\Gamma$-action on $S_{n}$. Thus, $\sum_{\sigma\in\Omega
_{\mathbf{p}}}\sigma$ is an orbit sum of this action. In view of
(\ref{pf.lem.Fi/Fi-1.L3.nap=}), this means that $\nabla_{\mathbf{p}}$ is an
orbit sum of this action. Hence, as a left $\mathcal{A}$-module, $F\left(
Q_{i}\right)  $ is generated by $\nabla_{\mathbf{p}}$ (since we have shown
that $F\left(  Q_{i}\right)  $ is generated by any \textbf{one} of the orbit
sums). This proves Lemma \ref{lem.Fi/Fi-1.L3} \textbf{(b)}. \medskip

\textbf{(c)} Let
\begin{equation}
\vartheta:=\sum_{p_{\ell}\in J_{\ell}}\nabla_{\left(  p_{2},p_{3},\ldots
,p_{m}\right)  }.\label{pf.lem.Fi/Fi-1.L3.c.1}%
\end{equation}
We then must show that $\vartheta\in F_{i-1}$.

We have%
\[
\vartheta=\sum_{p_{\ell}\in J_{\ell}}\underbrace{\nabla_{\left(  p_{2}%
,p_{3},\ldots,p_{m}\right)  }}_{\substack{\in F\left(  Q_{i}\right)
\\\text{(since Lemma \ref{lem.Fi/Fi-1.L3} \textbf{(b)} shows that }%
\nabla_{\mathbf{q}}\in F\left(  Q_{i}\right)  \\\text{for any }\mathbf{q}\in
J_{2}\times J_{3}\times\cdots\times J_{m}\text{)}}}\in F\left(  Q_{i}\right)
\]
(since $F\left(  Q_{i}\right)  $ is a $\mathbf{k}$-module). On the other hand,%
\begin{align}
\vartheta &  =\sum_{p_{\ell}\in J_{\ell}}\nabla_{\left(  p_{2},p_{3}%
,\ldots,p_{m}\right)  }\nonumber\\
&  =\sum_{p_{\ell}\in J_{\ell}}\ \ \sum_{\substack{\sigma\in S_{n}%
;\\\sigma\left(  J_{k}\right)  =J_{k}\text{ for each }k\in\left[  m\right]
;\\\sigma\left(  i_{k-1}\right)  =p_{k}\text{ for each }k\in\left[
2,m\right]  }}\sigma\ \ \ \ \ \ \ \ \ \ \left(  \text{by the definition of
}\nabla_{\left(  p_{2},p_{3},\ldots,p_{m}\right)  }\right)  \nonumber\\
&  =\sum_{p_{\ell}\in J_{\ell}}\ \ \sum_{\substack{\sigma\in S_{n}%
;\\\sigma\left(  J_{k}\right)  =J_{k}\text{ for each }k\in\left[  m\right]
;\\\sigma\left(  i_{k-1}\right)  =p_{k}\text{ for each }k\in\left[
2,m\right]  \setminus\left\{  \ell\right\}  ;\\\sigma\left(  i_{\ell
-1}\right)  =p_{\ell}}}\sigma\nonumber\\
&  \ \ \ \ \ \ \ \ \ \ \ \ \ \ \ \ \ \ \ \ \left(
\begin{array}
[c]{c}%
\text{here, we have split up the}\\
\text{condition \textquotedblleft}\sigma\left(  i_{k-1}\right)  =p_{k}\text{
for each }k\in\left[  2,m\right]  \text{\textquotedblright}\\
\text{under the second summation sign}\\
\text{into two: one for }k\neq\ell\text{ and one for }k=\ell
\end{array}
\right)  \nonumber\\
&  =\sum_{\substack{\sigma\in S_{n};\\\sigma\left(  J_{k}\right)  =J_{k}\text{
for each }k\in\left[  m\right]  ;\\\sigma\left(  i_{k-1}\right)  =p_{k}\text{
for each }k\in\left[  2,m\right]  \setminus\left\{  \ell\right\}
;\\\sigma\left(  i_{\ell-1}\right)  \in J_{\ell}}}\sigma
\label{pf.lem.Fi/Fi-1.L3.c.5}%
\end{align}
(here, we have subsumed the two summation signs into one by removing the
variable $p_{\ell}$). The condition \textquotedblleft$\sigma\left(  i_{\ell
-1}\right)  \in J_{\ell}$\textquotedblright\ under the summation sign in
(\ref{pf.lem.Fi/Fi-1.L3.c.5}) is redundant, since it follows from the
condition \textquotedblleft$\sigma\left(  J_{k}\right)  =J_{k}$ for each
$k\in\left[  m\right]  $\textquotedblright\ (indeed, the latter condition
implies that $\sigma\left(  J_{\ell}\right)  =J_{\ell}$ and therefore
$\sigma\left(  \underbrace{i_{\ell-1}}_{\in J_{\ell}}\right)  \in\sigma\left(
J_{\ell}\right)  =J_{\ell}$). Hence, we can remove this condition. Thus,
(\ref{pf.lem.Fi/Fi-1.L3.c.5}) rewrites as
\begin{equation}
\vartheta=\sum_{\substack{\sigma\in S_{n};\\\sigma\left(  J_{k}\right)
=J_{k}\text{ for each }k\in\left[  m\right]  ;\\\sigma\left(  i_{k-1}\right)
=p_{k}\text{ for each }k\in\left[  2,m\right]  \setminus\left\{  \ell\right\}
}}\sigma.\label{pf.lem.Fi/Fi-1.L3.c.2}%
\end{equation}
However, the two conditions \textquotedblleft$\sigma\left(  J_{k}\right)
=J_{k}$ for each $k\in\left[  m\right]  $\textquotedblright\ and
\textquotedblleft$\sigma\left(  i_{k-1}\right)  =p_{k}$ for each $k\in\left[
2,m\right]  \setminus\left\{  \ell\right\}  $\textquotedblright\ under the
summation sign in (\ref{pf.lem.Fi/Fi-1.L3.c.2}) remain unchanged if we replace
$\sigma$ by $\sigma s_{i_{\ell-1}}$ (since this replacement merely swaps the
values of $\sigma$ on $i_{\ell-1}$ and $i_{\ell-1}+1$, but this does not break
any of the two conditions\footnote{Here we use the fact that the two elements
$i_{\ell-1}$ and $i_{\ell-1}+1$ lie in the same $J_{k}$ (namely, in $J_{\ell
}=\left[  i_{\ell-1},\ i_{\ell}-1\right]  $). This is because $Q_{i}$ is
lacunar, so that $i_{\ell-1}<i_{\ell}-1$.}). Hence, the set of the
permutations $\sigma$ over which we sum in (\ref{pf.lem.Fi/Fi-1.L3.c.2}) is
fixed under right multiplication by $s_{i_{\ell-1}}$. Therefore, the whole sum
is fixed under right multiplication by $s_{i_{\ell-1}}$. Because of
(\ref{pf.lem.Fi/Fi-1.L3.c.2}), this shows that $\vartheta s_{i_{\ell-1}%
}=\vartheta$. Combining this with $\vartheta\in F\left(  Q_{i}\right)  $, we
obtain%
\[
\vartheta\in\left\{  t\in F\left(  Q_{i}\right)  \ \mid\ ts_{i_{\ell-1}%
}=t\right\}  \subseteq F_{i-1}%
\]
(by Lemma \ref{lem.Fi/Fi-1.L2}, applied to $k=\ell-1$). This proves Lemma
\ref{lem.Fi/Fi-1.L3} \textbf{(c)}.
\end{proof}

\subsection{Linear algebra lemmas}

We shall furthermore use two facts from linear algebra over any commutative
ring $\mathbf{k}$:

\begin{lemma}
\label{lem.surj=bij}Let $s\in\mathbb{N}$. Let $M$ and $N$ be two free
$\mathbf{k}$-modules of rank $s$. Then, any surjective $\mathbf{k}$-linear map
$\rho:M\rightarrow N$ is an isomorphism.
\end{lemma}

\begin{proof}
This is a well-known folklore result, and follows easily from the known fact
(\textquotedblleft Orzech's theorem\textquotedblright\ in one of its simplest
forms -- see, e.g., \cite[Exercise 2.5.18 (a)]{GriRei}, or \cite[Corollary
0.2]{Grinberg-Orzech} for a more general result) that any surjective
endomorphism of a free $\mathbf{k}$-module of finite rank is an isomorphism.
For the sake of self-containedness, let me nevertheless give a direct proof:

Let $\rho:M\rightarrow N$ be a surjective $\mathbf{k}$-linear map. We must
show that $\rho$ is an isomorphism.

Pick bases $\left(  e_{1},e_{2},\ldots,e_{s}\right)  $ and $\left(
f_{1},f_{2},\ldots,f_{s}\right)  $ of $M$ and $N$ (these exist, since $M$ and
$N$ are free of rank $s$). The surjectivity of $\rho$ then shows that every
basis vector $f_{i}$ of $N$ lies in the image of $\rho$. That is, every
$f_{i}$ can be written as $\rho\left(  g_{i}\right)  $ for some vector
$g_{i}\in M$. Choose such vectors $g_{i}$, and let $\tau:N\rightarrow M$ be
the $\mathbf{k}$-linear map that sends the basis vectors $f_{1},f_{2}%
,\ldots,f_{s}$ to $g_{1},g_{2},\ldots,g_{s}$, respectively. Then, the
composition $\rho\circ\tau:N\rightarrow N$ sends each vector $f_{i}$ to
$f_{i}$ (since $f_{i}\overset{\tau}{\mapsto}g_{i}\overset{\rho}{\mapsto}%
\rho\left(  g_{i}\right)  =f_{i}$), and thus is the identity map
$\operatorname*{id}\nolimits_{N}$ (since $f_{1},f_{2},\ldots,f_{s}$ form a
basis of $N$).

Now, let $A\in\mathbf{k}^{s\times s}$ be the matrix that represents the linear
map $\rho:M\rightarrow N$ with respect to our bases of $M$ and $N$. Likewise,
let $B\in\mathbf{k}^{s\times s}$ be the matrix that represents the linear map
$\tau:N\rightarrow M$ with respect to our bases of $N$ and $M$. Then, $AB$ is
the matrix that represents the linear map $\rho\circ\tau:N\rightarrow N$ with
respect to our basis of $N$. Therefore, $AB$ is the identity matrix $I_{s}$
(since $\rho\circ\tau$ is the identity map $\operatorname*{id}\nolimits_{N}$,
which is represented by the identity matrix $I_{s}$). Hence, $\det\left(
AB\right)  =\det\left(  I_{s}\right)  =1$ and thus $1=\det\left(  AB\right)
=\det A\cdot\det B$. This shows that $\det A$ is invertible (with inverse
$\det B$). Hence, the matrix $A$ is invertible (with inverse $\dfrac{1}{\det
A}\operatorname*{adj}A$, by the well-known identity $\det A\cdot I_{n}%
=A\cdot\operatorname*{adj}A=\operatorname*{adj}A\cdot A$). In other words, the
$\mathbf{k}$-linear map $\rho$ is invertible (since it is represented by the
matrix $A$), thus an isomorphism. This proves Lemma \ref{lem.surj=bij}.
\end{proof}

\begin{lemma}
\label{lem.tensor-of-quots}Let $V_{1},V_{2},\ldots,V_{m}$ be any $\mathbf{k}%
$-modules. For each $\ell\in\left[  m\right]  $, let $W_{\ell}$ be a
$\mathbf{k}$-submodule of $V_{\ell}$. For each $\ell\in\left[  m\right]  $, we
consider the $\mathbf{k}$-submodule
\[
\underbrace{V_{1}\otimes V_{2}\otimes\cdots\otimes W_{\ell}\otimes
\cdots\otimes V_{m}}_{\substack{\text{This means the tensor product }%
V_{1}\otimes V_{2}\otimes\cdots\otimes V_{m}\text{,}\\\text{in which the }%
\ell\text{-th factor is replaced by }W_{\ell}}}\text{ of }V_{1}\otimes
V_{2}\otimes\cdots\otimes V_{m}.
\]


Then, there is a canonical $\mathbf{k}$-module isomorphism%
\begin{align*}
&  \left(  V_{1}\otimes V_{2}\otimes\cdots\otimes V_{m}\right)  \diagup
\sum_{\ell=1}^{m}\ \ \underbrace{\left(  V_{1}\otimes V_{2}\otimes
\cdots\otimes W_{\ell}\otimes\cdots\otimes V_{m}\right)  }%
_{\substack{\text{This means the tensor product }V_{1}\otimes V_{2}%
\otimes\cdots\otimes V_{m}\text{,}\\\text{in which the }\ell\text{-th factor
is replaced by }W_{\ell}}}\\
&  \cong\left(  V_{1}/W_{1}\right)  \otimes\left(  V_{2}/W_{2}\right)
\otimes\cdots\otimes\left(  V_{m}/W_{m}\right)  .
\end{align*}

\end{lemma}

\begin{noncompile}
\textbf{(b)} If $V_{1},V_{2},\ldots,V_{m}$ are furthermore left modules over
some $\mathbf{k}$-algebras $A_{1},A_{2},\ldots,A_{m}$, and if $W_{1}%
,W_{2},\ldots,W_{m}$ are their submodules (i.e., each $W_{\ell}$ is a left
$A_{\ell}$-submodule of $V_{\ell}$), then this isomorphism is a left
$A_{1}\otimes A_{2}\otimes\cdots\otimes A_{m}$-module isomorphism.
\end{noncompile}

\begin{proof}
We construct both the isomorphism and its inverse using the universal
properties of tensor products and quotients:

\begin{itemize}
\item There is a canonical $\mathbf{k}$-linear map%
\[
\Phi:V_{1}\otimes V_{2}\otimes\cdots\otimes V_{m}\rightarrow\left(
V_{1}/W_{1}\right)  \otimes\left(  V_{2}/W_{2}\right)  \otimes\cdots
\otimes\left(  V_{m}/W_{m}\right)  ,
\]
sending each pure tensor $v_{1}\otimes v_{2}\otimes\cdots\otimes v_{m}$ to
$\overline{v_{1}}\otimes\overline{v_{2}}\otimes\cdots\otimes\overline{v_{m}}$.
This $\mathbf{k}$-linear map $\Phi$ is easily seen to vanish on the submodule
$\sum_{\ell=1}^{m}\left(  V_{1}\otimes V_{2}\otimes\cdots\otimes W_{\ell
}\otimes\cdots\otimes V_{m}\right)  $, and thus factors through the quotient
module. Hence, we obtain a $\mathbf{k}$-linear map%
\begin{align*}
\overline{\Phi}  &  :\left(  V_{1}\otimes V_{2}\otimes\cdots\otimes
V_{m}\right)  \diagup\sum_{\ell=1}^{m}\left(  V_{1}\otimes V_{2}\otimes
\cdots\otimes W_{\ell}\otimes\cdots\otimes V_{m}\right) \\
&  \rightarrow\left(  V_{1}/W_{1}\right)  \otimes\left(  V_{2}/W_{2}\right)
\otimes\cdots\otimes\left(  V_{m}/W_{m}\right)
\end{align*}
sending each $\overline{v_{1}\otimes v_{2}\otimes\cdots\otimes v_{m}}$ to
$\overline{v_{1}}\otimes\overline{v_{2}}\otimes\cdots\otimes\overline{v_{m}}$.

\item Conversely, there is a canonical $\mathbf{k}$-linear map%
\begin{align*}
\Psi &  :\left(  V_{1}/W_{1}\right)  \otimes\left(  V_{2}/W_{2}\right)
\otimes\cdots\otimes\left(  V_{m}/W_{m}\right) \\
&  \rightarrow\left(  V_{1}\otimes V_{2}\otimes\cdots\otimes V_{m}\right)
\diagup\sum_{\ell=1}^{m}\left(  V_{1}\otimes V_{2}\otimes\cdots\otimes
W_{\ell}\otimes\cdots\otimes V_{m}\right)
\end{align*}
sending each $\overline{v_{1}}\otimes\overline{v_{2}}\otimes\cdots
\otimes\overline{v_{m}}$ to $\overline{v_{1}\otimes v_{2}\otimes\cdots\otimes
v_{m}}$. To show that this map is well-defined, we need to check that
$\overline{v_{1}\otimes v_{2}\otimes\cdots\otimes v_{m}}$ depends only on the
residue classes $\overline{v_{i}}$ rather than on the $v_{i}$ themselves (this
is easy: replacing $v_{i}$ by $v_{i}^{\prime}$ with $v_{i}-v_{i}^{\prime}\in
W_{i}$ only changes $v_{1}\otimes v_{2}\otimes\cdots\otimes v_{m}$ by an
element of $V_{1}\otimes V_{2}\otimes\cdots\otimes W_{i}\otimes\cdots\otimes
V_{m}$) and that this dependence is multilinear (this is again easy).
\end{itemize}

Clearly, the maps $\overline{\Phi}$ and $\Psi$ are mutually inverse, hence
isomorphisms. Thus, Lemma \ref{lem.tensor-of-quots} is proved.
\end{proof}

\begin{noncompile}
\textbf{(b)} This follows easily from the construction in part \textbf{(a)}.
\end{noncompile}

For our specific needs, we specialize Lemma \ref{lem.tensor-of-quots} to the
case $W_{1}=0$:

\begin{lemma}
\label{lem.tensor-of-quots2}Let $V_{1},V_{2},\ldots,V_{m}$ be any $\mathbf{k}%
$-modules with $m\geq1$. For each $\ell\in\left[  2,m\right]  $, let $W_{\ell
}$ be a $\mathbf{k}$-submodule of $V_{\ell}$. For each $\ell\in\left[
2,m\right]  $, we consider the $\mathbf{k}$-submodule
\[
\underbrace{V_{1}\otimes V_{2}\otimes\cdots\otimes W_{\ell}\otimes
\cdots\otimes V_{m}}_{\substack{\text{This means the tensor product }%
V_{1}\otimes V_{2}\otimes\cdots\otimes V_{m}\text{,}\\\text{in which the }%
\ell\text{-th factor is replaced by }W_{\ell}}}\text{ of }V_{1}\otimes
V_{2}\otimes\cdots\otimes V_{m}.
\]


Then, there is a canonical $\mathbf{k}$-module isomorphism%
\begin{align*}
&  \left(  V_{1}\otimes V_{2}\otimes\cdots\otimes V_{m}\right)  \diagup
\sum_{\ell=2}^{m}\ \ \underbrace{\left(  V_{1}\otimes V_{2}\otimes
\cdots\otimes W_{\ell}\otimes\cdots\otimes V_{m}\right)  }%
_{\substack{\text{This means the tensor product }V_{1}\otimes V_{2}%
\otimes\cdots\otimes V_{m}\text{,}\\\text{in which the }\ell\text{-th factor
is replaced by }W_{\ell}}}\\
&  \cong V_{1}\otimes\left(  V_{2}/W_{2}\right)  \otimes\left(  V_{3}%
/W_{3}\right)  \otimes\cdots\otimes\left(  V_{m}/W_{m}\right)  .
\end{align*}

\end{lemma}

\begin{noncompile}
\textbf{(b)} If $V_{1},V_{2},\ldots,V_{m}$ are furthermore left modules over
some $\mathbf{k}$-algebras $A_{1},A_{2},\ldots,A_{m}$, and if each $W_{\ell}$
with $\ell\in\left[  2,m\right]  $ is a left $A_{\ell}$-submodule of $V_{\ell
}$, then this isomorphism is a left $A_{1}\otimes A_{2}\otimes\cdots\otimes
A_{m}$-module isomorphism.
\end{noncompile}

\begin{proof}
Apply Lemma \ref{lem.tensor-of-quots} to $W_{1}=0$, and observe that
$V_{1}/0\cong V_{1}$.
\end{proof}

\subsection{Proof of Theorem \ref{thm.Fi/Fi-1.as-ind}}

We can now prove Theorem \ref{thm.Fi/Fi-1.as-ind}:

\begin{proof}
[Proof of Theorem \ref{thm.Fi/Fi-1.as-ind}.]We shall use the notations of
Lemma \ref{lem.Fi/Fi-1.L3}. Note that each $k\in\left[  m\right]  $ satisfies
$J_{k}=\left[  i_{k-1},\ i_{k}-1\right]  $ and thus
\begin{equation}
\left\vert J_{k}\right\vert =i_{k}-i_{k-1}=j_{k}.
\label{pf.thm.Fi/Fi-1.as-ind.jj}%
\end{equation}
Explicitly, there is a bijection%
\begin{align}
\left[  j_{k}\right]   &  \rightarrow J_{k},\nonumber\\
x  &  \mapsto i_{k-1}-1+x \label{pf.thm.Fi/Fi-1.as-ind.bij}%
\end{align}
for each $k\in\left[  m\right]  $.

Consider the tensor product $\mathcal{H}_{j_{1}}\otimes\mathcal{N}_{j_{2}%
}\otimes\mathcal{N}_{j_{3}}\otimes\cdots\otimes\mathcal{N}_{j_{m}}$. We recall
that the trivial representation $\mathcal{H}_{j_{1}}=\mathbf{k}$ has a
$1$-element basis $\left(  1\right)  $, while each natural representation
$\mathcal{N}_{j_{k}}$ has basis $\left(  e_{p}\right)  _{p\in\left[
j_{k}\right]  }=\left(  e_{1},e_{2},\ldots,e_{j_{k}}\right)  $. However, by
abuse of notation, we shall rename the latter basis of $\mathcal{N}_{j_{k}}$
as $\left(  e_{p}\right)  _{p\in J_{k}}=\left(  e_{i_{k-1}},e_{i_{k-1}%
+1},\ldots,e_{i_{k}-1}\right)  $ instead (by shifting all subscripts up by
$i_{k-1}-1$, that is, renaming each basis vector $e_{x}$ as $e_{i_{k-1}-1+x}%
$). Note that this can be done because $j_{k}=i_{k}-i_{k-1}$.

Having renamed the basis vectors of the $\mathbf{k}$-module $\mathcal{N}%
_{j_{k}}$, let us also replace the symmetric group $S_{j_{k}}$ acting on this
module accordingly. Namely, we reinterpret the symmetric group $S_{j_{k}}$
acting on $\mathcal{N}_{j_{k}}$ as the symmetric group $S_{J_{k}}$ using the
bijection (\ref{pf.thm.Fi/Fi-1.as-ind.bij}) between the corresponding sets
$\left[  j_{k}\right]  $ and $J_{k}$. Thus, the left action of $S_{j_{k}}$ on
$\mathcal{N}_{j_{k}}$ becomes a left action of $S_{J_{k}}$ instead; it is
still a permutation action (given on our now-renamed basis by the formula
$\sigma e_{p}=e_{\sigma\left(  p\right)  }$ for each $p\in J_{k}$ and
$\sigma\in S_{J_{k}}$). With these reinterpretations, the parabolic embedding
$S_{j_{1}}\times S_{j_{2}}\times\cdots\times S_{j_{m}}\rightarrow S_{n}$
becomes the usual embedding $S_{J_{1}}\times S_{J_{2}}\times\cdots\times
S_{J_{m}}\rightarrow S_{n}$, which simply combines the $m$ permutations
without any need for shifting (i.e., any $m$-tuple $\left(  \sigma_{1}%
,\sigma_{2},\ldots,\sigma_{m}\right)  \in S_{J_{1}}\times S_{J_{2}}%
\times\cdots\times S_{J_{m}}$ is identified with the permutation $\sigma\in
S_{n}$ that sends each element $x\in J_{k}$ to $\sigma_{k}\left(  x\right)  $
for each $k\in\left[  m\right]  $).

For each $\mathbf{p}=\left(  p_{2},p_{3},\ldots,p_{m}\right)  \in J_{2}\times
J_{3}\times\cdots\times J_{m}$, we have%
\begin{align*}
\nabla_{\mathbf{p}}  &  \in F\left(  Q_{i}\right)  \ \ \ \ \ \ \ \ \ \ \left(
\text{by Lemma \ref{lem.Fi/Fi-1.L3} \textbf{(b)}}\right) \\
&  \subseteq F_{i}\ \ \ \ \ \ \ \ \ \ \left(  \text{since }F_{i}=F\left(
Q_{1}\right)  +F\left(  Q_{2}\right)  +\cdots+F\left(  Q_{i}\right)  \right)
\end{align*}
and thus $\overline{\nabla_{\mathbf{p}}}\in F_{i}/F_{i-1}$ (where
$\overline{\nabla_{\mathbf{p}}}$ denotes the residue class of $\nabla
_{\mathbf{p}}\in F_{i}$ in the quotient $F_{i}/F_{i-1}$). Hence, we can define
a $\mathbf{k}$-linear map%
\begin{align*}
\Phi:\mathcal{H}_{j_{1}}\otimes\mathcal{N}_{j_{2}}\otimes\mathcal{N}_{j_{3}%
}\otimes\cdots\otimes\mathcal{N}_{j_{m}}  &  \rightarrow F_{i}/F_{i-1},\\
1\otimes e_{p_{2}}\otimes e_{p_{3}}\otimes\cdots\otimes e_{p_{m}}  &
\mapsto\overline{\nabla_{\mathbf{p}}}\\
\ \ \ \ \ \ \ \ \ \ \text{for any }\mathbf{p}  &  =\left(  p_{2},p_{3}%
,\ldots,p_{m}\right)  \in J_{2}\times J_{3}\times\cdots\times J_{m}.
\end{align*}
(This map is defined by linearity, since the pure tensors of the form
$1\otimes e_{p_{2}}\otimes e_{p_{3}}\otimes\cdots\otimes e_{p_{m}}$ with
$\mathbf{p}=\left(  p_{2},p_{3},\ldots,p_{m}\right)  \in J_{2}\times
J_{3}\times\cdots\times J_{m}$ form a basis of the $\mathbf{k}$-module
$\mathcal{H}_{j_{1}}\otimes\mathcal{N}_{j_{2}}\otimes\mathcal{N}_{j_{3}%
}\otimes\cdots\otimes\mathcal{N}_{j_{m}}$.) Consider this map $\Phi$.

For each $\ell\in\left[  2,m\right]  $, we can consider the $\mathbf{k}%
$-submodule $\mathcal{H}_{j_{1}}\otimes\mathcal{N}_{j_{2}}\otimes
\mathcal{N}_{j_{3}}\otimes\cdots\otimes\mathcal{D}_{j_{\ell}}\otimes
\cdots\otimes\mathcal{N}_{j_{m}}$ of $\mathcal{H}_{j_{1}}\otimes
\mathcal{N}_{j_{2}}\otimes\mathcal{N}_{j_{3}}\otimes\cdots\otimes
\mathcal{N}_{j_{m}}$, in which its $\ell$-th factor $\mathcal{N}_{j_{\ell}}$
is replaced by its submodule $\mathcal{D}_{j_{\ell}}=\left\{  \left(
a,a,\ldots,a\right)  \ \mid\ a\in\mathbf{k}\right\}  $. We claim that the map
$\Phi$ sends this submodule to $0$. Indeed, this submodule is spanned by sums
of the form%
\[
\sum_{p_{\ell}\in J_{\ell}}1\otimes e_{p_{2}}\otimes e_{p_{3}}\otimes
\cdots\otimes e_{p_{m}}%
\]
(for fixed $p_{2},p_{3},\ldots,p_{\ell-1},p_{\ell+1},\ldots,p_{m}$ in the
respective intervals $J_{k}$)\ \ \ \ \footnote{\textit{Proof.} The submodule
$\mathcal{D}_{j_{\ell}}$ is spanned by the single vector
\[
\left(  1,1,\ldots,1\right)  =e_{i_{\ell-1}}+e_{i_{\ell-1}+1}+\cdots
+e_{i_{\ell}-1}=\sum_{p_{\ell}\in J_{\ell}}e_{p_{\ell}},
\]
and thus the tensor product $\mathcal{H}_{j_{1}}\otimes\mathcal{N}_{j_{2}%
}\otimes\mathcal{N}_{j_{3}}\otimes\cdots\otimes\mathcal{D}_{j_{\ell}}%
\otimes\cdots\otimes\mathcal{N}_{j_{m}}$ is spanned by the pure tensors of the
form%
\begin{align*}
&  1\otimes e_{p_{2}}\otimes e_{p_{3}}\otimes\cdots\otimes e_{p_{\ell-1}%
}\otimes\left(  \sum_{p_{\ell}\in J_{\ell}}e_{p_{\ell}}\right)  \otimes
e_{p_{\ell+1}}\otimes\cdots\otimes e_{p_{m}}\\
&  =\sum_{p_{\ell}\in J_{\ell}}1\otimes e_{p_{2}}\otimes e_{p_{3}}%
\otimes\cdots\otimes e_{p_{m}}\ \ \ \ \ \ \ \ \ \ \text{for fixed }p_{2}%
,p_{3},\ldots,p_{\ell-1},p_{\ell+1},\ldots,p_{m}.
\end{align*}
}, and the map $\Phi$ sends such sums to%
\[
\sum_{p_{\ell}\in J_{\ell}}\overline{\nabla_{\left(  p_{2},p_{3},\ldots
,p_{m}\right)  }}=\overline{\sum_{p_{\ell}\in J_{\ell}}\nabla_{\left(
p_{2},p_{3},\ldots,p_{m}\right)  }}=0_{F_{i}/F_{i-1}},
\]
since Lemma \ref{lem.Fi/Fi-1.L3} \textbf{(c)} shows that $\sum_{p_{\ell}\in
J_{\ell}}\nabla_{\left(  p_{2},p_{3},\ldots,p_{m}\right)  }\in F_{i-1}$.

Thus, the $\mathbf{k}$-linear map%
\[
\Phi:\mathcal{H}_{j_{1}}\otimes\mathcal{N}_{j_{2}}\otimes\mathcal{N}_{j_{3}%
}\otimes\cdots\otimes\mathcal{N}_{j_{m}}\rightarrow F_{i}/F_{i-1}%
\]
sends all the $\mathbf{k}$-submodules $\mathcal{H}_{j_{1}}\otimes
\mathcal{N}_{j_{2}}\otimes\mathcal{N}_{j_{3}}\otimes\cdots\otimes
\mathcal{D}_{j_{\ell}}\otimes\cdots\otimes\mathcal{N}_{j_{m}}$ for $\ell
\in\left[  2,m\right]  $ to $0$. By linearity, we can thus conclude that
$\Phi$ also sends their sum \newline$\sum_{\ell=2}^{m}\left(  \mathcal{H}%
_{j_{1}}\otimes\mathcal{N}_{j_{2}}\otimes\mathcal{N}_{j_{3}}\otimes
\cdots\otimes\mathcal{D}_{j_{\ell}}\otimes\cdots\otimes\mathcal{N}_{j_{m}%
}\right)  $ to $0$. Therefore, $\Phi$ factors through the quotient
$\mathbf{k}$-module%
\begin{align*}
&  \left(  \mathcal{H}_{j_{1}}\otimes\mathcal{N}_{j_{2}}\otimes\mathcal{N}%
_{j_{3}}\otimes\cdots\otimes\mathcal{N}_{j_{m}}\right)  \diagup\sum_{\ell
=2}^{m}\left(  \mathcal{H}_{j_{1}}\otimes\mathcal{N}_{j_{2}}\otimes
\mathcal{N}_{j_{3}}\otimes\cdots\otimes\mathcal{D}_{j_{\ell}}\otimes
\cdots\otimes\mathcal{N}_{j_{m}}\right) \\
&  \cong\mathcal{H}_{j_{1}}\otimes\left(  \mathcal{N}_{j_{2}}/\mathcal{D}%
_{j_{2}}\right)  \otimes\left(  \mathcal{N}_{j_{3}}/\mathcal{D}_{j_{3}%
}\right)  \otimes\cdots\otimes\left(  \mathcal{N}_{j_{m}}/\mathcal{D}_{j_{m}%
}\right) \\
&  \ \ \ \ \ \ \ \ \ \ \ \ \ \ \ \ \ \ \ \ \left(
\begin{array}
[c]{c}%
\text{by Lemma \ref{lem.tensor-of-quots2}, applied to }V_{1}=\mathcal{H}%
_{j_{1}}\text{ and }V_{\ell}=\mathcal{N}_{j_{\ell}}\text{ for }\ell>1\\
\text{and }W_{\ell}=\mathcal{D}_{j_{\ell}}\text{ for }\ell>1
\end{array}
\right) \\
&  =\mathcal{H}_{j_{1}}\otimes\mathcal{Z}_{j_{2}}\otimes\mathcal{Z}_{j_{3}%
}\otimes\cdots\otimes\mathcal{Z}_{j_{m}}\ \ \ \ \ \ \ \ \ \ \left(
\text{since }\mathcal{N}_{p}/\mathcal{D}_{p}=\mathcal{Z}_{p}\text{ for each
}p>0\right)  .
\end{align*}
Thus, we obtain a $\mathbf{k}$-linear map%
\begin{align*}
\overline{\Phi}:\mathcal{H}_{j_{1}}\otimes\mathcal{Z}_{j_{2}}\otimes
\mathcal{Z}_{j_{3}}\otimes\cdots\otimes\mathcal{Z}_{j_{m}}  &  \rightarrow
F_{i}/F_{i-1},\\
1\otimes\overline{e_{p_{2}}}\otimes\overline{e_{p_{3}}}\otimes\cdots
\otimes\overline{e_{p_{m}}}  &  \mapsto\overline{\nabla_{\mathbf{p}}}\\
\ \ \ \ \ \ \ \ \ \ \text{for any }\mathbf{p}  &  =\left(  p_{2},p_{3}%
,\ldots,p_{m}\right)  \in J_{2}\times J_{3}\times\cdots\times J_{m}.
\end{align*}
Consider this map $\overline{\Phi}$. Using Lemma \ref{lem.Fi/Fi-1.L3}
\textbf{(a)}, it is easy to see that this map $\overline{\Phi}$ is $S_{j_{1}%
}\times S_{j_{2}}\times\cdots\times S_{j_{m}}$%
-equivariant\footnote{\textit{Proof.} Let $\tau=\left(  \tau_{1},\tau
_{2},\ldots,\tau_{m}\right)  \in S_{j_{1}}\times S_{j_{2}}\times\cdots\times
S_{j_{m}}$ be any $m$-tuple, and let $\mathbf{p}=\left(  p_{2},p_{3}%
,\ldots,p_{m}\right)  \in J_{2}\times J_{3}\times\cdots\times J_{m}$. We shall
show that%
\[
\overline{\Phi}\left(  \tau\cdot\left(  1\otimes\overline{e_{p_{2}}}%
\otimes\overline{e_{p_{3}}}\otimes\cdots\otimes\overline{e_{p_{m}}}\right)
\right)  =\tau\cdot\overline{\Phi}\left(  1\otimes\overline{e_{p_{2}}}%
\otimes\overline{e_{p_{3}}}\otimes\cdots\otimes\overline{e_{p_{m}}}\right)  .
\]
By linearity, this will entail that the map $\overline{\Phi}$ is $S_{j_{1}%
}\times S_{j_{2}}\times\cdots\times S_{j_{m}}$-equivariant (since elements of
the form $1\otimes\overline{e_{p_{2}}}\otimes\overline{e_{p_{3}}}\otimes
\cdots\otimes\overline{e_{p_{m}}}$ span $\mathcal{H}_{j_{1}}\otimes
\mathcal{Z}_{j_{2}}\otimes\mathcal{Z}_{j_{3}}\otimes\cdots\otimes
\mathcal{Z}_{j_{m}}$).
\par
Indeed, as we mentioned at the beginning of our proof, we regard each
$S_{j_{k}}$ as $S_{J_{k}}$, so that the permutations $\tau_{1},\tau_{2}%
,\ldots,\tau_{m}$ act not on the sets $\left[  j_{1}\right]  ,\left[
j_{2}\right]  ,\ldots,\left[  j_{m}\right]  $ but rather on the sets
$J_{1},J_{2},\ldots,J_{m}$. The embedding of $S_{J_{1}}\times S_{J_{2}}%
\times\cdots\times S_{J_{m}}$ into $S_{n}$ is the usual one, so that our
$m$-tuple $\tau=\left(  \tau_{1},\tau_{2},\ldots,\tau_{m}\right)  $ is equated
with the permutation $\tau\in S_{n}$ given by%
\begin{equation}
\tau\left(  x\right)  =\tau_{k}\left(  x\right)  \ \ \ \ \ \ \ \ \ \ \text{for
each }k\in\left[  m\right]  \text{ and }x\in J_{k}.
\label{pf.thm.Fi/Fi-1.as-ind.fn.six=}%
\end{equation}
As in Lemma \ref{lem.Fi/Fi-1.L3} \textbf{(a)}, we set%
\[
\tau\mathbf{p}:=\left(  \tau_{2}\left(  p_{2}\right)  ,\tau_{3}\left(
p_{3}\right)  ,\ldots,\tau_{m}\left(  p_{m}\right)  \right)  =\left(
\tau\left(  p_{2}\right)  ,\tau\left(  p_{3}\right)  ,\ldots,\tau\left(
p_{m}\right)  \right)  .
\]
\par
Now, we have $\tau=\left(  \tau_{1},\tau_{2},\ldots,\tau_{m}\right)  $ and
thus%
\begin{align*}
\tau\cdot\left(  1\otimes\overline{e_{p_{2}}}\otimes\overline{e_{p_{3}}%
}\otimes\cdots\otimes\overline{e_{p_{m}}}\right)   &  =\tau_{1}1\otimes
\tau_{2}\overline{e_{p_{2}}}\otimes\tau_{3}\overline{e_{p_{3}}}\otimes
\cdots\otimes\tau_{m}\overline{e_{p_{m}}}\\
&  =1\otimes\overline{e_{\tau_{2}\left(  p_{2}\right)  }}\otimes
\overline{e_{\tau_{3}\left(  p_{3}\right)  }}\otimes\cdots\otimes
\overline{e_{\tau_{m}\left(  p_{m}\right)  }}.
\end{align*}
Therefore,%
\begin{align*}
&  \overline{\Phi}\left(  \tau\cdot\left(  1\otimes\overline{e_{p_{2}}}%
\otimes\overline{e_{p_{3}}}\otimes\cdots\otimes\overline{e_{p_{m}}}\right)
\right) \\
&  =\overline{\Phi}\left(  1\otimes\overline{e_{\tau_{2}\left(  p_{2}\right)
}}\otimes\overline{e_{\tau_{3}\left(  p_{3}\right)  }}\otimes\cdots
\otimes\overline{e_{\tau_{m}\left(  p_{m}\right)  }}\right) \\
&  =\overline{\nabla_{\tau\mathbf{p}}}\ \ \ \ \ \ \ \ \ \ \left(  \text{by the
definition of }\overline{\Phi}\text{, since }\tau\mathbf{p}=\left(  \tau
_{2}\left(  p_{2}\right)  ,\tau_{3}\left(  p_{3}\right)  ,\ldots,\tau
_{m}\left(  p_{m}\right)  \right)  \right) \\
&  =\overline{\tau\nabla_{\mathbf{p}}}\ \ \ \ \ \ \ \ \ \ \left(  \text{since
Lemma \ref{lem.Fi/Fi-1.L3} \textbf{(a)} yields }\nabla_{\tau\mathbf{p}}%
=\tau\nabla_{\mathbf{p}}\right) \\
&  =\tau\cdot\overline{\nabla_{\mathbf{p}}}=\tau\cdot\overline{\Phi}\left(
1\otimes\overline{e_{p_{2}}}\otimes\overline{e_{p_{3}}}\otimes\cdots
\otimes\overline{e_{p_{m}}}\right)
\end{align*}
(since the definition of $\overline{\Phi}$ yields $\overline{\nabla
_{\mathbf{p}}}=\overline{\Phi}\left(  1\otimes\overline{e_{p_{2}}}%
\otimes\overline{e_{p_{3}}}\otimes\cdots\otimes\overline{e_{p_{m}}}\right)
$). This is precisely what we wanted to show. Hence, we have proved that the
map $\overline{\Phi}$ is $S_{j_{1}}\times S_{j_{2}}\times\cdots\times
S_{j_{m}}$-equivariant.}, and thus is a left $\mathbf{k}\left[  S_{j_{1}%
}\times S_{j_{2}}\times\cdots\times S_{j_{m}}\right]  $-module morphism.

But the definition of an induction product yields%
\begin{align*}
&  \mathcal{H}_{j_{1}}\ast\mathcal{Z}_{j_{2}}\ast\mathcal{Z}_{j_{3}}\ast
\cdots\ast\mathcal{Z}_{j_{m}}\\
&  =\operatorname*{Ind}\nolimits_{S_{j_{1}}\times S_{j_{2}}\times\cdots\times
S_{j_{m}}}^{S_{n}}\left(  \mathcal{H}_{j_{1}}\otimes\mathcal{Z}_{j_{2}}%
\otimes\mathcal{Z}_{j_{3}}\otimes\cdots\otimes\mathcal{Z}_{j_{m}}\right) \\
&  =\underbrace{\mathbf{k}\left[  S_{n}\right]  }_{=\mathcal{A}}%
\otimes_{\mathbf{k}\left[  S_{j_{1}}\times S_{j_{2}}\times\cdots\times
S_{j_{m}}\right]  }\left(  \mathcal{H}_{j_{1}}\otimes\mathcal{Z}_{j_{2}%
}\otimes\mathcal{Z}_{j_{3}}\otimes\cdots\otimes\mathcal{Z}_{j_{m}}\right)
\ \ \ \ \ \ \ \ \ \ \left(  \text{by (\ref{eq.IndGH.def})}\right) \\
&  =\mathcal{A}\otimes_{\mathbf{k}\left[  S_{j_{1}}\times S_{j_{2}}%
\times\cdots\times S_{j_{m}}\right]  }\left(  \mathcal{H}_{j_{1}}%
\otimes\mathcal{Z}_{j_{2}}\otimes\mathcal{Z}_{j_{3}}\otimes\cdots
\otimes\mathcal{Z}_{j_{m}}\right)  .
\end{align*}
Hence, we can define a left $\mathcal{A}$-module morphism%
\begin{align*}
\Psi:\mathcal{H}_{j_{1}}\ast\mathcal{Z}_{j_{2}}\ast\mathcal{Z}_{j_{3}}%
\ast\cdots\ast\mathcal{Z}_{j_{m}}  &  \rightarrow F_{i}/F_{i-1},\\
a\otimes_{\mathbf{k}\left[  S_{j_{1}}\times S_{j_{2}}\times\cdots\times
S_{j_{m}}\right]  }v  &  \mapsto a\cdot\overline{\Phi}\left(  v\right)
\end{align*}
(this is well-defined, since $\overline{\Phi}:\mathcal{H}_{j_{1}}%
\otimes\mathcal{Z}_{j_{2}}\otimes\mathcal{Z}_{j_{3}}\otimes\cdots
\otimes\mathcal{Z}_{j_{m}}\rightarrow F_{i}/F_{i-1}$ is a left $\mathbf{k}%
\left[  S_{j_{1}}\times S_{j_{2}}\times\cdots\times S_{j_{m}}\right]  $-module
morphism). Explicitly, $\Psi$ is given by%
\begin{align*}
&  \Psi\left(  a\otimes_{\mathbf{k}\left[  S_{j_{1}}\times S_{j_{2}}%
\times\cdots\times S_{j_{m}}\right]  }\left(  1\otimes\overline{e_{p_{2}}%
}\otimes\overline{e_{p_{3}}}\otimes\cdots\otimes\overline{e_{p_{m}}}\right)
\right) \\
&  =a\cdot\overline{\Phi}\left(  1\otimes\overline{e_{p_{2}}}\otimes
\overline{e_{p_{3}}}\otimes\cdots\otimes\overline{e_{p_{m}}}\right) \\
&  =a\cdot\overline{\nabla_{\mathbf{p}}}\ \ \ \ \ \ \ \ \ \ \text{for any
}a\in\mathcal{A}\text{ and }\mathbf{p}=\left(  p_{2},p_{3},\ldots
,p_{m}\right)  \in J_{2}\times J_{3}\times\cdots\times J_{m}%
\end{align*}
(by the definition of $\overline{\Phi}$). Hence, using Lemma
\ref{lem.Fi/Fi-1.L3} \textbf{(b)}, it is easy to see that this map $\Psi$ is
surjective\footnote{\textit{Proof.} The map $\Psi$ is left $\mathcal{A}%
$-linear. Hence, its image is a left $\mathcal{A}$-submodule of $F_{i}%
/F_{i-1}$.
\par
By the definition of $F_{i}$, we have%
\begin{align*}
F_{i}  &  =F\left(  Q_{1}\right)  +F\left(  Q_{2}\right)  +\cdots+F\left(
Q_{i}\right) \\
&  =\underbrace{F\left(  Q_{1}\right)  +F\left(  Q_{2}\right)  +\cdots
+F\left(  Q_{i-1}\right)  }_{=F_{i-1}}+\,F\left(  Q_{i}\right)  =F_{i-1}%
+F\left(  Q_{i}\right)  .
\end{align*}
Hence, the composition of canonical maps
\begin{equation}
F\left(  Q_{i}\right)  \overset{\text{inclusion}}{\longrightarrow}%
F_{i}\overset{\text{projection}}{\longrightarrow}F_{i}/F_{i-1}
\label{pf.thm.Fi/Fi-1.as-ind.fn4.comp}%
\end{equation}
is surjective.
\par
But Lemma \ref{lem.Fi/Fi-1.L3} \textbf{(b)} shows that the left $\mathcal{A}%
$-module $F\left(  Q_{i}\right)  $ is generated by a single element of the
form $\nabla_{\mathbf{p}}$. Hence, the quotient $\mathcal{A}$-module
$F_{i}/F_{i-1}$ is generated by a single element of the form $\overline
{\nabla_{\mathbf{p}}}$ (since the map (\ref{pf.thm.Fi/Fi-1.as-ind.fn4.comp})
is surjective). But any such element of the form $\overline{\nabla
_{\mathbf{p}}}$ lies in the image of $\Psi$ (since we have $\overline
{\nabla_{\mathbf{p}}}=\Psi\left(  1\otimes_{\mathbf{k}\left[  S_{j_{1}}\times
S_{j_{2}}\times\cdots\times S_{j_{m}}\right]  }\left(  1\otimes\overline
{e_{p_{2}}}\otimes\overline{e_{p_{3}}}\otimes\cdots\otimes\overline{e_{p_{m}}%
}\right)  \right)  $ when $\mathbf{p}=\left(  p_{2},p_{3},\ldots,p_{m}\right)
$). Thus, the image of $\Psi$ must contain a generator of $F_{i}/F_{i-1}$, and
thus must be the entire $\mathcal{A}$-module $F_{i}/F_{i-1}$ (since this image
is a left $\mathcal{A}$-submodule of $F_{i}/F_{i-1}$). In other words, $\Psi$
is surjective.}.

We now know that $\Psi$ is a surjective left $\mathcal{A}$-module morphism
from $\mathcal{H}_{j_{1}}\ast\mathcal{Z}_{j_{2}}\ast\mathcal{Z}_{j_{3}}%
\ast\cdots\ast\mathcal{Z}_{j_{m}}$ to $F_{i}/F_{i-1}$. We shall now show that
$\Psi$ is an isomorphism.

From Lemma \ref{lem.Fi/Fi-1.dim}, we know that the $\mathbf{k}$-module
$F_{i}/F_{i-1}$ is free of rank%
\[
\dfrac{n!}{j_{1}!j_{2}!\cdots j_{m}!}\cdot\prod_{k=2}^{m}\left(
j_{k}-1\right)  .
\]
But the $\mathbf{k}$-module
\[
\mathcal{H}_{j_{1}}\ast\mathcal{Z}_{j_{2}}\ast\mathcal{Z}_{j_{3}}\ast
\cdots\ast\mathcal{Z}_{j_{m}}=\operatorname*{Ind}\nolimits_{S_{j_{1}}\times
S_{j_{2}}\times\cdots\times S_{j_{m}}}^{S_{n}}\left(  \mathcal{H}_{j_{1}%
}\otimes\mathcal{Z}_{j_{2}}\otimes\mathcal{Z}_{j_{3}}\otimes\cdots
\otimes\mathcal{Z}_{j_{m}}\right)
\]
is also free of rank\footnote{Here, we are denoting the rank of a free
$\mathbf{k}$-module $V$ by $\dim V$, and we are using the fact that an induced
representation $\operatorname*{Ind}\nolimits_{H}^{G}V$ is free of rank
$\dfrac{\left\vert G\right\vert }{\left\vert H\right\vert }\cdot\dim V$ (as a
$\mathbf{k}$-module) whenever $V$ is free (as a $\mathbf{k}$-module). (The
latter fact is an easy consequence of the fact that $\mathbf{k}\left[
G\right]  $ is a free right $\mathbf{k}\left[  H\right]  $-module of rank
$\dfrac{\left\vert G\right\vert }{\left\vert H\right\vert }$.)}%
\begin{align*}
&  \underbrace{\dfrac{\left\vert S_{n}\right\vert }{\left\vert S_{j_{1}}\times
S_{j_{2}}\times\cdots\times S_{j_{m}}\right\vert }}_{=\dfrac{n!}{j_{1}%
!j_{2}!\cdots j_{m}!}}\cdot\underbrace{\dim\left(  \mathcal{H}_{j_{1}}%
\otimes\mathcal{Z}_{j_{2}}\otimes\mathcal{Z}_{j_{3}}\otimes\cdots
\otimes\mathcal{Z}_{j_{m}}\right)  }_{\substack{=\dim\left(  \mathcal{H}%
_{j_{1}}\right)  \cdot\prod_{k=2}^{m}\dim\left(  \mathcal{Z}_{j_{k}}\right)
\\=1\cdot\prod_{k=2}^{m}\left(  j_{k}-1\right)  \\\text{(since }%
\mathcal{H}_{j_{1}}\text{ is free of rank }1\text{,}\\\text{whereas each
}\mathcal{Z}_{j_{k}}\text{ is free of rank }j_{k}-1\text{)}}}\\
&  =\dfrac{n!}{j_{1}!j_{2}!\cdots j_{m}!}\cdot\prod_{k=2}^{m}\left(
j_{k}-1\right)  .
\end{align*}
Thus, Lemma \ref{lem.surj=bij} (applied to $M=\mathcal{H}_{j_{1}}%
\ast\mathcal{Z}_{j_{2}}\ast\mathcal{Z}_{j_{3}}\ast\cdots\ast\mathcal{Z}%
_{j_{m}}$ and $N=F_{i}/F_{i-1}$ and $s=\dfrac{n!}{j_{1}!j_{2}!\cdots j_{m}%
!}\cdot\prod_{k=2}^{m}\left(  j_{k}-1\right)  $ and $\rho=\Psi$) shows that
the surjective $\mathbf{k}$-linear map $\Psi:\mathcal{H}_{j_{1}}%
\ast\mathcal{Z}_{j_{2}}\ast\mathcal{Z}_{j_{3}}\ast\cdots\ast\mathcal{Z}%
_{j_{m}}\rightarrow F_{i}/F_{i-1}$ must be an isomorphism. Since $\Psi$ is a
left $\mathcal{A}$-module morphism, we thus conclude that $\Psi$ is a left
$\mathcal{A}$-module isomorphism. Therefore, $F_{i}/F_{i-1}\cong%
\mathcal{H}_{j_{1}}\ast\mathcal{Z}_{j_{2}}\ast\mathcal{Z}_{j_{3}}\ast
\cdots\ast\mathcal{Z}_{j_{m}}$ as left $\mathcal{A}$-modules, i.e., as $S_{n}%
$-representations. Hence, Theorem \ref{thm.Fi/Fi-1.as-ind} is proved.
\end{proof}

\subsection{In terms of Littlewood--Richardson coefficients}

In the characteristic-$0$ case, we can restate the claim of Theorem
\ref{thm.Fi/Fi-1.as-ind} in terms of Littlewood--Richardson coefficients. Let
us first recount the bare minimum of symmetric function theory needed to state this.

We will use standard notations for (integer) partitions; in particular, the
size of a partition $\lambda$ will be denoted by $\left\vert \lambda
\right\vert $. We let $\operatorname*{Par}$ denote the set of all partitions.
We let $\Lambda$ be the ring of symmetric functions over $\mathbb{Z}$ (not
over $\mathbf{k}$); we refer to \cite[\S 2.1]{GriRei} or \cite[\S 4.3]%
{Sagan-SG} for its definition\footnote{Note that \cite[\S 4.3]{Sagan-SG} uses
$\mathbb{C}$ as the base ring, but everything works for any base ring.}. To
each partition $\lambda$ corresponds a special symmetric function $s_{\lambda
}\in\Lambda$ called the \emph{Schur function}; see \cite[(2.2.4)]{GriRei} or
\cite[\S 4.4]{Sagan-SG} or \cite[Definition 5.3]{Egge19} for its definition.
It is well-known (see \cite[(4.26) and Theorem 4.9.4]{Sagan-SG} or
\cite[Definition 2.5.8 and Corollary 2.6.12]{GriRei} or \cite[Theorem
10.40]{Egge19}) that a product $s_{\mu}s_{\nu}$ of two Schur functions (for
$\mu,\nu\in\operatorname*{Par}$) can always be written as an $\mathbb{N}%
$-linear combination of Schur functions -- i.e., there exist coefficients
$c_{\mu,\nu}^{\lambda}\in\mathbb{N}$ for all $\lambda,\mu,\nu\in
\operatorname*{Par}$ such that every two partitions $\mu$ and $\nu$ satisfy%
\begin{equation}
s_{\mu}s_{\nu}=\sum_{\lambda\in\operatorname*{Par}}c_{\mu,\nu}^{\lambda
}s_{\lambda}. \label{eq.LR.2}%
\end{equation}
These coefficients $c_{\mu,\nu}^{\lambda}$ are known as the
\emph{Littlewood--Richardson coefficients}. More generally, if $\mu_{1}%
,\mu_{2},\ldots,\mu_{k}$ are any $k$ partitions, then we can write the product
$s_{\mu_{1}}s_{\mu_{2}}\cdots s_{\mu_{k}}$ in the form%
\begin{equation}
s_{\mu_{1}}s_{\mu_{2}}\cdots s_{\mu_{k}}=\sum_{\lambda\in\operatorname*{Par}%
}c_{\mu_{1},\mu_{2},\ldots,\mu_{k}}^{\lambda}s_{\lambda} \label{eq.LR.k}%
\end{equation}
with coefficients $c_{\mu_{1},\mu_{2},\ldots,\mu_{k}}^{\lambda}\in\mathbb{N}$.
These \textquotedblleft$k$\emph{-Littlewood--Richardson coefficients}%
\textquotedblright\ $c_{\mu_{1},\mu_{2},\ldots,\mu_{k}}^{\lambda}$ are, in
fact, easily computed by recursion using the standard Littlewood--Richardson
coefficients $c_{\mu,\nu}^{\lambda}$: Namely, for $k=0$, we have%
\[
c^{\lambda}=\delta_{\lambda,\varnothing}\ \ \ \ \ \ \ \ \ \ \left(
\text{Kronecker delta}\right)  ;
\]
for $k=1$, we have%
\[
c_{\mu}^{\lambda}=\delta_{\lambda,\mu}\ \ \ \ \ \ \ \ \ \ \left(
\text{Kronecker delta}\right)  ;
\]
and for any higher $k$, we have%
\[
c_{\mu_{1},\mu_{2},\ldots,\mu_{k}}^{\lambda}=\sum_{\nu\in\operatorname*{Par}%
}c_{\mu_{1},\mu_{2},\ldots,\mu_{k-1}}^{\nu}c_{\nu,\mu_{k}}^{\lambda}%
\]
(since the product $s_{\mu_{1}}s_{\mu_{2}}\cdots s_{\mu_{k}}$ can be computed
as $\left(  s_{\mu_{1}}s_{\mu_{2}}\cdots s_{\mu_{k-1}}\right)  s_{\mu_{k}}$).

Note that any Schur function $s_{\lambda}$ is homogeneous of degree
$\left\vert \lambda\right\vert $. Hence, a Littlewood--Richardson coefficient
$c_{\mu,\nu}^{\lambda}$ is always $0$ unless $\left\vert \lambda\right\vert
=\left\vert \mu\right\vert +\left\vert \nu\right\vert $. Thus, we can rewrite
the equality (\ref{eq.LR.2}) as%
\begin{equation}
s_{\mu}s_{\nu}=\sum_{\substack{\lambda\in\operatorname*{Par};\\\left\vert
\lambda\right\vert =\left\vert \mu\right\vert +\left\vert \nu\right\vert
}}c_{\mu,\nu}^{\lambda}s_{\lambda}. \label{eq.LR.2n}%
\end{equation}
Likewise, we can rewrite (\ref{eq.LR.k}) as%
\begin{equation}
s_{\mu_{1}}s_{\mu_{2}}\cdots s_{\mu_{k}}=\sum_{\substack{\lambda
\in\operatorname*{Par};\\\left\vert \lambda\right\vert =\left\vert \mu
_{1}\right\vert +\left\vert \mu_{2}\right\vert +\cdots+\left\vert \mu
_{k}\right\vert }}c_{\mu_{1},\mu_{2},\ldots,\mu_{k}}^{\lambda}s_{\lambda}.
\label{eq.LR.kn}%
\end{equation}


We note that there is a second Littlewood--Richardson rule (\cite[Theorem
4.9.2]{Sagan-SG}, \cite[(2.6.4)]{GriRei}, \cite[Theorem 10.40]{Egge19}) that
decomposes a skew Schur function $s_{\lambda/\mu}$ into an $\mathbb{N}$-linear
combination of (straight) Schur functions $s_{\nu}$ as follows:%
\begin{equation}
s_{\lambda/\mu}=\sum_{\nu\in\operatorname*{Par}}c_{\mu,\nu}^{\lambda}s_{\nu}.
\label{eq.LR.lam/mu}%
\end{equation}
The formula (\ref{eq.LR.k}) can also be viewed as a particular case of that
second rule, since the product $s_{\mu_{1}}s_{\mu_{2}}\cdots s_{\mu_{k}}$ can
be written as the skew Schur function $s_{\mu_{1}\ast\mu_{2}\ast\cdots\ast
\mu_{k}}$ corresponding to the skew shape $\mu_{1}\ast\mu_{2}\ast\cdots\ast
\mu_{k}$ obtained by attaching the Young diagrams of $\mu_{1},\mu_{2}%
,\ldots,\mu_{k}$ to each other along their northeastern/southwestern corners
(see \cite[\S 1]{RemmelWhitney} for the precise definition; this claim follows
from \cite[Proposition 5.9]{Egge19}; cf. also \cite[Figure 7.2]{Stanley-EC2}).
Thus, a $k$-Littlewood--Richardson coefficient $c_{\mu_{1},\mu_{2},\ldots
,\mu_{k}}^{\lambda}$ can actually be rewritten as a (regular)
Littlewood--Richardson coefficient using (\ref{eq.LR.lam/mu}): If we write the
skew shape $\mu_{1}\ast\mu_{2}\ast\cdots\ast\mu_{k}$ as $\alpha/\beta$, then
\begin{equation}
c_{\mu_{1},\mu_{2},\ldots,\mu_{k}}^{\lambda}=c_{\beta,\lambda}^{\alpha}.
\label{eq.LR.c=c}%
\end{equation}


The same Littlewood--Richardson coefficients govern the decomposition of
induction products of Specht modules into Specht modules in characteristic
$0$. Namely, if $\mathbf{k}$ is a field of characteristic $0$, and if $\mu$
and $\nu$ are two partitions of respective sizes $i$ and $j$, then%
\begin{equation}
\mathcal{S}^{\mu}\ast\mathcal{S}^{\nu}\cong\bigoplus_{\substack{\lambda
\in\operatorname*{Par};\\\left\vert \lambda\right\vert =i+j}}\left(
\mathcal{S}^{\lambda}\right)  ^{\oplus c_{\mu,\nu}^{\lambda}}
\label{eq.LR.2Specht}%
\end{equation}
as $S_{i+j}$-modules\footnote{The notation $V^{\oplus k}$ means the direct sum
$V\oplus V\oplus\cdots\oplus V$ of $k$ copies of $V$.}. Indeed, this follows
from the Schur function equality (\ref{eq.LR.2n}) using the Frobenius
characteristic map \cite[Theorem 4.7.4]{Sagan-SG} (in fact, this map -- or,
rather, its inverse -- sends Schur functions $s_{\lambda}$ to Specht modules
$\mathcal{S}^{\lambda}$, while sending products of symmetric functions to
induction products of representations\footnote{For products with two factors,
this is proved in \cite[Theorem 4.7.4]{Sagan-SG} (using characters and
Frobenius reciprocity). For products with $k$ factors, it follows from the
two-factor case using Proposition \ref{prop.indprod.ass}.}). Likewise, if
$\mathbf{k}$ is a field of characteristic $0$, and if $\mu_{1},\mu_{2}%
,\ldots,\mu_{k}$ are $k$ partitions of respective sizes $i_{1},i_{2}%
,\ldots,i_{k}$, then%
\begin{equation}
\mathcal{S}^{\mu_{1}}\ast\mathcal{S}^{\mu_{2}}\ast\cdots\ast\mathcal{S}%
^{\mu_{k}}\cong\bigoplus_{\substack{\lambda\in\operatorname*{Par};\\\left\vert
\lambda\right\vert =i_{1}+i_{2}+\cdots+i_{k}}}\left(  \mathcal{S}^{\lambda
}\right)  ^{\oplus c_{\mu_{1},\mu_{2},\ldots,\mu_{k}}^{\lambda}}
\label{eq.LR.kSpecht}%
\end{equation}
as $S_{i_{1}+i_{2}+\cdots+i_{k}}$-modules. When $\mathbf{k}$ is a field of
characteristic $\neq0$, or, more generally, just any commutative ring, then
the isomorphisms (\ref{eq.LR.2Specht}) and (\ref{eq.LR.kSpecht}) still hold in
a weaker form, where the direct sums are replaced by filtrations whose
subquotients are Specht modules of the form $\mathcal{S}^{\lambda}$ (see
\cite[Theorem 9.7]{Clausen91} for the more general problem of finding such
filtrations of skew Specht modules\footnote{Indeed, the induction product
$\mathcal{S}^{\mu_{1}}\ast\mathcal{S}^{\mu_{2}}\ast\cdots\ast\mathcal{S}%
^{\mu_{k}}$ is easily seen to be isomorphic to the skew Specht module
$\mathcal{S}^{\alpha/\beta}$, where $\alpha/\beta=\mu_{1}\ast\mu_{2}\ast
\cdots\ast\mu_{k}$. Thus, \cite[Theorem 9.7]{Clausen91} (applied to this skew
diagram $\alpha/\beta$) yields a Specht series (i.e., a filtration whose
subquotients are Specht modules of the form $\mathcal{S}^{\lambda}$) for
$\mathcal{S}^{\mu_{1}}\ast\mathcal{S}^{\mu_{2}}\ast\cdots\ast\mathcal{S}%
^{\mu_{k}}$.}).

Now, Theorem \ref{thm.Fi/Fi-1.as-ind} gives rise to the following
decomposition of $F_{i}/F_{i-1}$:

\begin{corollary}
\label{cor.Fi/Fi-1.lr}Assume that $\mathbf{k}$ is a field of characteristic
$0$.

Let $i\in\left[  f_{n+1}\right]  $. Consider the lacunar subset $Q_{i}$ of
$\left[  n-1\right]  $. Write the set $Q_{i}\cup\left\{  n+1\right\}  $ as
$\left\{  i_{1}<i_{2}<\cdots<i_{m}\right\}  $, so that $i_{m}=n+1$.
Furthermore, set $i_{0}:=1$. Set $j_{k}:=i_{k}-i_{k-1}$ for each $k\in\left[
m\right]  $. (Note that $j_{1}\geq0$ and $j_{2},j_{3},\ldots,j_{m}>1$ and
$j_{1}+j_{2}+\cdots+j_{m}=n$; this follows from Lemma \ref{lem.Fi/Fi-1.sum-js}
(applied to $I=Q_{i}$).)

Let $\operatorname*{Par}\nolimits_{n}$ be the set of all partitions of $n$.
Then,%
\[
F_{i}/F_{i-1}\cong\bigoplus_{\lambda\in\operatorname*{Par}\nolimits_{n}%
}\left(  \mathcal{S}^{\lambda}\right)  ^{\oplus c_{\left(  j_{1}\right)
,\ \left(  j_{2}-1,1\right)  ,\ \left(  j_{3}-1,1\right)  ,\ \ldots,\ \left(
j_{m}-1,1\right)  }^{\lambda}}%
\]
(where the subscripts in $c_{\left(  j_{1}\right)  ,\ \left(  j_{2}%
-1,1\right)  ,\ \left(  j_{3}-1,1\right)  ,\ \ldots,\ \left(  j_{m}%
-1,1\right)  }^{\lambda}$ are the partition $\left(  j_{1}\right)  $ followed
by the partitions $\left(  j_{k}-1,1\right)  $ for all $k\in\left[
2,m\right]  $).
\end{corollary}

\begin{proof}
The set $Q_{i}\cup\left\{  n+1\right\}  $ is lacunar (since $Q_{i}$ is
lacunar, and since $Q_{i}\subseteq\left[  n-1\right]  $ ensures that $n+1$ is
larger than any element of $Q_{i}$ by at least $2$). In other words, every
$k>1$ satisfies $i_{k}>i_{k-1}+1$. Hence, every $k>1$ satisfies $j_{k}>1$
(since $j_{k}=i_{k}-i_{k-1}$) and therefore%
\begin{equation}
\mathcal{Z}_{j_{k}}\cong\mathcal{S}^{\left(  j_{k}-1,1\right)  }
\label{pf.cor.Fi/Fi-1.lr.Z}%
\end{equation}
(by Proposition \ref{prop.Z-iso} \textbf{(b)}, since $j_{k}$ is invertible in
$\mathbf{k}$).

Also, we have $\mathcal{H}_{p}\cong\mathcal{S}^{\left(  p\right)  }$ for each
$p\in\mathbb{N}$ (since both $\mathcal{H}_{p}$ and $\mathcal{S}^{\left(
p\right)  }$ are trivial $1$-dimensional representations of $S_{p}$). Thus,
$\mathcal{H}_{j_{1}}\cong\mathcal{S}^{\left(  j_{1}\right)  }$.

Now, Theorem \ref{thm.Fi/Fi-1.as-ind} yields%
\begin{align*}
F_{i}/F_{i-1}  &  \cong\underbrace{\mathcal{H}_{j_{1}}\ast\mathcal{Z}_{j_{2}%
}\ast\mathcal{Z}_{j_{3}}\ast\cdots\ast\mathcal{Z}_{j_{m}}}%
_{\substack{\text{the first factor is an }\mathcal{H}\text{,}\\\text{while all
others are }\mathcal{Z}\text{'s}}}\\
&  \cong\mathcal{S}^{\left(  j_{1}\right)  }\ast\mathcal{S}^{\left(
j_{2}-1,1\right)  }\ast\mathcal{S}^{\left(  j_{3}-1,1\right)  }\ast\cdots
\ast\mathcal{S}^{\left(  j_{m}-1,1\right)  }\\
&  \ \ \ \ \ \ \ \ \ \ \ \ \ \ \ \ \ \ \ \ \left(
\begin{array}
[c]{c}%
\text{since }\mathcal{H}_{j_{1}}\cong\mathcal{S}^{\left(  j_{1}\right)
}\text{, and since }\mathcal{Z}_{j_{k}}\cong\mathcal{S}^{\left(
j_{k}-1,1\right)  }\text{ for each }k>1\\
\text{(by (\ref{pf.cor.Fi/Fi-1.lr.Z}))}%
\end{array}
\right) \\
&  \cong\bigoplus_{\substack{\lambda\in\operatorname*{Par};\\\left\vert
\lambda\right\vert =j_{1}+j_{2}+\cdots+j_{m}}}\left(  \mathcal{S}^{\lambda
}\right)  ^{\oplus c_{\left(  j_{1}\right)  ,\ \left(  j_{2}-1,1\right)
,\ \left(  j_{3}-1,1\right)  ,\ \ldots,\ \left(  j_{m}-1,1\right)  }^{\lambda
}}\ \ \ \ \ \ \ \ \ \ \left(  \text{by (\ref{eq.LR.kSpecht})}\right) \\
&  =\bigoplus_{\substack{\lambda\in\operatorname*{Par};\\\left\vert
\lambda\right\vert =n}}\left(  \mathcal{S}^{\lambda}\right)  ^{\oplus
c_{\left(  j_{1}\right)  ,\ \left(  j_{2}-1,1\right)  ,\ \left(
j_{3}-1,1\right)  ,\ \ldots,\ \left(  j_{m}-1,1\right)  }^{\lambda}%
}\ \ \ \ \ \ \ \ \ \ \left(  \text{since }j_{1}+j_{2}+\cdots+j_{m}=n\right) \\
&  =\bigoplus_{\lambda\in\operatorname*{Par}\nolimits_{n}}\left(
\mathcal{S}^{\lambda}\right)  ^{\oplus c_{\left(  j_{1}\right)  ,\ \left(
j_{2}-1,1\right)  ,\ \left(  j_{3}-1,1\right)  ,\ \ldots,\ \left(
j_{m}-1,1\right)  }^{\lambda}}.
\end{align*}
This proves Corollary \ref{cor.Fi/Fi-1.lr}.
\end{proof}

%\subsection{In terms of the Frobenius characteristic map}


%We shall now restate Corollary \ref{cor.Fi/Fi-1.lr} in terms of symmetric
%functions, at least in the characteristic-$0$ case.


%For the rest of this section, we assume that $\mathbf{k}$ is a field of
%characteristic $0$. We let $\Lambda_{\mathbf{k}}$ denote the ring of symmetric
%functions over $\mathbf{k}$ (defined in \cite[\S 2.1]{GriRei}). For every
%finite-dimensional $S_{n}$-representation $V$ over $\mathbf{k}$, we define%
%\[
%\psi\left(  V\right)  :=\dfrac{1}{n!}\sum_{\sigma\in S_{n}}\chi_{V}\left(
%\sigma\right)  p_{\operatorname*{type}\sigma}\in\Lambda_{\mathbf{k}},
%\]
%where


%\begin{itemize}
%\item $\chi_{V}$ is the character of the representation $V$ (that is,
%$\chi_{V}\left(  a\right)  =\operatorname*{Tr}\left(  a\mid_{V}\right)  $ for
%any $a\in\mathcal{A}=\mathbf{k}\left[  S_{n}\right]  $);


%\item $\operatorname*{type}\sigma$ denotes the cycle type of a permutation
%$\sigma$;


%\item $p_{\mu}$ (where $\mu$ is a partition) is the $\mu$-th power-sum
%symmetric function (see \cite[Definition 2.2.1]{GriRei}).
%\end{itemize}


%What we just have called $\psi\left(  V\right)  $ is denoted
%$\operatorname*{ch}\left(  \chi_{V}\right)  $ in \cite[\S 4.4]{GriRei} (by
%\cite[Exercise 4.4.5(a)]{GriRei}). Our notation $\psi$ is taken from
%\cite[\S 7.3]{Fulton97} (indeed, our above formula $\psi\left(  V\right)
%=\dfrac{1}{n!}\sum_{\sigma\in S_{n}}\chi_{V}\left(  \sigma\right)
%p_{\operatorname*{type}\sigma}$ is easily seen to be equivalent to the formula
%\cite[\S 7.3, (12)]{Fulton97}, because of \cite[\S 7.2, Exercise 2]{Fulton97}).


%\begin{noncompile}
%$\sum_{\mu\vdash n}\dfrac{1}{z\left(  \mu\right)  }\chi_{V}\left(  C\left(
%\mu\right)  \right)  p_{\mu}$


%\textquotedblleft$\mu\vdash n$\textquotedblright\ means \textquotedblleft$\mu$
%is a partition of $n$\textquotedblright;


%$\chi_{V}\left(  C\left(  \mu\right)  \right)  $ denotes the character of $V$
%evaluated on a permutation $\sigma\in S_{n}$ that has cycle type $\mu$;


%$z\left(  \mu\right)  =\prod_{r\geq1}r^{m_{r}}\cdot m_{r}!$, where $m_{r}$ is
%the multiplicity of $r$ in the partition $\mu$.
%\end{noncompile}


%The \emph{representation ring of the symmetric groups} is a free $\mathbf{k}%
%$-module with basis consisting of all the isomorphism classes $\left[
%V\right]  $ of all irreducible $S_{n}$-representations $V$ over $\mathbf{k}$
%for all $n\in\mathbb{N}$. Addition is identified with direct sum, and the
%product is defined by the induction product (that is, $\left[  V\right]
%\cdot\left[  W\right]  :=\left[  V\ast W\right]  $ for all $V,W$).


%We shall use some basic notions from the representation theory of $S_{n}$ and
%from symmetric functions; the reader can find all prerequisites in
%\cite[Chapters 6 and 7]{Fulton97}. For any partition $\lambda$ of $n$, a
%Specht module $\mathcal{S}^{\lambda}$ is defined, which is an irreducible
%representation of $S_{n}$ with a basis indexed by standard tableaux of shape
%$\lambda$. (In \cite{Fulton97}, it is called $S^{\lambda}$.) Each
%$u\in\mathbf{k}\left[  S_{n}\right]  $ acts (on the left) on this Specht
%module $\mathcal{S}^{\lambda}$; we let $L_{\lambda}\left(  u\right)  $ denote
%this action (viewed as a $\mathbf{k}$-module endomorphism of $\mathcal{S}%
%^{\lambda}$).


%We let $\mathcal{R}$ denote the (inductive) representation ring of the
%symmetric groups (called $R$ in \cite[\S 7.3]{Fulton97}), and $\Lambda$ denote
%the ring of symmetric functions over $\mathbb{Z}$ (defined in \cite[\S 6.2]%
%{Fulton97}). The \emph{Frobenius characteristic map} is the ring isomorphism
%$\varphi:\Lambda\rightarrow\mathcal{R}$ defined in \cite[\S 7.3]{Fulton97},
%and the famous \emph{Schur function} $s_{\lambda}\in\Lambda$ corresponding to
%a partition $\lambda$ is the preimage of the Specht module $\mathcal{S}%
%^{\lambda}$ under this isomorphism $\varphi$.


%For each $m\in\mathbb{N}$, we let $h_{m}\in\Lambda$ denote the $m$-th complete
%homogeneous symmetric polynomial. For each $m>0$, we let $z_{m}\in\Lambda$
%denote the Schur function $s_{\left(  m-1,1\right)  }=h_{m-1}h_{1}-h_{m}%
%\in\Lambda$. (By definition, this is $0$ for $m=1$.)


%For each subset $I$ of $\left[  n\right]  $, we define a symmetric function
%$z_{I}:=h_{i_{1}-1}\prod_{j=2}^{k}z_{i_{j}-i_{j-1}}\in\Lambda$, where
%$i_{1},i_{2},\ldots,i_{k}$ are the elements of $I\cup\left\{  n+1\right\}  $
%in increasing order (so that $i_{k}=n+1$ and $I=\left\{  i_{1}<i_{2}%
%<\cdots<i_{k-1}\right\}  $). When this symmetric function $z_{I}$ is expanded
%in the basis $\left(  s_{\lambda}\right)  _{\lambda\text{ is a partition}}$ of
%$\Lambda$, the coefficient of a given Schur function $s_{\lambda}$ shall be
%called $c_{I}^{\lambda}$. This coefficient $c_{I}^{\lambda}$ is actually a
%Littlewood--Richardson coefficient (since $z_{I}$ is a skew Schur function),
%hence $\in\mathbb{N}$.


\section{\label{sec.specht-spec}The Specht module spectrum}

\subsection{\label{subsec.specht-spec.thm}The theorem}

We need a few more notations from \cite{s2b1}. For any subset $I$ of $\left[
n\right]  $, we define the following:

\begin{itemize}
\item We let $\widehat{I}$ be the set $\left\{  0\right\}  \cup I\cup\left\{
n+1\right\}  $. We shall refer to $\widehat{I}$ as the \emph{enclosure} of $I$.

For example, if $n=5$, then $\widehat{\left\{  2,3\right\}  }=\left\{
0,2,3,6\right\}  $.

\item For any $\ell\in\left[  n\right]  $, we let $m_{I,\ell}$ be the number%
\[
\left(  \text{smallest element of }\widehat{I}\text{ that is }\geq\ell\right)
-\ell\in\left[  0,n+1-\ell\right]  \subseteq\left[  0,n\right]  .
\]


For example, if $n=6$ and $I=\left\{  2,5\right\}  $, then%
\[
\left(  m_{I,1},\ m_{I,2},\ m_{I,3},\ m_{I,4},\ m_{I,5},\ m_{I,6}\right)
=\left(  1,\ 0,\ 2,\ 1,\ 0,\ 1\right)  .
\]
We note that an $\ell\in\left[  n\right]  $ satisfies $m_{I,\ell}=0$ if and
only if $\ell\in\widehat{I}$ (or, equivalently, $\ell\in I$).
\end{itemize}

We recall that any partition $\lambda$ of $n$ gives rise to an $S_{n}%
$-representation called the Specht module $\mathcal{S}^{\lambda}$. If
$\lambda$ is a partition of $n$, and if $a\in\mathcal{A}$, then the action of
$a$ on $\mathcal{S}^{\lambda}$ (that is, the $\mathbf{k}$-linear map
$\mathcal{S}^{\lambda}\rightarrow\mathcal{S}^{\lambda},\ w\mapsto aw$) will be
denoted by $L_{\lambda}\left(  a\right)  $.

\begin{definition}
\label{def.clI}Let $\lambda$ be a partition of $n$. Let $I$ be a lacunar
subset of $\left[  n-1\right]  $. Write the set $I\cup\left\{  n+1\right\}  $
as $\left\{  i_{1}<i_{2}<\cdots<i_{m}\right\}  $, so that $i_{m}=n+1$.
Furthermore, set $i_{0}:=1$. Set $j_{k}:=i_{k}-i_{k-1}$ for each $k\in\left[
m\right]  $. (Note that Lemma \ref{lem.Fi/Fi-1.sum-js} shows that $j_{1}\geq0$
and $j_{2},j_{3},\ldots,j_{m}>1$, hence $j_{2},j_{3},\ldots,j_{m}\geq2$.)

The $m$-Littlewood--Richardson coefficient $c_{\left(  j_{1}\right)
,\ \left(  j_{2}-1,1\right)  ,\ \left(  j_{3}-1,1\right)  ,\ \ldots,\ \left(
j_{m}-1,1\right)  }^{\lambda}$ (as defined in (\ref{eq.LR.k}), where the
subscripts are the partition $\left(  j_{1}\right)  $ followed by the
partitions $\left(  j_{k}-1,1\right)  $ for all $k\in\left[  2,m\right]  $)
will then be denoted by $c_{I}^{\lambda}$.
\end{definition}

As we recall from Corollary \ref{cor.Fi/Fi-1.lr}, if $I=Q_{i}$ for some
$i\in\left[  f_{n+1}\right]  $, then this coefficient $c_{I}^{\lambda}$ is the
multiplicity of the Specht module $\mathcal{S}^{\lambda}$ in the left
$\mathcal{A}$-module $F_{i}/F_{i-1}$ when $\mathbf{k}$ is a field of
characteristic $0$. Indeed, we can rewrite Corollary \ref{cor.Fi/Fi-1.lr} as
follows using Definition \ref{def.clI}:

\begin{corollary}
\label{cor.Fi/Fi-1.lr-rewr}Assume that $\mathbf{k}$ is a field of
characteristic $0$.

Let $i\in\left[  f_{n+1}\right]  $. Let $\operatorname*{Par}\nolimits_{n}$ be
the set of all partitions of $n$. Then,%
\[
F_{i}/F_{i-1}\cong\bigoplus_{\lambda\in\operatorname*{Par}\nolimits_{n}%
}\left(  \mathcal{S}^{\lambda}\right)  ^{\oplus c_{Q_{i}}^{\lambda}}.
\]

\end{corollary}

We shall now state our first main theorem:


\begin{theorem}
\label{thm.eigs}Let $\mathbf{k}$ be any field. Let $\lambda$ be a partition of
$n$. Let $\omega_{1},\omega_{2},\ldots,\omega_{n}\in\mathbf{k}$. For each
subset $I$ of $\left[  n\right]  $, we set%
\[
\omega_{I}:=\omega_{1}m_{I,1}+\omega_{2}m_{I,2}+\cdots+\omega_{n}m_{I,n}%
\in\mathbf{k}.
\]
Then:

\begin{enumerate}
\item[\textbf{(a)}] The eigenvalues of the operator $L_{\lambda}\left(
\omega_{1}t_{1}+\omega_{2}t_{2}+\cdots+\omega_{n}t_{n}\right)  $ on the Specht
module $\mathcal{S}^{\lambda}$ are the elements%
\begin{equation}
\omega_{I}\text{ for all lacunar subsets }I\subseteq\left[  n-1\right]  \text{
satisfying }c_{I}^{\lambda}\neq0,\nonumber
\end{equation}
and their respective algebraic multiplicities are the $c_{I}^{\lambda}$ in the
generic case (i.e., if no two $I$'s produce the same $\omega_{I}$; otherwise
the multiplicities of colliding eigenvalues should be added together).

\item[\textbf{(b)}] If all these $\omega_{I}$ (for all lacunar subsets
$I\subseteq\left[  n-1\right]  $ satisfying $c_{I}^{\lambda}\neq0$) are
distinct, then $L_{\lambda}\left(  \omega_{1}t_{1}+\omega_{2}t_{2}%
+\cdots+\omega_{n}t_{n}\right)  $ is diagonalizable.

\item[\textbf{(c)}] We have%
\[
\prod_{\substack{I\subseteq\left[  n-1\right]  \text{ is lacunar;}%
\\c_{I}^{\lambda}\neq0}}\left(  L_{\lambda}\left(  \omega_{1}t_{1}+\omega
_{2}t_{2}+\cdots+\omega_{n}t_{n}\right)  -\omega_{I}\operatorname*{id}%
\nolimits_{\mathcal{S}^{\lambda}}\right)  =0.
\]

\end{enumerate}
\end{theorem}



To prove this theorem, we will need a further theorem, which \textquotedblleft
maps\textquotedblright\ the Fibonacci filtration from $\mathcal{A}$ to a given
Specht module $\mathcal{S}^{\lambda}$:

\begin{theorem}
\label{thm.specht.filtr}Let $\mathbf{k}$ be any field of characteristic $0$.
Let $\lambda$ be a partition of $n$. Then, there exists a filtration
\[
0=F_{f_{n+1}}^{\lambda}\subseteq F_{f_{n+1}-1}^{\lambda}\subseteq
F_{f_{n+1}-2}^{\lambda}\subseteq\cdots\subseteq F_{2}^{\lambda}\subseteq
F_{1}^{\lambda}\subseteq F_{0}^{\lambda}=\mathcal{S}^{\lambda}%
\]
(note the \textquotedblleft backward\textquotedblright\ indexing!) of the
Specht module $\mathcal{S}^{\lambda}$ by left $\mathcal{T}$-submodules with
the following four properties:

\begin{enumerate}
\item Each subquotient $F_{i-1}^{\lambda}/F_{i}^{\lambda}$ (for $i\in\left[
f_{n+1}\right]  $) has dimension $c_{Q_{i}}^{\lambda}$ as a $\mathbf{k}%
$-vector space (see Definition \ref{def.clI} for the meaning of $c_{I}%
^{\lambda}$).

\item In particular, an $i\in\left[  f_{n+1}\right]  $ satisfies
$F_{i-1}^{\lambda}=F_{i}^{\lambda}$ if and only if $c_{Q_{i}}^{\lambda}=0$.

\item On each subquotient $F_{i-1}^{\lambda}/F_{i}^{\lambda}$ (for
$i\in\left[  f_{n+1}\right]  $), each element $t_{\ell}\in\mathcal{T}$ (for
$\ell\in\left[  n\right]  $) acts as multiplication by the scalar
$m_{Q_{i},\ell}$.

\item More generally, on each subquotient $F_{i-1}^{\lambda}/F_{i}^{\lambda}$
(for $i\in\left[  f_{n+1}\right]  $), each element $P\left(  t_{1}%
,t_{2},\ldots,t_{n}\right)  \in\mathcal{T}$ (where $P$ is a polynomial in $n$
noncommuting indeterminates over $\mathbf{k}$) acts as multiplication by the
scalar $P\left(  m_{Q_{i},1},m_{Q_{i},2},\ldots,m_{Q_{i},n}\right)  $.
\end{enumerate}
\end{theorem}

\subsection{Lemmas about $S_{n}$-representations}

In order to prove Theorem \ref{thm.eigs} and Theorem \ref{thm.specht.filtr},
we need a few lemmas about left $\mathcal{A}$-modules (i.e., representations
of $S_{n}$). We begin with something basic and well-known (see, e.g.,
\cite[last paragraph of \S 5.13]{Etingof-et-al}):

\begin{proposition}
\label{prop.specht.tensA}Assume that $\mathbf{k}$ is a field of characteristic
$0$. Let $\lambda$ and $\mu$ be two partitions of $n$. Then, the Specht
modules $\mathcal{S}^{\lambda}$ and $\mathcal{S}^{\mu}$ satisfy%
\[
\operatorname*{Hom}\nolimits_{\mathcal{A}}\left(  \mathcal{S}^{\lambda
},\mathcal{S}^{\mu}\right)  \cong%
\begin{cases}
\mathbf{k}, & \text{if }\lambda=\mu;\\
0, & \text{if }\lambda\neq\mu
\end{cases}
\ \ \ \ \ \ \ \ \ \ \text{as }\mathbf{k}\text{-vector spaces.}%
\]
(Here and in the following, \textquotedblleft$\operatorname*{Hom}%
\nolimits_{\mathcal{A}}$\textquotedblright\ always stands for the set of left
$\mathcal{A}$-module morphisms. This is always a $\mathbf{k}$-vector space,
but usually not an $\mathcal{A}$-module on any side.)
\end{proposition}

For the sake of completeness, we give a proof of this proposition in the
Appendix (Subsection \ref{subsec.apx.proofs.specht.tensA}).

\begin{lemma}
\label{lem.specht.exact}Assume that $\mathbf{k}$ is a field of characteristic
$0$. Let $\lambda$ be a partition of $n$. Define the contravariant functor
$\operatorname*{Hom}\nolimits_{\mathcal{A}}\left(  -,\mathcal{S}^{\lambda
}\right)  $ from the category of left $\mathcal{A}$-modules to the category of
$\mathbf{k}$-vector spaces that is given by%
\[
X\mapsto\operatorname*{Hom}\nolimits_{\mathcal{A}}\left(  X,\mathcal{S}%
^{\lambda}\right)  \ \ \ \ \ \ \ \ \ \ \text{on objects}%
\]
and likewise on morphisms. (This is a contravariant Hom functor.)

This contravariant functor $\operatorname*{Hom}\nolimits_{\mathcal{A}}\left(
-,\mathcal{S}^{\lambda}\right)  $ is exact (i.e., respects exact sequences).
\end{lemma}

\begin{proof}
The $\mathbf{k}$-algebra $\mathcal{A}=\mathbf{k}\left[  S_{n}\right]  $ is
semisimple (by Maschke's theorem, since $\mathbf{k}$ is a field of
characteristic $0$). Hence, every short exact sequence of left $\mathcal{A}%
$-modules is split. Consequently, any Hom functor from the category of left
$\mathcal{A}$-modules is exact\footnote{We remind the reader of the proof of
this fact: Let $\mathcal{F}$ be a contravariant Hom functor from the category
of left $\mathcal{A}$-modules (the case of a covariant Hom functor is
analogous). We must show that $\mathcal{F}$ is exact.
\par
Let $0\rightarrow X\rightarrow Y\rightarrow Z\rightarrow0$ be a short exact
sequence of left $\mathcal{A}$-modules. Then, this sequence is split (since
every short exact sequence of left $\mathcal{A}$-modules is split), and thus
is isomorphic to the obvious exact sequence $0\rightarrow X\rightarrow X\oplus
Z\rightarrow Z\rightarrow0$. The Hom functor $\mathcal{F}$ sends the latter
sequence to another short exact sequence $0\rightarrow\mathcal{F}\left(
X\right)  \rightarrow\mathcal{F}\left(  X\right)  \oplus\mathcal{F}\left(
Z\right)  \rightarrow\mathcal{F}\left(  Z\right)  \rightarrow0$ (since Hom
functors respect finite direct sums). But $\mathcal{F}$ is a functor; thus,
$\mathcal{F}$ sends any two isomorphic complexes to two isomorphic complexes.
Hence, because the two short exact sequences $0\rightarrow X\rightarrow
Y\rightarrow Z\rightarrow0$ and $0\rightarrow X\rightarrow X\oplus
Z\rightarrow Z\rightarrow0$ are isomorphic, their images under $\mathcal{F}$
must also be isomorphic. In other words, the image of the short exact sequence
$0\rightarrow X\rightarrow Y\rightarrow Z\rightarrow0$ under $\mathcal{F}$ is
isomorphic to the image of $0\rightarrow X\rightarrow X\oplus Z\rightarrow
Z\rightarrow0$ under $\mathcal{F}$. But the latter image is exact, as we have
shown. Hence, the former image is also exact.
\par
Thus, $\mathcal{F}$ sends the original short exact sequence $0\rightarrow
X\rightarrow Y\rightarrow Z\rightarrow0$ to a short exact sequence. Since this
is true for any short exact sequence $0\rightarrow X\rightarrow Y\rightarrow
Z\rightarrow0$, we thus have shown that the functor $\mathcal{F}$ respects
short exact sequences. In other words, $\mathcal{F}$ is exact, qed.}. Thus,
the contravariant Hom functor $\operatorname*{Hom}\nolimits_{\mathcal{A}%
}\left(  -,\mathcal{S}^{\lambda}\right)  $ is exact. This proves Lemma
\ref{lem.specht.exact}.
\end{proof}

Note that the Specht module $\mathcal{S}^{\lambda}$ in Lemma
\ref{lem.specht.exact} could be replaced by any left $\mathcal{A}$-module, but
we will use $\mathcal{S}^{\lambda}$ only. The same applies to the following lemma:

\begin{lemma}
\label{lem.specht.f1A}Let $\lambda$ be a partition of $n$. Let $J$ be a left
$\mathcal{A}$-submodule of $\mathcal{A}$ (that is, a left ideal of
$\mathcal{A}$). Then, there is a canonical $\mathbf{k}$-vector space
isomorphism%
\begin{align*}
\operatorname*{Hom}\nolimits_{\mathcal{A}}\left(  \mathcal{A}/J,\ \mathcal{S}%
^{\lambda}\right)   &  \rightarrow\left\{  v\in\mathcal{S}^{\lambda}%
\ \mid\ Jv=0\right\}  ,\\
f  &  \mapsto f\left(  \overline{1_{\mathcal{A}}}\right)  .
\end{align*}

\end{lemma}

\begin{proof}
The left $\mathcal{A}$-module morphisms from $\mathcal{A}/J$ to $\mathcal{S}%
^{\lambda}$ can be identified with the left $\mathcal{A}$-module morphisms
from $\mathcal{A}$ to $\mathcal{S}^{\lambda}$ that vanish on $J$. Thus, we
obtain a $\mathbf{k}$-vector space isomorphism%
\begin{align*}
\Phi:\operatorname*{Hom}\nolimits_{\mathcal{A}}\left(  \mathcal{A}%
/J,\ \mathcal{S}^{\lambda}\right)   &  \rightarrow\left\{  g\in
\operatorname*{Hom}\nolimits_{\mathcal{A}}\left(  \mathcal{A},\ \mathcal{S}%
^{\lambda}\right)  \ \mid\ g\left(  J\right)  =0\right\}  ,\\
f  &  \mapsto\left(  \mathcal{A}\rightarrow\mathcal{S}^{\lambda},\ a\mapsto
f\left(  \overline{a}\right)  \right)  .
\end{align*}


However, recall the well-known $\mathbf{k}$-vector space isomorphism
$\operatorname*{Hom}\nolimits_{A}\left(  A,M\right)  \cong M$ that holds for
any $\mathbf{k}$-algebra $A$ and any left $A$-module $M$. Thus, in particular,
the left $\mathcal{A}$-module morphisms from $\mathcal{A}$ to $\mathcal{S}%
^{\lambda}$ can be identified with the elements of $\mathcal{S}^{\lambda}$ via
the $\mathbf{k}$-vector space isomorphism
\begin{align*}
\Psi:\operatorname*{Hom}\nolimits_{\mathcal{A}}\left(  \mathcal{A}%
,\ \mathcal{S}^{\lambda}\right)   &  \rightarrow\mathcal{S}^{\lambda},\\
g  &  \mapsto g\left(  1_{\mathcal{A}}\right)  .
\end{align*}
This latter isomorphism $\Psi$ has the property that an arbitrary
left $\mathcal{A}$-module
morphism $g\in\operatorname*{Hom}\nolimits_{\mathcal{A}}\left(  \mathcal{A}%
,\ \mathcal{S}^{\lambda}\right)  $ satisfies $g\left(  J\right)  =0$ if and
only if its image $\Psi\left(  g\right)  $ satisfies $J\cdot\Psi\left(
g\right)  =0$ (since $g\left(  J\right)  =g\left(  J\cdot1_{\mathcal{A}%
}\right)  =J\cdot\underbrace{g\left(  1_{\mathcal{A}}\right)  }_{=\Psi\left(
g\right)  }=J\cdot\Psi\left(  g\right)  $). Thus, $\Psi$ can be restricted to
a $\mathbf{k}$-vector space isomorphism
\begin{align*}
\Psi^{\prime}:\left\{  g\in\operatorname*{Hom}\nolimits_{\mathcal{A}}\left(
\mathcal{A},\ \mathcal{S}^{\lambda}\right)  \ \mid\ g\left(  J\right)
=0\right\}   &  \rightarrow\left\{  v\in\mathcal{S}^{\lambda}\ \mid
\ Jv=0\right\}  ,\\
g  &  \mapsto g\left(  1_{\mathcal{A}}\right)  .
\end{align*}
The composition $\Psi^{\prime}\circ\Phi$ is thus a $\mathbf{k}$-vector space
isomorphism%
\begin{align*}
\operatorname*{Hom}\nolimits_{\mathcal{A}}\left(  \mathcal{A}/J,\ \mathcal{S}%
^{\lambda}\right)   &  \rightarrow\left\{  v\in\mathcal{S}^{\lambda}%
\ \mid\ Jv=0\right\}  ,\\
f  &  \mapsto f\left(  \overline{1_{\mathcal{A}}}\right)  .
\end{align*}
This is clearly canonical in $J$, so that Lemma \ref{lem.specht.f1A} is proved.
\end{proof}

\subsection{The proofs}

We are now ready to prove Theorem \ref{thm.specht.filtr} and Theorem
\ref{thm.eigs}, in this order.

\begin{proof}
[Proof of Theorem \ref{thm.specht.filtr}.]For each $i\in\left[  0,f_{n+1}%
\right]  $, we define a subset $F_{i}^{\lambda}$ of $\mathcal{S}^{\lambda}$ by%
\[
F_{i}^{\lambda}:=\left\{  v\in\mathcal{S}^{\lambda}\ \mid\ F_{i}v=0\right\}
.
\]
This subset $F_{i}^{\lambda}$ is actually a left $\mathcal{T}$-submodule of
$\mathcal{S}^{\lambda}$ (since Proposition \ref{prop.fibfilt.A} shows that
$F_{i}$ is a right $\mathcal{T}$-submodule of $\mathcal{A}$%
)\ \ \ \ \footnote{\textit{Proof.} To show that $F_{i}^{\lambda}$ is closed
under addition, we observe that any $v,w\in F_{i}^{\lambda}$ satisfy $v+w\in
F_{i}^{\lambda}$ (because $v,w\in F_{i}^{\lambda}$ entails $F_{i}v=0$ and
$F_{i}w=0$ and thus $F_{i}\left(  v+w\right)  \subseteq\underbrace{F_{i}%
v}_{=0}+\underbrace{F_{i}w}_{=0}=0$, so that $F_{i}\left(  v+w\right)  =0$ and
therefore $v+w\in F_{i}^{\lambda}$).
\par
To show that $F_{i}^{\lambda}$ is closed under the left $\mathcal{T}$-action,
we observe that any $t\in\mathcal{T}$ and $v\in F_{i}^{\lambda}$ satisfy
$tv\in F_{i}^{\lambda}$ (because $v\in F_{i}^{\lambda}$ entails $F_{i}v=0$;
but we have $F_{i}t\subseteq F_{i}$ since $F_{i}$ is a right $\mathcal{T}%
$-submodule of $\mathcal{A}$; thus $F_{i}\left(  tv\right)  =\underbrace{F_{i}%
t}_{\subseteq F_{i}}v\subseteq F_{i}v=0$ and therefore $F_{i}\left(
tv\right)  =0$, so that $tv\in F_{i}^{\lambda}$).
\par
An even simpler argument shows that $0\in F_{i}^{\lambda}$.}. Moreover,
because of%
\[
0=F_{0}\subseteq F_{1}\subseteq F_{2}\subseteq\cdots\subseteq F_{f_{n+1}%
}=\mathbf{k}\left[  S_{n}\right]  =\mathcal{A},
\]
we have%
\[
\mathcal{S}^{\lambda}=F_{0}^{\lambda}\supseteq F_{1}^{\lambda}\supseteq
F_{2}^{\lambda}\supseteq\cdots\supseteq F_{f_{n+1}}^{\lambda}=0.
\]
This is a left $\mathcal{T}$-module filtration of $\mathcal{S}^{\lambda}$,
albeit written backwards. We can rewrite it as%
\[
0=F_{f_{n+1}}^{\lambda}\subseteq F_{f_{n+1}-1}^{\lambda}\subseteq
F_{f_{n+1}-2}^{\lambda}\subseteq\cdots\subseteq F_{2}^{\lambda}\subseteq
F_{1}^{\lambda}\subseteq F_{0}^{\lambda}=\mathcal{S}^{\lambda}.
\]
Our goal is to show that this filtration satisfies the four properties 1, 2, 3
and 4 claimed in Theorem \ref{thm.specht.filtr}.

For this purpose, we fix $i\in\left[  f_{n+1}\right]  $. First, we shall show
property 3. We must show that each element $t_{\ell}\in\mathcal{T}$ acts on
$F_{i-1}^{\lambda}/F_{i}^{\lambda}$ as multiplication by the scalar
$m_{Q_{i},\ell}$. So we let $\ell\in\left[  n\right]  $ and $\overline{v}\in
F_{i-1}^{\lambda}/F_{i}^{\lambda}$ (with $v\in F_{i-1}^{\lambda}$) be
arbitrary. We must show that $t_{\ell}\overline{v}=m_{Q_{i},\ell}\overline{v}$
in $F_{i-1}^{\lambda}/F_{i}^{\lambda}$.

We have $v\in F_{i-1}^{\lambda}$. In other words, $v\in\mathcal{S}^{\lambda}$
and $F_{i-1}v=0$ (by the definition of $F_{i-1}^{\lambda}$).

Theorem \ref{thm.t-simultri} \textbf{(c)} yields $F_{i}\cdot\left(  t_{\ell
}-m_{Q_{i},\ell}\right)  \subseteq F_{i-1}$. Hence,%
\[
\underbrace{F_{i}\cdot\left(  t_{\ell}-m_{Q_{i},\ell}\right)  }_{\subseteq
F_{i-1}}v\subseteq F_{i-1}v=0,
\]
so that $F_{i}\cdot\left(  t_{\ell}-m_{Q_{i},\ell}\right)  v=0$. In other
words, $\left(  t_{\ell}-m_{Q_{i},\ell}\right)  v\in F_{i}^{\lambda}$ (by the
definition of $F_{i}^{\lambda}$). In other words, $t_{\ell}v-m_{Q_{i},\ell
}v\in F_{i}^{\lambda}$. In other words, $\overline{t_{\ell}v}=\overline
{m_{Q_{i},\ell}v}$ in $F_{i-1}^{\lambda}/F_{i}^{\lambda}$. In other words,
$t_{\ell}\overline{v}=m_{Q_{i},\ell}\overline{v}$ in $F_{i-1}^{\lambda}%
/F_{i}^{\lambda}$. Thus, the proof of property 3 is complete.

Property 4 follows immediately from property 3.

Let us now prove property 1.

Note that $F_{i}$ is a left $\mathcal{A}$-submodule of $\mathcal{A}$ (by
Proposition \ref{prop.fibfilt.A}). Lemma \ref{lem.specht.f1A} shows that
whenever $J$ is a left $\mathcal{A}$-submodule of $\mathcal{A}$ (that is, a
left ideal of $\mathcal{A}$), there is a canonical $\mathbf{k}$-vector space
isomorphism%
\begin{align*}
\operatorname*{Hom}\nolimits_{\mathcal{A}}\left(  \mathcal{A}/J,\ \mathcal{S}%
^{\lambda}\right)   &  \rightarrow\left\{  v\in\mathcal{S}^{\lambda}%
\ \mid\ Jv=0\right\}  ,\\
f  &  \mapsto f\left(  \overline{1_{\mathcal{A}}}\right)  .
\end{align*}
Thus, we have%
\begin{equation}
\left\{  v\in\mathcal{S}^{\lambda}\ \mid\ Jv=0\right\}  \cong%
\operatorname*{Hom}\nolimits_{\mathcal{A}}\left(  \mathcal{A}/J,\ \mathcal{S}%
^{\lambda}\right)  \label{pf.thm.specht.filtr.caniso}%
\end{equation}
canonically for each left $\mathcal{A}$-submodule $J$ of $\mathcal{A}$. Now,
the definition of $F_{i}^{\lambda}$ yields%
\begin{equation}
F_{i}^{\lambda}=\left\{  v\in\mathcal{S}^{\lambda}\ \mid\ F_{i}v=0\right\}
\cong\operatorname*{Hom}\nolimits_{\mathcal{A}}\left(  \mathcal{A}%
/F_{i},\ \mathcal{S}^{\lambda}\right)  \label{pf.thm.specht.filtr.canisoi}%
\end{equation}
canonically (by (\ref{pf.thm.specht.filtr.caniso})) and similarly%
\begin{equation}
F_{i-1}^{\lambda}\cong\operatorname*{Hom}\nolimits_{\mathcal{A}}\left(
\mathcal{A}/F_{i-1},\ \mathcal{S}^{\lambda}\right)  .
\label{pf.thm.specht.filtr.canisoi-1}%
\end{equation}


However, Lemma \ref{lem.specht.exact} shows that the contravariant functor
$\operatorname*{Hom}\nolimits_{\mathcal{A}}\left(  -,\mathcal{S}^{\lambda
}\right)  $ from the category of left $\mathcal{A}$-modules to the category of
$\mathbf{k}$-vector spaces is exact. Hence, applying this contravariant
functor to the exact sequence%
\[
0\rightarrow F_{i}/F_{i-1}\rightarrow\mathcal{A}/F_{i-1}\rightarrow
\mathcal{A}/F_{i}\rightarrow0
\]
of left $\mathcal{A}$-modules, we obtain an exact sequence%
\[
0\rightarrow\operatorname*{Hom}\nolimits_{\mathcal{A}}\left(  \mathcal{A}%
/F_{i},\ \mathcal{S}^{\lambda}\right)  \rightarrow\operatorname*{Hom}%
\nolimits_{\mathcal{A}}\left(  \mathcal{A}/F_{i-1},\ \mathcal{S}^{\lambda
}\right)  \rightarrow\operatorname*{Hom}\nolimits_{\mathcal{A}}\left(
F_{i}/F_{i-1},\ \mathcal{S}^{\lambda}\right)  \rightarrow0
\]
of $\mathbf{k}$-vector spaces. In view of (\ref{pf.thm.specht.filtr.canisoi-1}%
) and (\ref{pf.thm.specht.filtr.canisoi}), we can rewrite this latter exact
sequence as%
\[
0\rightarrow F_{i}^{\lambda}\rightarrow F_{i-1}^{\lambda}\rightarrow
\operatorname*{Hom}\nolimits_{\mathcal{A}}\left(  F_{i}/F_{i-1},\ \mathcal{S}%
^{\lambda}\right)  \rightarrow0.
\]
The arrow $F_{i}^{\lambda}\rightarrow F_{i-1}^{\lambda}$ here is the canonical
inclusion (since the isomorphisms in (\ref{pf.thm.specht.filtr.canisoi-1}) and
(\ref{pf.thm.specht.filtr.canisoi}) are the canonical ones), and thus we
obtain%
\begin{equation}
\operatorname*{Hom}\nolimits_{\mathcal{A}}\left(  F_{i}/F_{i-1},\ \mathcal{S}%
^{\lambda}\right)  \cong F_{i-1}^{\lambda}/F_{i}^{\lambda}
\label{pf.thm.specht.filtr.iso5}%
\end{equation}
from the exactness of our sequence.

However, Corollary \ref{cor.Fi/Fi-1.lr-rewr} says that%
\begin{equation}
F_{i}/F_{i-1}\cong\bigoplus_{\nu\in\operatorname*{Par}\nolimits_{n}}\left(
\mathcal{S}^{\nu}\right)  ^{\oplus c_{Q_{i}}^{\nu}}.
\label{pf.thm.specht.filtr.iso6}%
\end{equation}
Hence,%
\begin{align}
\operatorname*{Hom}\nolimits_{\mathcal{A}}\left(  F_{i}/F_{i-1},\ \mathcal{S}%
^{\lambda}\right)   &  \cong\operatorname*{Hom}\nolimits_{\mathcal{A}}\left(
\bigoplus_{\nu\in\operatorname*{Par}\nolimits_{n}}\left(  \mathcal{S}^{\nu
}\right)  ^{\oplus c_{Q_{i}}^{\nu}},\ \mathcal{S}^{\lambda}\right) \nonumber\\
&  \cong\bigoplus_{\nu\in\operatorname*{Par}\nolimits_{n}}\left(
\operatorname*{Hom}\nolimits_{\mathcal{A}}\left(  \mathcal{S}^{\nu
},\mathcal{S}^{\lambda}\right)  \right)  ^{\oplus c_{Q_{i}}^{\nu}}
\label{pf.thm.specht.filtr.iso6Hom}%
\end{align}
(since Hom functors respect finite direct sums).

But each $\nu\in\operatorname*{Par}\nolimits_{n}$ satisfies%
\[
\operatorname*{Hom}\nolimits_{\mathcal{A}}\left(  \mathcal{S}^{\nu
},\mathcal{S}^{\lambda}\right)  \cong%
\begin{cases}
\mathbf{k}, & \text{if }\nu=\lambda;\\
0, & \text{if }\nu\neq\lambda
\end{cases}
\ \ \ \ \ \ \ \ \ \ \left(  \text{by Proposition \ref{prop.specht.tensA}%
}\right)
\]
and thus
\begin{equation}
\left(  \operatorname*{Hom}\nolimits_{\mathcal{A}}\left(  \mathcal{S}^{\nu
},\mathcal{S}^{\lambda}\right)  \right)  ^{\oplus c_{Q_{i}}^{\nu}}\cong\left(
%
\begin{cases}
\mathbf{k}, & \text{if }\nu=\lambda;\\
0, & \text{if }\nu\neq\lambda
\end{cases}
\right)  ^{\oplus c_{Q_{i}}^{\nu}}\cong%
\begin{cases}
\mathbf{k}^{\oplus c_{Q_{i}}^{\nu}}, & \text{if }\nu=\lambda;\\
0, & \text{if }\nu\neq\lambda.
\end{cases}
\nonumber
\end{equation}
Thus, we can rewrite (\ref{pf.thm.specht.filtr.iso6Hom}) as%
\[
\operatorname*{Hom}\nolimits_{\mathcal{A}}\left(  F_{i}/F_{i-1},\ \mathcal{S}%
^{\lambda}\right)  \cong\bigoplus_{\nu\in\operatorname*{Par}\nolimits_{n}}%
\begin{cases}
\mathbf{k}^{\oplus c_{Q_{i}}^{\nu}}, & \text{if }\nu=\lambda;\\
0, & \text{if }\nu\neq\lambda
\end{cases}
\ \ \ =\mathbf{k}^{\oplus c_{Q_{i}}^{\lambda}}.
\]
Comparing this with (\ref{pf.thm.specht.filtr.iso5}), we see that%
\[
F_{i-1}^{\lambda}/F_{i}^{\lambda}\cong\mathbf{k}^{\oplus c_{Q_{i}}^{\lambda}%
}.
\]


Thus, the $\mathbf{k}$-vector space $F_{i-1}^{\lambda}/F_{i}^{\lambda}$ has
dimension $c_{Q_{i}}^{\lambda}$. This proves property 1.

Property 2 follows immediately from property 1 (since $F_{i-1}^{\lambda}%
=F_{i}^{\lambda}$ is equivalent to $\dim\left(  F_{i-1}^{\lambda}%
/F_{i}^{\lambda}\right)  =0$). Hence, our proof of Theorem
\ref{thm.specht.filtr} is complete.
\end{proof}

\begin{proof}
[Proof of Theorem \ref{thm.eigs}.]In the following, the symbol $\dim$ will
always refer to the dimension of a $\mathbf{k}$-vector space, even if some
other module structures are present. Thus, in particular, if $X$ is a
$\mathcal{T}$-module, then $\dim X$ will mean the dimension of $X$ as a
$\mathbf{k}$-vector space. \medskip

\textbf{(a)} Let us first assume that $\mathbf{k}$ is a field of
characteristic $0$. We shall later extend this to the general case.

Theorem \ref{thm.specht.filtr} shows that there exists a filtration
\begin{equation}
0=F_{f_{n+1}}^{\lambda}\subseteq F_{f_{n+1}-1}^{\lambda}\subseteq
F_{f_{n+1}-2}^{\lambda}\subseteq\cdots\subseteq F_{2}^{\lambda}\subseteq
F_{1}^{\lambda}\subseteq F_{0}^{\lambda}=\mathcal{S}^{\lambda}
\label{pf.thm.eigs.filt}%
\end{equation}
of the Specht module $\mathcal{S}^{\lambda}$ by left $\mathcal{T}$-submodules
with the four properties 1, 2, 3 and 4 stated in Theorem
\ref{thm.specht.filtr}. Consider this filtration. Fix any basis $\left(
v_{1},v_{2},\ldots,v_{s}\right)  $ of $\mathcal{S}^{\lambda}$ that conforms
with this filtration (i.e., a basis that begins with a basis of $F_{f_{n+1}%
-1}^{\lambda}$, then extends it to a basis of $F_{f_{n+1}-2}^{\lambda}$, then
extends it to a basis of $F_{f_{n+1}-3}^{\lambda}$, and so on), so that each
$F_{i}^{\lambda}$ is spanned by $v_{1},v_{2},\ldots,v_{j\left(  i\right)  }$
for some $j\left(  i\right)  \in\left[  0,s\right]  $. Note that the
inclusions in (\ref{pf.thm.eigs.filt}) yield%
\[
0=j\left(  f_{n+1}\right)  \leq j\left(  f_{n+1}-1\right)  \leq j\left(
f_{n+1}-2\right)  \leq\cdots\leq j\left(  2\right)  \leq j\left(  1\right)
\leq j\left(  0\right)  =s.
\]
Note that each $i\in\left[  f_{n+1}\right]  $ satisfies $j\left(  i\right)
=\dim\left(  F_{i}^{\lambda}\right)  $ and $j\left(  i-1\right)  =\dim\left(
F_{i-1}^{\lambda}\right)  $ and thus
\begin{align}
j\left(  i-1\right)  -j\left(  i\right)   &  =\dim\left(  F_{i-1}^{\lambda
}\right)  -\dim\left(  F_{i}^{\lambda}\right)  =\dim\left(  F_{i-1}^{\lambda
}/F_{i}^{\lambda}\right) \nonumber\\
&  =c_{Q_{i}}^{\lambda} \label{pf.thm.eigs.j-j}%
\end{align}
(by property 1 of our filtration).

The operator $L_{\lambda}\left(  \omega_{1}t_{1}+\omega_{2}t_{2}+\cdots
+\omega_{n}t_{n}\right)  $ preserves the filtration (\ref{pf.thm.eigs.filt})
(since this filtration is a filtration by left $\mathcal{T}$-submodules, but
$\omega_{1}t_{1}+\omega_{2}t_{2}+\cdots+\omega_{n}t_{n}\in\mathcal{T}$).
Hence, the matrix $M$ that represents this operator with respect to the basis
$\left(  v_{1},v_{2},\ldots,v_{s}\right)  $ is block-upper-triangular with
blocks of sizes
\[
j\left(  i-1\right)  -j\left(  i\right)  \ \ \ \ \ \ \ \ \ \ \text{for all
}i\in\left[  f_{n+1}\right]
\]
(because, e.g., the fact that the operator preserves $F_{i}^{\lambda}$ means
that the first $j\left(  i\right)  $ columns of the matrix $M$ have zeroes
everywhere below the $j\left(  i\right)  $-th row). Moreover, property 4 of
our filtration shows that the element $\omega_{1}t_{1}+\omega_{2}t_{2}%
+\cdots+\omega_{n}t_{n}\in\mathcal{T}$ acts as multiplication by the scalar
\[
\omega_{1}m_{Q_{i},1}+\omega_{2}m_{Q_{i},2}+\cdots+\omega_{n}m_{Q_{i}%
,n}=\omega_{Q_{i}}\ \ \ \ \ \ \ \ \ \ \left(  \text{by the definition of
}\omega_{Q_{i}}\right)
\]
on each subquotient $F_{i-1}^{\lambda}/F_{i}^{\lambda}$. In other words, for
each $v\in F_{i-1}^{\lambda}$, the vector $\overline{v}\in F_{i-1}^{\lambda
}/F_{i}^{\lambda}$ satisfies%
\begin{equation}
\left(  \omega_{1}t_{1}+\omega_{2}t_{2}+\cdots+\omega_{n}t_{n}\right)
\cdot\overline{v}=\omega_{Q_{i}}\overline{v}, \label{pf.thm.eigs.3}%
\end{equation}
and therefore%
\begin{align*}
&  \left(  L_{\lambda}\left(  \omega_{1}t_{1}+\omega_{2}t_{2}+\cdots
+\omega_{n}t_{n}\right)  \right)  \left(  v\right) \\
&  =\left(  \omega_{1}t_{1}+\omega_{2}t_{2}+\cdots+\omega_{n}t_{n}\right)
\cdot v\\
&  =\omega_{Q_{i}}v+\left(  \text{some element of }F_{i}^{\lambda}\right)
\ \ \ \ \ \ \ \ \ \ \left(  \text{by (\ref{pf.thm.eigs.3})}\right) \\
&  =\omega_{Q_{i}}v+\left(  \text{some linear combination of }v_{1}%
,v_{2},\ldots,v_{j\left(  i\right)  }\right)
\end{align*}
(since $F_{i}^{\lambda}$ is spanned by $v_{1},v_{2},\ldots,v_{j\left(
i\right)  }$). We can apply this in particular to $v=v_{k}$ for each
$k\in\left[  j\left(  i-1\right)  \right]  $ (since $F_{i-1}^{\lambda}$ is
spanned by $v_{1},v_{2},\ldots,v_{j\left(  i-1\right)  }$), and conclude that%
\begin{align*}
&  \left(  L_{\lambda}\left(  \omega_{1}t_{1}+\omega_{2}t_{2}+\cdots
+\omega_{n}t_{n}\right)  \right)  \left(  v_{k}\right) \\
&  =\omega_{Q_{i}}v_{k}+\left(  \text{some linear combination of }v_{1}%
,v_{2},\ldots,v_{j\left(  i\right)  }\right)
\end{align*}
for each $k\in\left[  j\left(  i-1\right)  \right]  $.

Thus, the matrix $M$ that represents the operator $L_{\lambda}\left(
\omega_{1}t_{1}+\omega_{2}t_{2}+\cdots+\omega_{n}t_{n}\right)  $ with respect
to the basis $\left(  v_{1},v_{2},\ldots,v_{s}\right)  $ is not only
block-upper-triangular, but also has the property that its $i$-th diagonal
block (for each $i\in\left[  f_{n+1}\right]  $, counted from the end) is the
scalar matrix $\omega_{Q_{i}}\cdot I_{j\left(  i-1\right)  -j\left(  i\right)
}=\omega_{Q_{i}}\cdot I_{c_{Q_{i}}^{\lambda}}$ (by (\ref{pf.thm.eigs.j-j})).
Consequently, the matrix $M$ is upper-triangular, and its diagonal entries are
the elements $\omega_{Q_{i}}$ for all $i\in\left[  f_{n+1}\right]  $, with
each $\omega_{Q_{i}}$ appearing $c_{Q_{i}}^{\lambda}$ times (this means that
if $c_{Q_{i}}^{\lambda}=0$, then $\omega_{Q_{i}}$ does not appear at all).

Of course, this allows us to read off the eigenvalues of this matrix $M$, and
thus of the operator $L_{\lambda}\left(  \omega_{1}t_{1}+\omega_{2}%
t_{2}+\cdots+\omega_{n}t_{n}\right)  $ (since the eigenvalues of a triangular
matrix are just its diagonal entries). We conclude that the eigenvalues of
$L_{\lambda}\left(  \omega_{1}t_{1}+\omega_{2}t_{2}+\cdots+\omega_{n}%
t_{n}\right)  $ are the elements $\omega_{Q_{i}}$ for all $i\in\left[
f_{n+1}\right]  $, with each $\omega_{Q_{i}}$ appearing with algebraic
multiplicity $c_{Q_{i}}^{\lambda}$. Since $Q_{1},Q_{2},\ldots,Q_{f_{n+1}}$ are
just the lacunar subsets of $\left[  n-1\right]  $, we can rewrite this as
follows: The eigenvalues of $L_{\lambda}\left(  \omega_{1}t_{1}+\omega
_{2}t_{2}+\cdots+\omega_{n}t_{n}\right)  $ are the elements $\omega_{I}$ for
all lacunar subsets $I\subseteq\left[  n-1\right]  $, with each $\omega_{I}$
appearing with algebraic multiplicity $c_{I}^{\lambda}$. We can restrict this
list to those lacunar subsets $I\subseteq\left[  n-1\right]  $ that satisfy
$c_{I}^{\lambda}\neq0$ (since an eigenvalue $\omega_{I}$ that appears with
algebraic multiplicity $c_{I}^{\lambda}=0$ simply does not appear at all).

This proves Theorem \ref{thm.eigs} \textbf{(a)} in the case when $\mathbf{k}$
is a field of characteristic $0$. It remains to extend the proof to the case
when $\mathbf{k}$ is an arbitrary field. But there is a standard trick for
this: We recast our result as a polynomial identity. Namely, Theorem
\ref{thm.eigs} \textbf{(a)} is saying that
\[
\underbrace{\det\left(  x\operatorname*{id}\nolimits_{\mathcal{S}^{\lambda}%
}-\,L_{\lambda}\left(  \omega_{1}t_{1}+\omega_{2}t_{2}+\cdots+\omega_{n}%
t_{n}\right)  \right)  }_{\substack{\text{This is the characteristic
polynomial of the}\\\text{endomorphism }L_{\lambda}\left(  \omega_{1}%
t_{1}+\omega_{2}t_{2}+\cdots+\omega_{n}t_{n}\right)  \text{ of }%
\mathcal{S}^{\lambda}}}=\prod_{I\subseteq\left[  n-1\right]  \text{ lacunar}%
}\left(  x-\omega_{I}\right)  ^{c_{I}^{\lambda}}%
\]
in the polynomial ring $\mathbf{k}\left[  x\right]  $. This is a polynomial
identity in the indeterminates $x,\omega_{1},\omega_{2},\ldots,\omega_{n}$
(since the Specht module $\mathcal{S}^{\lambda}$ has a basis consisting of the
standard polytabloids, and the action of $S_{n}$ on this basis is independent
of the base field $\mathbf{k}$). Thus, knowing that this identity holds
whenever $\mathbf{k}$ is a field of characteristic $0$, we can immediately
conclude that it holds for all fields $\mathbf{k}$ (and even all commutative
rings $\mathbf{k}$). This proves Theorem \ref{thm.eigs} \textbf{(a)} in the
general case. \medskip

\textbf{(c)} Again, let us first assume that $\mathbf{k}$ is a field of
characteristic $0$. Recall the filtration (\ref{pf.thm.eigs.filt}) constructed
in the proof of part \textbf{(a)}. As we saw in that proof, the element
$\omega_{1}t_{1}+\omega_{2}t_{2}+\cdots+\omega_{n}t_{n}\in\mathcal{T}$ acts as
multiplication by the scalar $\omega_{Q_{i}}$ on each subquotient
$F_{i-1}^{\lambda}/F_{i}^{\lambda}$ of that filtration. In other words, for
each $i\in\left[  f_{n+1}\right]  $, we have%
\[
\left(  \omega_{1}t_{1}+\omega_{2}t_{2}+\cdots+\omega_{n}t_{n}\right)
\overline{v}=\omega_{Q_{i}}\overline{v}\ \ \ \ \ \ \ \ \ \ \text{for each
}\overline{v}\in F_{i-1}^{\lambda}/F_{i}^{\lambda},
\]
that is,%
\[
\left(  \omega_{1}t_{1}+\omega_{2}t_{2}+\cdots+\omega_{n}t_{n}\right)
v-\omega_{Q_{i}}v\in F_{i}^{\lambda}\ \ \ \ \ \ \ \ \ \ \text{for each }v\in
F_{i-1}^{\lambda}.
\]
In other words, for each $i\in\left[  f_{n+1}\right]  $, we have%
\[
\left(  L_{\lambda}\left(  \omega_{1}t_{1}+\omega_{2}t_{2}+\cdots+\omega
_{n}t_{n}\right)  -\omega_{Q_{i}}\operatorname*{id}\nolimits_{\mathcal{S}%
^{\lambda}}\right)  F_{i-1}^{\lambda}\subseteq F_{i}^{\lambda}.
\]
Hence, the operator%
\[
\prod_{i\in\left[  f_{n+1}\right]  }\left(  L_{\lambda}\left(  \omega_{1}%
t_{1}+\omega_{2}t_{2}+\cdots+\omega_{n}t_{n}\right)  -\omega_{Q_{i}%
}\operatorname*{id}\nolimits_{\mathcal{S}^{\lambda}}\right)  \in
\operatorname*{End}\nolimits_{\mathbf{k}}\left(  \mathcal{S}^{\lambda}\right)
\]
\footnote{This product is well-defined (and does not depend on the order of
its factors), since all its factors (being polynomials in $L_{\lambda}\left(
\omega_{1}t_{1}+\omega_{2}t_{2}+\cdots+\omega_{n}t_{n}\right)  $) commute.}
sends the whole $\mathcal{S}^{\lambda}$ to $0$ (because its first factor sends
$\mathcal{S}^{\lambda}=F_{0}^{\lambda}$ down to $F_{1}^{\lambda}$, then its
second factor sends $F_{1}^{\lambda}$ further down to $F_{2}^{\lambda}$, then
its third factor sends $F_{2}^{\lambda}$ onward to $F_{3}^{\lambda}$, and so
on, until the last factor sends $F_{f_{n+1}-1}^{\lambda}$ down to $F_{f_{n+1}%
}^{\lambda}=0$). Moreover, for this to hold, we do not actually need all the
$f_{n+1}$ factors of this product, but rather only those factors that
correspond to the numbers $i\in\left[  f_{n+1}\right]  $ satisfying $c_{Q_{i}%
}^{\lambda}\neq0$ (because if $c_{Q_{i}}^{\lambda}=0$, then property 2 of our
filtration shows that $F_{i-1}^{\lambda}=F_{i}^{\lambda}$, and thus we don't
need to apply the $L_{\lambda}\left(  \omega_{1}t_{1}+\omega_{2}t_{2}%
+\cdots+\omega_{n}t_{n}\right)  -\omega_{Q_{i}}\operatorname*{id}%
\nolimits_{\mathcal{S}^{\lambda}}$ factor to send us from $F_{i-1}^{\lambda}$
down into $F_{i}^{\lambda}$). Hence, the operator
\[
\prod_{\substack{i\in\left[  f_{n+1}\right]  ;\\c_{Q_{i}}^{\lambda}\neq
0}}\left(  L_{\lambda}\left(  \omega_{1}t_{1}+\omega_{2}t_{2}+\cdots
+\omega_{n}t_{n}\right)  -\omega_{Q_{i}}\operatorname*{id}%
\nolimits_{\mathcal{S}^{\lambda}}\right)  \in\operatorname*{End}%
\nolimits_{\mathbf{k}}\left(  \mathcal{S}^{\lambda}\right)
\]
sends the whole $\mathcal{S}^{\lambda}$ to $0$ as well. In other words,
\[
\prod_{\substack{i\in\left[  f_{n+1}\right]  ;\\c_{Q_{i}}^{\lambda}\neq
0}}\left(  L_{\lambda}\left(  \omega_{1}t_{1}+\omega_{2}t_{2}+\cdots
+\omega_{n}t_{n}\right)  -\omega_{Q_{i}}\operatorname*{id}%
\nolimits_{\mathcal{S}^{\lambda}}\right)  =0.
\]
Since $Q_{1},Q_{2},\ldots,Q_{f_{n+1}}$ are just the lacunar subsets of
$\left[  n-1\right]  $, we can rewrite this as%
\[
\prod_{\substack{I\subseteq\left[  n-1\right]  \text{ is lacunar;}%
\\c_{I}^{\lambda}\neq0}}\left(  L_{\lambda}\left(  \omega_{1}t_{1}+\omega
_{2}t_{2}+\cdots+\omega_{n}t_{n}\right)  -\omega_{I}\operatorname*{id}%
\nolimits_{\mathcal{S}^{\lambda}}\right)  =0.
\]
This proves Theorem \ref{thm.eigs} \textbf{(c)} in the case when $\mathbf{k}$
is a field of characteristic $0$. Just as in our proof of part \textbf{(a)},
we can derive the general case from this case by a polynomial identity
argument (treating $\omega_{1},\omega_{2},\ldots,\omega_{n}$ as
indeterminates, and now considering polynomials with values in
$\operatorname*{End}\nolimits_{\mathbf{k}}\left(  \mathcal{S}^{\lambda
}\right)  $, which can be encoded as tuples of usual polynomials). \medskip

\textbf{(b)} Theorem \ref{thm.eigs} \textbf{(c)} shows that the endomorphism
$L_{\lambda}\left(  \omega_{1}t_{1}+\omega_{2}t_{2}+\cdots+\omega_{n}%
t_{n}\right)  \in\operatorname*{End}\nolimits_{\mathbf{k}}\left(
\mathcal{S}^{\lambda}\right)  $ is annihilated by the polynomial
$\prod_{\substack{I\subseteq\left[  n-1\right]  \text{ is lacunar;}%
\\c_{I}^{\lambda}\neq0}}\left(  x-\omega_{I}\right)  \in\mathbf{k}\left[
x\right]  $ (meaning that the polynomial vanishes when we substitute the
endomorphism for $x$). But it is well-known that a linear endomorphism (of a
finite-dimensional $\mathbf{k}$-vector space) that is annihilated by a
polynomial of the form $\prod\left(  x-r\right)  $ with pairwise distinct
scalars $r$ is always diagonalizable. Hence, if the $\omega_{I}$ in the above
polynomial are pairwise distinct, then the endomorphism $L_{\lambda}\left(
\omega_{1}t_{1}+\omega_{2}t_{2}+\cdots+\omega_{n}t_{n}\right)  $ is
diagonalizable. This proves Theorem \ref{thm.eigs} \textbf{(b)}.
\end{proof}


\section{Final remarks}

Thus we have computed the eigenvalues -- and their algebraic multiplicities --
for the action of any one-sided cycle shuffle $\omega_{1}t_{1}+\omega_{2}%
t_{2}+\cdots+\omega_{n}t_{n}$ on any Specht module $\mathcal{S}^{\lambda}$.
With a trivial amount of work, we could extend this analysis to the action of
any element of $\mathcal{T}$ (that is, of any noncommutative polynomial in
$t_{1},t_{2},\ldots,t_{n}$). This automatically allows us to identify the
eigenvalues of such elements on any $S_{n}$-representation $V$, as long as the
decomposition of $V$ into Specht modules is known.

The proof of our result was achieved in a rather roundabout way: We did no
work in the Specht modules $\mathcal{S}^{\lambda}$ themselves. Instead, we
used a filtration of $\mathcal{A}$ (the Fibonacci filtration) whose
subquotients $F_{i}/F_{i-1}$ we were able to decompose into Specht modules
(Theorem \ref{thm.Fi/Fi-1.as-ind}). Then, we \textquotedblleft
projected\textquotedblright\ this filtration onto each Specht module
$\mathcal{S}^{\lambda}$ (Theorem \ref{thm.specht.filtr}) and used the
semisimplicity of $\mathcal{A}$ (actually, the complete reducibility of
$\mathcal{S}^{\lambda}$ would have sufficed) to triangularize the action of
$\mathcal{T}$ on $\mathcal{S}^{\lambda}$. This is in contrast to other
instances of similar questions, such as the recent \cite{2407.08644}, where
the solution requires significant exploration of the inner life of
$\mathcal{S}^{\lambda}$.

Our above method for proving Theorems \ref{thm.specht.filtr} and
\ref{thm.eigs} -- in which we used the Fibonacci filtration to triangularize
$L_{\lambda}\left(  \omega_{1}t_{1}+\omega_{2}t_{2}+\cdots+\omega_{n}%
t_{n}\right)  $ -- is partly generalizable:

\begin{proposition}
\label{prop.reflT-gen}Let $A$ be a $\mathbf{k}$-algebra\footnotemark, and let
$T$ be a $\mathbf{k}$-subalgebra of $A$. Let%
\begin{equation}
0=F_{0}\subseteq F_{1}\subseteq F_{2}\subseteq\cdots\subseteq F_{m}=A
\label{eq.prop.reflT-gen.filA}%
\end{equation}
be a filtration of $A$ by $\left(  A,T\right)  $-subbimodules. Let $V$ be any
left $A$-module. If $B$ is any $\left(  A,T\right)  $-subbimodule of $A$, then
we can define a left $T$-submodule
\[
V^{B}:=\left\{  v\in V\ \mid\ Bv=0\right\}
\]
of $V$. Then, we have a filtration%
\begin{equation}
0=V^{F_{m}}\subseteq V^{F_{m-1}}\subseteq V^{F_{m-2}}\subseteq\cdots\subseteq
V^{F_{0}}=V \label{eq.prop.reflT-gen.filV}%
\end{equation}
of $V$ by left $T$-submodules.

\begin{enumerate}
\item[\textbf{(a)}] Its subquotients $V^{F_{i-1}}/V^{F_{i}}$ can be
canonically embedded into $\operatorname*{Hom}\nolimits_{A}\left(
F_{i}/F_{i-1},\ V\right)  $ as left $T$-modules. Thus, if some element $t\in
T$ acts triangularly from the right on the filtration
(\ref{eq.prop.reflT-gen.filA}) (meaning that it acts as a scalar on each
subquotient $F_{i}/F_{i-1}$), then it also acts triangularly from the left on
the filtration (\ref{eq.prop.reflT-gen.filV}).

\item[\textbf{(b)}] If $\mathbf{k}$ is a field and the algebra $A$ is
semisimple, then these embeddings $V^{F_{i-1}}/V^{F_{i}}\rightarrow
\operatorname*{Hom}\nolimits_{A}\left(  F_{i}/F_{i-1},\ V\right)  $ are
isomorphisms. Thus, in this case, knowing the dimensions of the Hom-spaces
$\operatorname*{Hom}\nolimits_{A}\left(  F_{i},\ V\right)  $ allows us to
compute the multiplicities of eigenvalues for a triangular $t\in T$ acting on
$V$.
\end{enumerate}
\end{proposition}

\footnotetext{Recall that $\mathbf{k}$ is an arbitrary commutative ring.}

\begin{proof}
This is implicit in our above proofs of Theorems \ref{thm.specht.filtr} and
\ref{thm.eigs} (where $A$, $T$ and $V$ were taken to be $\mathcal{A}$,
$\mathcal{T}$ and $\mathcal{S}^{\lambda}$, and where the submodules $V^{F_{i}%
}$ were called $F_{i}^{\lambda}$).
\end{proof}

\appendix



\section{\label{sec.apx.proofs}Omitted proofs}

In this appendix, we give proofs to some folklore facts about representations
of symmetric groups.

\subsection{\label{subsec.apx.pf.prop.Z-iso}Proof of Proposition
\ref{prop.Z-iso}}

\begin{proof}
[Proof of Proposition \ref{prop.Z-iso} (sketched).]We will use the notations
from \cite{sga}.

Let $\left(  e_{1},e_{2},\ldots,e_{n}\right)  $ be the standard basis of the
natural representation $\mathcal{N}_{n}=\mathbf{k}^{n}$. The symmetric group
$S_{n}$ acts on it by the rule $\sigma\cdot e_{i}=e_{\sigma\left(  i\right)
}$ for all $\sigma\in S_{n}$ and $i\in\left[  n\right]  $. The
subrepresentation $\mathcal{D}_{n}$ is then spanned by the single vector
$e_{1}+e_{2}+\cdots+e_{n}=\left(  1,1,\ldots,1\right)  $.

Consider the Young diagram $Y\left(  \left(  n-1,1\right)  \right)  $ of the
partition $\left(  n-1,1\right)  $. This Young diagram gives rise to a Young
module $\mathcal{M}^{\left(  n-1,1\right)  }:=\mathcal{M}^{Y\left(  \left(
n-1,1\right)  \right)  }$ and a Specht module $\mathcal{S}^{\left(
n-1,1\right)  }:=\mathcal{S}^{Y\left(  \left(  n-1,1\right)  \right)  }$.

The Young module $\mathcal{M}^{\left(  n-1,1\right)  }$ has a basis formed by
the $n$-tabloids of shape $Y\left(  \left(  n-1,1\right)  \right)  $. We will
use the short notation $\underline{\overline{k}}$ for the $n$-tabloid of shape
$Y\left(  \left(  n-1,1\right)  \right)  $ that has the entry $k$ in cell
$\left(  2,1\right)  $, as in \cite[Example 5.3.16]{sga}. For instance,%
\[
\text{if }n=5\text{, then }\underline{\overline{3}}%
=\ytableausetup{tabloids}\ytableaushort{1245,3}\ytableausetup{notabloids}\ \ .
\]


The Specht module $\mathcal{S}^{\left(  n-1,1\right)  }$ is the submodule of
$\mathcal{M}^{\left(  n-1,1\right)  }$ spanned by the polytabloids
$\mathbf{e}_{T}=\underline{\overline{T\left(  2,1\right)  }}%
-\underline{\overline{T\left(  1,1\right)  }}$, or, equivalently, by the
differences $\underline{\overline{k}}-\underline{\overline{\ell}}$ for
$k\neq\ell$ in $\left[  n\right]  $. (See \cite[Example 5.4.2]{sga}.)

The $S_{n}$-representation $\mathcal{M}^{\left(  n-1,1\right)  }$ is
isomorphic to the natural representation $\mathcal{N}_{n}$ via the isomorphism%
\begin{align*}
\mathcal{M}^{\left(  n-1,1\right)  }  &  \rightarrow\mathcal{N}_{n},\\
\underline{\overline{k}}  &  \mapsto e_{k}.
\end{align*}
Under this isomorphism, the Specht module $\mathcal{S}^{\left(  n-1,1\right)
}$ becomes the span of the differences $e_{k}-e_{\ell}$ for $k\neq\ell$ in
$\left[  n\right]  $. Hence, $\mathcal{S}^{\left(  n-1,1\right)  }$ is
isomorphic to the zero-sum subrepresentation of $\mathcal{N}_{n}$ (called
$R\left(  \mathbf{k}^{n}\right)  $ in \cite[\S 4.2.5]{sga}). \medskip

\begin{noncompile}
\textit{Old proof of Proposition \ref{prop.Z-iso} \textbf{(a)}:} We define the
$\mathbf{k}$-bilinear form%
\begin{align*}
f:\mathcal{N}_{n}\times\mathcal{N}_{n}  &  \rightarrow\mathbf{k},\\
\left(  e_{i},e_{j}\right)   &  \mapsto\delta_{i,j}\ \ \ \ \ \ \ \ \ \ \left(
\text{Kronecker delta}\right)  .
\end{align*}
This form $f$ is $S_{n}$-invariant and nondegenerate in the strongest sense
(i.e., it gives rise to an isomorphism $\mathcal{N}_{n}\rightarrow
\mathcal{N}_{n}^{\ast}$). Restricting it to the subset $\mathcal{S}^{\left(
n-1,1\right)  }\times\mathcal{N}_{n}$, we obtain a $\mathbf{k}$-bilinear form%
\begin{align*}
f^{\prime}:\mathcal{S}^{\left(  n-1,1\right)  }\times\mathcal{N}_{n}  &
\rightarrow\mathbf{k},\\
\left(  \underline{\overline{k}}-\underline{\overline{\ell}},e_{j}\right)   &
\mapsto\delta_{k,j}-\delta_{\ell,j},
\end{align*}
which gives rise to a $\mathbf{k}$-module morphism
\begin{align*}
f_{R}^{\prime}:\mathcal{N}_{n}  &  \rightarrow\left(  \mathcal{S}^{\left(
n-1,1\right)  }\right)  ^{\ast},\\
v  &  \mapsto\left(  u\mapsto f^{\prime}\left(  u,v\right)  \right)  .
\end{align*}
This morphism $f_{R}^{\prime}$ is a morphism of $S_{n}$-representations (since
the form $f$ and thus the form $f^{\prime}$ is $S_{n}$-invariant), i.e., a
left $\mathcal{A}$-module morphism. Moreover, its kernel is%
\begin{align*}
\operatorname*{Ker}\left(  f_{R}^{\prime}\right)   &  =\left\{  v\in
\mathcal{N}_{n}\ \mid\ f^{\prime}\left(  u,v\right)  =0\text{ for all }%
u\in\mathcal{S}^{\left(  n-1,1\right)  }\right\} \\
&  =\left\{  v\in\mathcal{N}_{n}\ \mid\ f^{\prime}\left(  \underline{\overline
{k}}-\underline{\overline{\ell}},v\right)  =0\text{ for all }k\neq\ell\right\}
\\
&  =\left\{  v\in\mathcal{N}_{n}\ \mid\ v_{k}-v_{\ell}=0\text{ for all }%
k\neq\ell\right\} \\
&  \ \ \ \ \ \ \ \ \ \ \ \ \ \ \ \ \ \ \ \ \left(
\begin{array}
[c]{c}%
\text{where }v_{i}\text{ denotes the }i\text{-th entry of}\\
\text{the vector }v\in\mathcal{N}_{n}=\mathbf{k}^{n}%
\end{array}
\right) \\
&  =\mathcal{D}_{n}.
\end{align*}
Hence, it factors through the quotient $\mathcal{N}_{n}/\mathcal{D}%
_{n}=\mathcal{Z}_{n}$. Thus, we obtain a left $\mathcal{A}$-module morphism%
\begin{align*}
f_{R}^{\prime\prime}:\mathcal{Z}_{n}  &  \rightarrow\left(  \mathcal{S}%
^{\left(  n-1,1\right)  }\right)  ^{\ast},\\
\overline{v}  &  \mapsto\left(  u\mapsto f^{\prime}\left(  u,v\right)
\right)  .
\end{align*}
This morphism $f_{R}^{\prime\prime}$ is injective (since it is obtained from
$f_{R}^{\prime}$ by factoring out the whole kernel $\operatorname*{Ker}\left(
f_{R}^{\prime}\right)  =\mathcal{D}_{n}$) and surjective (since it sends each
basis vector $\overline{e_{i}}$ to the linear form on $\mathcal{S}^{\left(
n-1,1\right)  }$ that sends each vector to its $i$-th coordinate; but the
latter forms span $\left(  \mathcal{S}^{\left(  n-1,1\right)  }\right)
^{\ast}$), hence bijective. Thus, it is a left $\mathcal{A}$-module
isomorphism, i.e., an isomorphism of $S_{n}$-representations. This proves
Proposition \ref{prop.Z-iso} \textbf{(a)}. \medskip
\end{noncompile}

\textbf{(a)} Let $\left(  e_{1}^{\ast},e_{2}^{\ast},\ldots,e_{n}^{\ast
}\right)  $ be the dual basis of the standard basis $\left(  e_{1}%
,e_{2},\ldots,e_{n}\right)  $ of $\mathcal{N}_{n}$. Thus, each $i\in\left[
n\right]  $ satisfies $e_{i}^{\ast}\in\mathcal{N}_{n}^{\ast}$, and each
$i,j\in\left[  n\right]  $ satisfy $e_{i}^{\ast}\left(  e_{j}\right)
=\delta_{i,j}$ (Kronecker delta). Furthermore, it is straightforward to see
that $\sigma\left(  e_{i}^{\ast}\right)  =e_{\sigma\left(  i\right)  }^{\ast}$
for each $\sigma\in S_{n}$ and $i\in\left[  n\right]  $.

Now, consider the $\mathbf{k}$-linear map%
\begin{align*}
\Phi:\mathcal{M}^{\left(  n-1,1\right)  }  &  \rightarrow\mathcal{N}_{n}%
^{\ast},\\
\underline{\overline{k}}  &  \mapsto e_{k}^{\ast}\ \ \ \ \ \ \ \ \ \ \left(
\text{for all }k\in\left[  n\right]  \right)  .
\end{align*}
This map $\Phi$ is invertible (since $\left(  \underline{\overline{k}}\right)
_{k\in\left[  n\right]  }$ and $\left(  e_{k}^{\ast}\right)  _{k\in\left[
n\right]  }$ are bases of the $\mathbf{k}$-modules $\mathcal{M}^{\left(
n-1,1\right)  }$ and $\mathcal{N}_{n}^{\ast}$, respectively) and $S_{n}%
$-equivariant (since each $k\in\left[  n\right]  $ and $\sigma\in S_{n}$
satisfy $\sigma\left(  \underline{\overline{k}}\right)  =\underline{\overline
{\sigma\left(  k\right)  }}$ and $\sigma\left(  e_{k}^{\ast}\right)
=e_{\sigma\left(  k\right)  }^{\ast}$). Hence, this map $\Phi$ is an
isomorphism of $S_{n}$-representations.

Now, consider the subrepresentation $\mathcal{S}^{\left(  n-1,1\right)
}=\operatorname*{span}\nolimits_{\mathbf{k}}\left\{  \underline{\overline{k}%
}-\underline{\overline{\ell}}\ \mid\ k\neq\ell\right\}  $ of $\mathcal{M}%
^{\left(  n-1,1\right)  }$. The image of this subrepresentation under $\Phi$
is
\begin{align}
&  \Phi\left(  \mathcal{S}^{\left(  n-1,1\right)  }\right) \nonumber\\
&  =\Phi\left(  \operatorname*{span}\nolimits_{\mathbf{k}}\left\{
\underline{\overline{k}}-\underline{\overline{\ell}}\ \mid\ k\neq\ell\right\}
\right)  \ \ \ \ \ \ \ \ \ \ \left(  \text{since }\mathcal{S}^{\left(
n-1,1\right)  }=\operatorname*{span}\nolimits_{\mathbf{k}}\left\{
\underline{\overline{k}}-\underline{\overline{\ell}}\ \mid\ k\neq\ell\right\}
\right) \nonumber\\
&  =\operatorname*{span}\nolimits_{\mathbf{k}}\left\{  \Phi\left(
\underline{\overline{k}}-\underline{\overline{\ell}}\right)  \ \mid\ k\neq
\ell\right\}  \ \ \ \ \ \ \ \ \ \ \left(  \text{since }\Phi\text{ is
}\mathbf{k}\text{-linear}\right) \nonumber\\
&  =\operatorname*{span}\nolimits_{\mathbf{k}}\left\{  e_{k}^{\ast}-e_{\ell
}^{\ast}\ \mid\ k\neq\ell\right\}  \label{pf.prop.Z-iso.a.2}%
\end{align}
(since every $k,\ell\in\left[  n\right]  $ satisfy $\Phi\left(
\underline{\overline{k}}-\underline{\overline{\ell}}\right)  =\Phi\left(
\underline{\overline{k}}\right)  -\Phi\left(  \underline{\overline{\ell}%
}\right)  =e_{k}^{\ast}-e_{\ell}^{\ast}$ by the definition of $\Phi$).

Now, let us recall some general properties of dual modules. If $V$ is any
$\mathbf{k}$-module, and if $W$ is a $\mathbf{k}$-submodule of $V$, then we
let $W^{\perp}$ denote the $\mathbf{k}$-submodule $\left\{  f\in V^{\ast
}\ \mid\ f\left(  W\right)  =0\right\}  $ of $V^{\ast}$. It is well-known that
there is a canonical isomorphism $W^{\perp}\cong\left(  V/W\right)  ^{\ast}$
in this situation (since any linear map $f\in W^{\perp}$ annihilates $W$ and
thus can be factored through $V/W$ by the first isomorphism theorem).
Moreover, if $V$ is a representation of a group $G$, and if $W$ is a
subrepresentation of $V$, then $W^{\perp}$ is a subrepresentation of $V^{\ast
}$, and the above-mentioned isomorphism $W^{\perp}\cong\left(  V/W\right)
^{\ast}$ is an isomorphism of $G$-representations. Applying this fact to
$G=S_{n}$ and $V=\mathcal{N}_{n}$ and $W=\mathcal{D}_{n}$, we obtain the
isomorphism%
\begin{equation}
\mathcal{D}_{n}^{\perp}\cong\left(  \mathcal{N}_{n}/\mathcal{D}_{n}\right)
^{\ast} \label{pf.prop.Z-iso.a.3}%
\end{equation}
of $S_{n}$-representations.

Now we shall show that $\operatorname*{span}\nolimits_{\mathbf{k}}\left\{
e_{k}^{\ast}-e_{\ell}^{\ast}\ \mid\ k\neq\ell\right\}  =\mathcal{D}_{n}%
^{\perp}$. Indeed,
\begin{align}
\mathcal{D}_{n}^{\perp}  &  =\left\{  f\in\mathcal{N}_{n}^{\ast}%
\ \mid\ f\left(  \mathcal{D}_{n}\right)  =0\right\}
\ \ \ \ \ \ \ \ \ \ \left(  \text{by the definition of }\mathcal{D}_{n}%
^{\perp}\right) \nonumber\\
&  =\left\{  f\in\mathcal{N}_{n}^{\ast}\ \mid\ f\left(  \left(  1,1,\ldots
,1\right)  \right)  =0\right\}  \label{pf.prop.Z-iso.a.4}%
\end{align}
(since the $\mathbf{k}$-module $\mathcal{D}_{n}$ is spanned by the single
vector $\left(  1,1,\ldots,1\right)  $, and thus a $\mathbf{k}$-linear map
$f\in\mathcal{N}_{n}^{\ast}$ satisfies $f\left(  \mathcal{D}_{n}\right)  =0$
if and only if it satisfies $f\left(  \left(  1,1,\ldots,1\right)  \right)
=0$). For any $k\neq\ell$, we have $\left(  e_{k}^{\ast}-e_{\ell}^{\ast
}\right)  \left(  \left(  1,1,\ldots,1\right)  \right)  =1-1=0$ and thus
$e_{k}^{\ast}-e_{\ell}^{\ast}\in\mathcal{D}_{n}^{\perp}$ (by
(\ref{pf.prop.Z-iso.a.4})). Hence, by linearity, it follows that
\begin{equation}
\operatorname*{span}\nolimits_{\mathbf{k}}\left\{  e_{k}^{\ast}-e_{\ell}%
^{\ast}\ \mid\ k\neq\ell\right\}  \subseteq\mathcal{D}_{n}^{\perp}.
\label{pf.prop.Z-iso.a.6}%
\end{equation}


Let us now show the converse inclusion. Indeed, let $f\in\mathcal{D}%
_{n}^{\perp}$. Then, $f\in\mathcal{N}_{n}^{\ast}$ and $f\left(  \left(
1,1,\ldots,1\right)  \right)  =0$ (by (\ref{pf.prop.Z-iso.a.4})). But $\left(
e_{1}^{\ast},e_{2}^{\ast},\ldots,e_{n}^{\ast}\right)  $ is a basis of
$\mathcal{N}_{n}^{\ast}$; hence, $f$ can be written as a $\mathbf{k}$-linear
combination $f=\alpha_{1}e_{1}^{\ast}+\alpha_{2}e_{2}^{\ast}+\cdots+\alpha
_{n}e_{n}^{\ast}$ with $\alpha_{1},\alpha_{2},\ldots,\alpha_{n}\in\mathbf{k}$
(since $f\in\mathcal{N}_{n}^{\ast}$). Consider these $\alpha_{1},\alpha
_{2},\ldots,\alpha_{n}$. From $f=\alpha_{1}e_{1}^{\ast}+\alpha_{2}e_{2}^{\ast
}+\cdots+\alpha_{n}e_{n}^{\ast}$, we obtain
\begin{align*}
f\left(  \left(  1,1,\ldots,1\right)  \right)   &  =\left(  \alpha_{1}%
e_{1}^{\ast}+\alpha_{2}e_{2}^{\ast}+\cdots+\alpha_{n}e_{n}^{\ast}\right)
\left(  \left(  1,1,\ldots,1\right)  \right) \\
&  =\alpha_{1}1+\alpha_{2}1+\cdots+\alpha_{n}1=\alpha_{1}+\alpha_{2}%
+\cdots+\alpha_{n}.
\end{align*}
Thus, $\alpha_{1}+\alpha_{2}+\cdots+\alpha_{n}=f\left(  \left(  1,1,\ldots
,1\right)  \right)  =0$, so that $\alpha_{1}=-\alpha_{2}-\alpha_{3}%
-\cdots-\alpha_{n}$. Now,%
\begin{align*}
f  &  =\underbrace{\alpha_{1}}_{=-\alpha_{2}-\alpha_{3}-\cdots-\alpha_{n}%
}e_{1}^{\ast}+\alpha_{2}e_{2}^{\ast}+\cdots+\alpha_{n}e_{n}^{\ast}\\
&  =\left(  -\alpha_{2}-\alpha_{3}-\cdots-\alpha_{n}\right)  e_{1}^{\ast
}+\alpha_{2}e_{2}^{\ast}+\alpha_{3}e_{3}^{\ast}+\cdots+\alpha_{n}e_{n}^{\ast
}\\
&  =\alpha_{2}\underbrace{\left(  e_{2}^{\ast}-e_{1}^{\ast}\right)  }%
_{\in\left\{  e_{k}^{\ast}-e_{\ell}^{\ast}\ \mid\ k\neq\ell\right\}
}+\,\alpha_{3}\underbrace{\left(  e_{3}^{\ast}-e_{1}^{\ast}\right)  }%
_{\in\left\{  e_{k}^{\ast}-e_{\ell}^{\ast}\ \mid\ k\neq\ell\right\}  }%
+\cdots+\alpha_{n}\underbrace{\left(  e_{n}^{\ast}-e_{1}^{\ast}\right)  }%
_{\in\left\{  e_{k}^{\ast}-e_{\ell}^{\ast}\ \mid\ k\neq\ell\right\}  }\\
&  \in\operatorname*{span}\nolimits_{\mathbf{k}}\left\{  e_{k}^{\ast}-e_{\ell
}^{\ast}\ \mid\ k\neq\ell\right\}  .
\end{align*}


Forget that we fixed $f$. We thus have shown that $f\in\operatorname*{span}%
\nolimits_{\mathbf{k}}\left\{  e_{k}^{\ast}-e_{\ell}^{\ast}\ \mid\ k\neq
\ell\right\}  $ for each $f\in\mathcal{D}_{n}^{\perp}$. In other words,
$\mathcal{D}_{n}^{\perp}\subseteq\operatorname*{span}\nolimits_{\mathbf{k}%
}\left\{  e_{k}^{\ast}-e_{\ell}^{\ast}\ \mid\ k\neq\ell\right\}  $. Combining
this with (\ref{pf.prop.Z-iso.a.6}), we obtain%
\[
\mathcal{D}_{n}^{\perp}=\operatorname*{span}\nolimits_{\mathbf{k}}\left\{
e_{k}^{\ast}-e_{\ell}^{\ast}\ \mid\ k\neq\ell\right\}  .
\]
Comparing this with (\ref{pf.prop.Z-iso.a.2}), we are led to%
\begin{align*}
\Phi\left(  \mathcal{S}^{\left(  n-1,1\right)  }\right)   &  =\mathcal{D}%
_{n}^{\perp}\cong\left(  \mathcal{N}_{n}/\mathcal{D}_{n}\right)  ^{\ast
}\ \ \ \ \ \ \ \ \ \ \left(  \text{by (\ref{pf.prop.Z-iso.a.3})}\right) \\
&  \cong\mathcal{Z}_{n}^{\ast}\ \ \ \ \ \ \ \ \ \ \left(  \text{since
}\mathcal{N}_{n}/\mathcal{D}_{n}=\mathcal{Z}_{n}\right)  ,
\end{align*}
and this is an isomorphism of $S_{n}$-representations (as we saw above). On
the other hand, however, we also have $\Phi\left(  \mathcal{S}^{\left(
n-1,1\right)  }\right)  \cong\mathcal{S}^{\left(  n-1,1\right)  }$ as $S_{n}%
$-representations (since $\Phi$ is an isomorphism of $S_{n}$-representations).
Thus,%
\[
\mathcal{S}^{\left(  n-1,1\right)  }\cong\Phi\left(  \mathcal{S}^{\left(
n-1,1\right)  }\right)  \cong\mathcal{Z}_{n}^{\ast}%
\ \ \ \ \ \ \ \ \ \ \text{as }S_{n}\text{-representations.}%
\]
Taking duals, we thus obtain%
\begin{equation}
\left(  \mathcal{S}^{\left(  n-1,1\right)  }\right)  ^{\ast}\cong\left(
\mathcal{Z}_{n}^{\ast}\right)  ^{\ast}\ \ \ \ \ \ \ \ \ \ \text{as }%
S_{n}\text{-representations.} \label{pf.prop.Z-iso.a.dual}%
\end{equation}


But the $\mathbf{k}$-module $\mathcal{Z}_{n}=\mathcal{N}_{n}/\mathcal{D}%
_{n}=\mathbf{k}^{n}/\operatorname*{span}\nolimits_{\mathbf{k}}\left\{  \left(
1,1,\ldots,1\right)  \right\}  $ has a finite basis (namely, $\left(
\overline{e_{1}},\overline{e_{2}},\ldots,\overline{e_{n-1}}\right)  $, as can
be easily seen from basic linear algebra). Hence, \cite[Proposition 5.19.22
\textbf{(b)}]{sga} shows that $\left(  \mathcal{Z}_{n}^{\ast}\right)  ^{\ast
}\cong\mathcal{Z}_{n}$ as $S_{n}$-representations. In view of this, we can
rewrite (\ref{pf.prop.Z-iso.a.dual}) as follows:%
\[
\left(  \mathcal{S}^{\left(  n-1,1\right)  }\right)  ^{\ast}\cong%
\mathcal{Z}_{n}\ \ \ \ \ \ \ \ \ \ \text{as }S_{n}\text{-representations.}%
\]
This proves Proposition \ref{prop.Z-iso} \textbf{(a)}. \medskip

\textbf{(b)} Assume that $n$ is invertible in $\mathbf{k}$. Then,
\cite[Proposition 4.2.28 \textbf{(a)}]{sga} says that $\mathbf{k}^{n}=R\left(
\mathbf{k}^{n}\right)  \oplus D\left(  \mathbf{k}^{n}\right)  $, where we are
using the notations of \cite{sga}. In our notations, this is saying that
$\mathcal{N}_{n}=\mathcal{S}^{\left(  n-1,1\right)  }\oplus\mathcal{D}_{n}$
(since the submodule $\mathcal{S}^{\left(  n-1,1\right)  }%
=\operatorname*{span}\nolimits_{\mathbf{k}}\left\{  \underline{\overline{k}%
}-\underline{\overline{\ell}}\ \mid\ k\neq\ell\right\}  $ of $\mathcal{M}%
^{\left(  n-1,1\right)  }$ corresponds to the zero-sum subrepresentation
$R\left(  \mathbf{k}^{n}\right)  =\operatorname*{span}\nolimits_{\mathbf{k}%
}\left\{  e_{k}-e_{\ell}\ \mid\ k\neq\ell\right\}  $ of $\mathbf{k}^{n}$).
Hence, $\mathcal{S}^{\left(  n-1,1\right)  }\cong\mathcal{N}_{n}%
/\mathcal{D}_{n}=\mathcal{Z}_{n}$ as $S_{n}$-representations. This proves
Proposition \ref{prop.Z-iso} \textbf{(b)}.

Alternatively, we can prove Proposition \ref{prop.Z-iso} \textbf{(b)}
directly: Let
\[
\underline{\overline{\operatorname*{avg}}}:=\dfrac{1}{n}\left(
\underline{\overline{1}}+\underline{\overline{2}}+\cdots+\underline{\overline
{n}}\right)  =\dfrac{1}{n}\sum_{j=1}^{n}\underline{\overline{j}}\in
\mathcal{M}^{\left(  n-1,1\right)  }.
\]
Thus, $n\cdot\underline{\overline{\operatorname*{avg}}}=\sum_{j=1}%
^{n}\underline{\overline{j}}=\sum_{k=1}^{n}\underline{\overline{k}}$. Hence,%
\begin{align}
&  \left(  \underline{\overline{1}}-\underline{\overline{\operatorname*{avg}}%
}\right)  +\left(  \underline{\overline{2}}-\underline{\overline
{\operatorname*{avg}}}\right)  +\cdots+\left(  \underline{\overline{n}%
}-\underline{\overline{\operatorname*{avg}}}\right) \nonumber\\
&  =\sum_{k=1}^{n}\left(  \underline{\overline{k}}-\underline{\overline
{\operatorname*{avg}}}\right)  =\sum_{k=1}^{n}\underline{\overline{k}}%
-n\cdot\underline{\overline{\operatorname*{avg}}}=0 \label{pf.prop.Z-iso.b.2}%
\end{align}
(since $n\cdot\underline{\overline{\operatorname*{avg}}}=\sum_{k=1}%
^{n}\underline{\overline{k}}$).

It is furthermore easy to see that each $i\in\left[  n\right]  $ satisfies
$\underline{\overline{i}}-\underline{\overline{\operatorname*{avg}}}%
\in\mathcal{S}^{\left(  n-1,1\right)  }$ (this follows from an easy
computation\footnote{\textit{Proof.} Let $i\in\left[  n\right]  $. Then, from
$\underline{\overline{\operatorname*{avg}}}=\dfrac{1}{n}\sum_{j=1}%
^{n}\underline{\overline{j}}$, we obtain%
\[
\underline{\overline{i}}-\underline{\overline{\operatorname*{avg}}%
}=\underline{\overline{i}}-\dfrac{1}{n}\sum_{j=1}^{n}\underline{\overline{j}%
}=\dfrac{1}{n}\sum_{j=1}^{n}\underbrace{\left(  \underline{\overline{i}%
}-\underline{\overline{j}}\right)  }_{\substack{\in\mathcal{S}^{\left(
n-1,1\right)  }\\\text{(since }\mathcal{S}^{\left(  n-1,1\right)
}=\operatorname*{span}\nolimits_{\mathbf{k}}\left\{  \underline{\overline{k}%
}-\underline{\overline{\ell}}\ \mid\ k\neq\ell\right\}  \text{)}}%
}\in\mathcal{S}^{\left(  n-1,1\right)  }%
\]
(since $\mathcal{S}^{\left(  n-1,1\right)  }$ is a $\mathbf{k}$-module),
qed.}). Moreover, the list $\left(  \underline{\overline{1}}%
-\underline{\overline{\operatorname*{avg}}},\ \underline{\overline{2}%
}-\underline{\overline{\operatorname*{avg}}},\ \ldots,\ \underline{\overline
{n-1}}-\underline{\overline{\operatorname*{avg}}}\right)  $ is a basis of the
$\mathbf{k}$-module $\mathcal{S}^{\left(  n-1,1\right)  }$ (this is an easy
exercise in linear algebra\footnote{For the sake of completeness, here is a
\textit{proof:}
\par
\begin{itemize}
\item The $n-1$ vectors $\underline{\overline{1}}-\underline{\overline
{\operatorname*{avg}}},\ \underline{\overline{2}}-\underline{\overline
{\operatorname*{avg}}},\ \ldots,\ \underline{\overline{n-1}}%
-\underline{\overline{\operatorname*{avg}}}$ are $\mathbf{k}$-linearly
independent.
\par
\textit{Proof:} Let $\alpha_{1},\alpha_{2},\ldots,\alpha_{n-1}\in\mathbf{k}$
be scalars satisfying $\sum_{k=1}^{n-1}\alpha_{k}\left(  \underline{\overline
{k}}-\underline{\overline{\operatorname*{avg}}}\right)  =0$. We must show that
all coefficients $\alpha_{1},\alpha_{2},\ldots,\alpha_{n-1}$ are $0$.
\par
Recall that $\mathcal{M}^{\left(  n-1,1\right)  }$ is a free $\mathbf{k}%
$-module with basis $\left(  \underline{\overline{1}},\underline{\overline{2}%
},\ldots,\underline{\overline{n}}\right)  $. Comparing the coefficients of
$\underline{\overline{n}}$ on both sides of the equality $\sum_{k=1}%
^{n-1}\alpha_{k}\left(  \underline{\overline{k}}-\underline{\overline
{\operatorname*{avg}}}\right)  =0$, we obtain $-\sum_{k=1}^{n-1}\alpha
_{k}\cdot\dfrac{1}{n}=0$ (since the coefficient of $\underline{\overline{n}}$
in $\underline{\overline{\operatorname*{avg}}}$ is $\dfrac{1}{n}$, whereas the
coefficient of $\underline{\overline{n}}$ in $\underline{\overline{k}}$ for
any $k\in\left[  n-1\right]  $ is $0$). Multiplying the latter equality by
$-n$, we find $\sum_{k=1}^{n-1}\alpha_{k}=0$. Hence,%
\[
\sum_{k=1}^{n-1}\alpha_{k}\left(  \underline{\overline{k}}%
-\underline{\overline{\operatorname*{avg}}}\right)  =\sum_{k=1}^{n-1}%
\alpha_{k}\underline{\overline{k}}-\underbrace{\sum_{k=1}^{n-1}\alpha_{k}%
}_{=0}\underline{\overline{\operatorname*{avg}}}=\sum_{k=1}^{n-1}\alpha
_{k}\underline{\overline{k}},
\]
so that $\sum_{k=1}^{n-1}\alpha_{k}\underline{\overline{k}}=\sum_{k=1}%
^{n-1}\alpha_{k}\left(  \underline{\overline{k}}-\underline{\overline
{\operatorname*{avg}}}\right)  =0$. Since the basis vectors
$\underline{\overline{1}},\underline{\overline{2}},\ldots,\underline{\overline
{n-1}}$ are $\mathbf{k}$-linearly independent, this entails that all
coefficients $\alpha_{1},\alpha_{2},\ldots,\alpha_{n-1}$ are $0$. This
completes the proof of the linear independence of the vectors
$\underline{\overline{1}}-\underline{\overline{\operatorname*{avg}}%
},\ \underline{\overline{2}}-\underline{\overline{\operatorname*{avg}}%
},\ \ldots,\ \underline{\overline{n-1}}-\underline{\overline
{\operatorname*{avg}}}$.
\par
\item The $n-1$ vectors $\underline{\overline{1}}-\underline{\overline
{\operatorname*{avg}}},\ \underline{\overline{2}}-\underline{\overline
{\operatorname*{avg}}},\ \ldots,\ \underline{\overline{n-1}}%
-\underline{\overline{\operatorname*{avg}}}$ span the $\mathbf{k}$-module
$\mathcal{S}^{\left(  n-1,1\right)  }$.
\par
\textit{Proof:} Consider the two subsets%
\begin{align*}
W  &  :=\left\{  \underline{\overline{i}}-\underline{\overline
{\operatorname*{avg}}}\ \mid\ i\in\left[  n\right]  \right\}  =\left\{
\underline{\overline{1}}-\underline{\overline{\operatorname*{avg}}%
},\ \underline{\overline{2}}-\underline{\overline{\operatorname*{avg}}%
},\ \ldots,\ \underline{\overline{n}}-\underline{\overline{\operatorname*{avg}%
}}\right\}  \ \ \ \ \ \ \ \ \ \ \text{and}\\
W^{\prime}  &  :=\left\{  \underline{\overline{i}}-\underline{\overline
{\operatorname*{avg}}}\ \mid\ i\in\left[  n-1\right]  \right\}  =\left\{
\underline{\overline{1}}-\underline{\overline{\operatorname*{avg}}%
},\ \underline{\overline{2}}-\underline{\overline{\operatorname*{avg}}%
},\ \ldots,\ \underline{\overline{n-1}}-\underline{\overline
{\operatorname*{avg}}}\right\}
\end{align*}
of $\mathcal{M}^{\left(  n-1,1\right)  }$. Then, $W\subseteq\mathcal{S}%
^{\left(  n-1,1\right)  }$ (since each $i\in\left[  n\right]  $ satisfies
$\underline{\overline{i}}-\underline{\overline{\operatorname*{avg}}}%
\in\mathcal{S}^{\left(  n-1,1\right)  }$), so that $\operatorname*{span}%
W\subseteq\mathcal{S}^{\left(  n-1,1\right)  }$ (by linearity). Similarly,
$\operatorname*{span}W^{\prime}\subseteq\mathcal{S}^{\left(  n-1,1\right)  }$.
Also, the set $W$ differs from $W^{\prime}$ only in the extra element
$\underline{\overline{n}}-\underline{\overline{\operatorname*{avg}}}$. Thus,
$W=W^{\prime}\cup\left\{  \underline{\overline{n}}-\underline{\overline
{\operatorname*{avg}}}\right\}  $. Therefore,%
\[
\operatorname*{span}W=\operatorname*{span}\left(  W^{\prime}\cup\left\{
\underline{\overline{n}}-\underline{\overline{\operatorname*{avg}}}\right\}
\right)  =\operatorname*{span}W^{\prime}+\operatorname*{span}\left\{
\underline{\overline{n}}-\underline{\overline{\operatorname*{avg}}}\right\}
.
\]
However, solving the equation (\ref{pf.prop.Z-iso.b.2}) for
$\underline{\overline{n}}-\underline{\overline{\operatorname*{avg}}}$, we find%
\begin{align*}
\underline{\overline{n}}-\underline{\overline{\operatorname*{avg}}}  &
=-\left(  \underline{\overline{1}}-\underline{\overline{\operatorname*{avg}}%
}\right)  -\left(  \underline{\overline{2}}-\underline{\overline
{\operatorname*{avg}}}\right)  -\cdots-\left(  \underline{\overline{n-1}%
}-\underline{\overline{\operatorname*{avg}}}\right) \\
&  \in\operatorname*{span}\underbrace{\left\{  \underline{\overline{1}%
}-\underline{\overline{\operatorname*{avg}}},\ \underline{\overline{2}%
}-\underline{\overline{\operatorname*{avg}}},\ \ldots,\ \underline{\overline
{n-1}}-\underline{\overline{\operatorname*{avg}}}\right\}  }_{=W^{\prime}%
}=\operatorname*{span}W^{\prime}.
\end{align*}
Therefore, $\operatorname*{span}\left\{  \underline{\overline{n}%
}-\underline{\overline{\operatorname*{avg}}}\right\}  \subseteq
\operatorname*{span}W^{\prime}$, so that%
\[
\operatorname*{span}W=\operatorname*{span}W^{\prime}%
+\underbrace{\operatorname*{span}\left\{  \underline{\overline{n}%
}-\underline{\overline{\operatorname*{avg}}}\right\}  }_{\subseteq
\operatorname*{span}W^{\prime}}\subseteq\operatorname*{span}W^{\prime
}+\operatorname*{span}W^{\prime}\subseteq\operatorname*{span}W^{\prime}.
\]
\par
Recall that the $\mathbf{k}$-module $\mathcal{S}^{\left(  n-1,1\right)  }$ is
spanned by the differences $\underline{\overline{k}}-\underline{\overline
{\ell}}$ with $k\neq\ell$. But each of the latter differences
$\underline{\overline{k}}-\underline{\overline{\ell}}$ belongs to
$\operatorname*{span}W$ (indeed, $\underline{\overline{k}}%
-\underline{\overline{\ell}}=\underbrace{\left(  \underline{\overline{k}%
}-\underline{\overline{\operatorname*{avg}}}\right)  }_{\in W}%
-\underbrace{\left(  \underline{\overline{\ell}}-\underline{\overline
{\operatorname*{avg}}}\right)  }_{\in W}\in W-W\subseteq\operatorname*{span}%
W$). Combining the previous two sentences, we conclude (by linearity) that
$\mathcal{S}^{\left(  n-1,1\right)  }\subseteq\operatorname*{span}%
W\subseteq\operatorname*{span}W^{\prime}$, so that $\mathcal{S}^{\left(
n-1,1\right)  }=\operatorname*{span}W^{\prime}$ (since $\operatorname*{span}%
W^{\prime}\subseteq\mathcal{S}^{\left(  n-1,1\right)  }$). In other words, the
$n-1$ vectors $\underline{\overline{1}}-\underline{\overline
{\operatorname*{avg}}},\ \underline{\overline{2}}-\underline{\overline
{\operatorname*{avg}}},\ \ldots,\ \underline{\overline{n-1}}%
-\underline{\overline{\operatorname*{avg}}}$ span the $\mathbf{k}$-module
$\mathcal{S}^{\left(  n-1,1\right)  }$ (since these $n-1$ vectors are the
elements of $W^{\prime}$).
\end{itemize}
\par
So we know that the $n-1$ vectors $\underline{\overline{1}}%
-\underline{\overline{\operatorname*{avg}}},\ \underline{\overline{2}%
}-\underline{\overline{\operatorname*{avg}}},\ \ldots,\ \underline{\overline
{n-1}}-\underline{\overline{\operatorname*{avg}}}$ are $\mathbf{k}$-linearly
independent and span the $\mathbf{k}$-module $\mathcal{S}^{\left(
n-1,1\right)  }$. Hence, they are a basis of $\mathcal{S}^{\left(
n-1,1\right)  }$.}). Now, the $\mathbf{k}$-linear map%
\begin{align*}
\varrho:\mathcal{N}_{n}  &  \rightarrow\mathcal{S}^{\left(  n-1,1\right)  },\\
e_{i}  &  \mapsto\underline{\overline{i}}-\underline{\overline
{\operatorname*{avg}}}\ \ \ \ \ \ \ \ \ \ \left(  \text{for all }i\in\left[
n\right]  \right)
\end{align*}
is well-defined (since each $i\in\left[  n\right]  $ satisfies
$\underline{\overline{i}}-\underline{\overline{\operatorname*{avg}}}%
\in\mathcal{S}^{\left(  n-1,1\right)  }$) and $S_{n}$-equivariant (since we
can easily see that $\sigma\cdot\underline{\overline{\operatorname*{avg}}%
}=\underline{\overline{\operatorname*{avg}}}$ for each $\sigma\in S_{n}$). It
furthermore sends $e_{1}+e_{2}+\cdots+e_{n}$ to $\left(  \underline{\overline
{1}}-\underline{\overline{\operatorname*{avg}}}\right)  +\left(
\underline{\overline{2}}-\underline{\overline{\operatorname*{avg}}}\right)
+\cdots+\left(  \underline{\overline{n}}-\underline{\overline
{\operatorname*{avg}}}\right)  =0$ (by (\ref{pf.prop.Z-iso.b.2})), and thus
vanishes on the submodule $\mathcal{D}_{n}$ (by linearity, since
$\mathcal{D}_{n}$ is spanned by $e_{1}+e_{2}+\cdots+e_{n}$). Hence, it factors
through an $S_{n}$-equivariant $\mathbf{k}$-linear map $\varrho^{\prime
}:\mathcal{N}_{n}/\mathcal{D}_{n}\rightarrow\mathcal{S}^{\left(  n-1,1\right)
}$. This latter map $\varrho^{\prime}$ sends each residue class $\overline
{e_{i}}$ to $\underline{\overline{i}}-\underline{\overline{\operatorname*{avg}%
}}$, and thus is invertible (since it sends the basis $\left(  \overline
{e_{1}},\overline{e_{2}},\ldots,\overline{e_{n-1}}\right)  $ of $\mathcal{N}%
_{n}/\mathcal{D}_{n}$ to the basis $\left(  \underline{\overline{1}%
}-\underline{\overline{\operatorname*{avg}}},\ \underline{\overline{2}%
}-\underline{\overline{\operatorname*{avg}}},\ \ldots,\ \underline{\overline
{n-1}}-\underline{\overline{\operatorname*{avg}}}\right)  $ of $\mathcal{S}%
^{\left(  n-1,1\right)  }$). Hence, $\varrho^{\prime}$ is an isomorphism of
$S_{n}$-representations (since it is $S_{n}$-equivariant). Therefore,
$\mathcal{S}^{\left(  n-1,1\right)  }\cong\mathcal{N}_{n}/\mathcal{D}%
_{n}=\mathcal{Z}_{n}$ as $S_{n}$-representations. Proposition \ref{prop.Z-iso}
\textbf{(b)} is now proved again.
\end{proof}

\subsection{\label{subsec.apx.pf.prop.indprod.ass}Proof of Proposition
\ref{prop.indprod.ass}}

\begin{proof}
[Proof of Proposition \ref{prop.indprod.ass}.]We shall use the following
general facts about induced representations (\cite[Exercise 4.1.2]{GriRei} and
\cite[Exercise 4.1.3]{GriRei}, respectively\footnote{The facts are stated in
\cite{GriRei} only for $\mathbf{k}=\mathbb{C}$, but the proofs work equally
well for any $\mathbf{k}$.}):

\begin{itemize}
\item \textit{Transitivity of induction:} Let $G$ be a group. Let $H$ be a
subgroup of $G$. Let $I$ be a subgroup of $H$. Let $U$ be a representation of
$I$. Then,
\begin{equation}
\operatorname*{Ind}\nolimits_{H}^{G}\operatorname*{Ind}\nolimits_{I}^{H}%
U\cong\operatorname*{Ind}\nolimits_{I}^{G}U. \label{pf.prop.indprod.ass.trind}%
\end{equation}


\item \textit{Monoidality of induction:} Let $G_{1}$ and $G_{2}$ be two
groups. Let $H_{1}$ be a subgroup of $G_{1}$, and let $H_{2}$ be a subgroup of
$G_{2}$. Let $W_{1}$ be a representation of $H_{1}$, and let $W_{2}$ be a
representation of $H_{2}$. Then,%
\begin{equation}
\operatorname*{Ind}\nolimits_{H_{1}\times H_{2}}^{G_{1}\times G_{2}}\left(
W_{1}\otimes W_{2}\right)  \cong\left(  \operatorname*{Ind}\nolimits_{H_{1}%
}^{G_{1}}W_{1}\right)  \otimes\left(  \operatorname*{Ind}\nolimits_{H_{2}%
}^{G_{2}}W_{2}\right)  . \label{pf.prop.indprod.ass.mind}%
\end{equation}

\end{itemize}

Now, set $G_{1}:=S_{n_{1}+n_{2}+\cdots+n_{i}}$ and $H_{1}:=S_{n_{1}}\times
S_{n_{2}}\times\cdots\times S_{n_{i}}$ and $W_{1}:=U_{1}\otimes U_{2}%
\otimes\cdots\otimes U_{i}$ (a representation of $H_{1}$) and $G_{2}%
:=S_{n_{i+1}+n_{i+2}+\cdots+n_{k}}$ and $H_{2}:=S_{n_{i+1}}\times S_{n_{i+2}%
}\times\cdots\times S_{n_{k}}$ and $W_{2}:=U_{i+1}\otimes U_{i+2}\otimes
\cdots\otimes U_{k}$ (a representation of $H_{2}$). Then, by the definition of
induction products, we have%
\[
U_{1}\ast U_{2}\ast\cdots\ast U_{i}=\operatorname*{Ind}\nolimits_{S_{n_{1}%
}\times S_{n_{2}}\times\cdots\times S_{n_{i}}}^{S_{n_{1}+n_{2}+\cdots+n_{i}}%
}\left(  U_{1}\otimes U_{2}\otimes\cdots\otimes U_{i}\right)
=\operatorname*{Ind}\nolimits_{H_{1}}^{G_{1}}W_{1}%
\]
(by the definitions of $G_{1}$ and $H_{1}$ and $W_{1}$) and likewise%
\[
U_{i+1}\ast U_{i+2}\ast\cdots\ast U_{k}=\operatorname*{Ind}\nolimits_{H_{2}%
}^{G_{2}}W_{2}.
\]
Thus,
\begin{align*}
&  \underbrace{\left(  U_{1}\ast U_{2}\ast\cdots\ast U_{i}\right)
}_{=\operatorname*{Ind}\nolimits_{H_{1}}^{G_{1}}W_{1}}\ast\underbrace{\left(
U_{i+1}\ast U_{i+2}\ast\cdots\ast U_{k}\right)  }_{=\operatorname*{Ind}%
\nolimits_{H_{2}}^{G_{2}}W_{2}}\\
&  =\left(  \operatorname*{Ind}\nolimits_{H_{1}}^{G_{1}}W_{1}\right)
\ast\left(  \operatorname*{Ind}\nolimits_{H_{2}}^{G_{2}}W_{2}\right) \\
&  =\operatorname*{Ind}\nolimits_{S_{n_{1}+n_{2}+\cdots+n_{i}}\times
S_{n_{i+1}+n_{i+2}+\cdots+n_{k}}}^{S_{n_{1}+n_{2}+\cdots+n_{k}}}\left(
\left(  \operatorname*{Ind}\nolimits_{H_{1}}^{G_{1}}W_{1}\right)
\otimes\left(  \operatorname*{Ind}\nolimits_{H_{2}}^{G_{2}}W_{2}\right)
\right) \\
&  \ \ \ \ \ \ \ \ \ \ \ \ \ \ \ \ \ \ \ \ \left(  \text{by the definition of
an induction product}\right) \\
&  =\operatorname*{Ind}\nolimits_{G_{1}\times G_{2}}^{S_{n_{1}+n_{2}%
+\cdots+n_{k}}}\underbrace{\left(  \left(  \operatorname*{Ind}\nolimits_{H_{1}%
}^{G_{1}}W_{1}\right)  \otimes\left(  \operatorname*{Ind}\nolimits_{H_{2}%
}^{G_{2}}W_{2}\right)  \right)  }_{\substack{\cong\operatorname*{Ind}%
\nolimits_{H_{1}\times H_{2}}^{G_{1}\times G_{2}}\left(  W_{1}\otimes
W_{2}\right)  \\\text{(by (\ref{pf.prop.indprod.ass.mind}))}}%
}\ \ \ \ \ \ \ \ \ \ \left(
\begin{array}
[c]{c}%
\text{by the definitions}\\
\text{of }G_{1}\text{ and }G_{2}%
\end{array}
\right) \\
&  \cong\operatorname*{Ind}\nolimits_{G_{1}\times G_{2}}^{S_{n_{1}%
+n_{2}+\cdots+n_{k}}}\left(  \operatorname*{Ind}\nolimits_{H_{1}\times H_{2}%
}^{G_{1}\times G_{2}}\left(  W_{1}\otimes W_{2}\right)  \right) \\
&  \cong\operatorname*{Ind}\nolimits_{H_{1}\times H_{2}}^{S_{n_{1}%
+n_{2}+\cdots+n_{k}}}\left(  W_{1}\otimes W_{2}\right)
\ \ \ \ \ \ \ \ \ \ \left(  \text{by an application of
(\ref{pf.prop.indprod.ass.trind})}\right) \\
&  \cong\operatorname*{Ind}\nolimits_{S_{n_{1}}\times S_{n_{2}}\times
\cdots\times S_{n_{k}}}^{S_{n_{1}+n_{2}+\cdots+n_{k}}}\left(  U_{1}\otimes
U_{2}\otimes\cdots\otimes U_{k}\right)
\end{align*}
(since the definitions of $H_{1}$ and $H_{2}$ yield%
\begin{align*}
H_{1}\times H_{2}  &  =\left(  S_{n_{1}}\times S_{n_{2}}\times\cdots\times
S_{n_{i}}\right)  \times\left(  S_{n_{i+1}}\times S_{n_{i+2}}\times
\cdots\times S_{n_{k}}\right) \\
&  \cong S_{n_{1}}\times S_{n_{2}}\times\cdots\times S_{n_{k}}%
\end{align*}
and%
\begin{align*}
W_{1}\otimes W_{2}  &  =\left(  U_{1}\otimes U_{2}\otimes\cdots\otimes
U_{i}\right)  \otimes\left(  U_{i+1}\otimes U_{i+2}\otimes\cdots\otimes
U_{k}\right) \\
&  \cong U_{1}\otimes U_{2}\otimes\cdots\otimes U_{k},
\end{align*}
and both of these isomorphisms are canonical and \textquotedblleft fit
together\textquotedblright\ in that the latter isomorphism respects the action
of the former groups\footnote{Alternatively, you can argue directly: Both
$\operatorname*{Ind}\nolimits_{H_{1}\times H_{2}}^{S_{n_{1}+n_{2}+\cdots
+n_{k}}}\left(  W_{1}\otimes W_{2}\right)  $ and $\operatorname*{Ind}%
\nolimits_{S_{n_{1}}\times S_{n_{2}}\times\cdots\times S_{n_{k}}}%
^{S_{n_{1}+n_{2}+\cdots+n_{k}}}\left(  U_{1}\otimes U_{2}\otimes\cdots\otimes
U_{k}\right)  $ can be written as%
\[
\mathbf{k}\left[  S_{n_{1}+n_{2}+\cdots+n_{k}}\right]  \otimes\left(
U_{1}\otimes U_{2}\otimes\cdots\otimes U_{k}\right)
\]
modulo the relation
\[
\sigma\left(  \sigma_{1}\ast\sigma_{2}\ast\cdots\ast\sigma_{k}\right)
\otimes\left(  u_{1}\otimes u_{2}\otimes\cdots\otimes u_{k}\right)
-\sigma\otimes\left(  \sigma_{1}u_{1}\otimes\sigma_{2}u_{2}\otimes
\cdots\otimes\sigma_{k}u_{k}\right)
\]
for all $\sigma\in S_{n_{1}+n_{2}+\cdots+n_{k}}$ and $\sigma_{i}\in S_{n_{i}}$
and $u_{i}\in U_{i}$.}). In view of%
\[
U_{1}\ast U_{2}\ast\cdots\ast U_{k}=\operatorname*{Ind}\nolimits_{S_{n_{1}%
}\times S_{n_{2}}\times\cdots\times S_{n_{k}}}^{S_{n_{1}+n_{2}+\cdots+n_{k}}%
}\left(  U_{1}\otimes U_{2}\otimes\cdots\otimes U_{k}\right)
\]
(by the definition of the induction product), we can rewrite this as%
\[
\left(  U_{1}\ast U_{2}\ast\cdots\ast U_{i}\right)  \ast\left(  U_{i+1}\ast
U_{i+2}\ast\cdots\ast U_{k}\right)  \cong U_{1}\ast U_{2}\ast\cdots\ast
U_{k}.
\]
This proves Proposition \ref{prop.indprod.ass}.
\end{proof}

\subsection{\label{subsec.apx.proofs.specht.tensA}Proof of Proposition
\ref{prop.specht.tensA}}

\begin{proof}
[Proof of Proposition \ref{prop.specht.tensA}.]The following proof works more
generally when $\mathbf{k}$ is a commutative ring in which $n!$ is invertible.

We shall use the notations of \cite[Chapter 5]{sga}. Pick any $n$-tableau $P$
of shape $\lambda$ and any $n$-tableau $Q$ of shape $\mu$. Consider the
corresponding Young symmetrizers $\mathbf{E}_{P}$ and $\mathbf{E}_{Q}$ (as
defined in \cite[Definition 5.11.1]{sga}). From \cite[(228)]{sga}, we have
$\mathcal{S}^{\lambda}\cong\mathcal{A}\mathbf{E}_{P}$ and $\mathcal{S}^{\mu
}\cong\mathcal{A}\mathbf{E}_{Q}$. Hence, $\operatorname*{Hom}%
\nolimits_{\mathcal{A}}\left(  \mathcal{S}^{\lambda},\mathcal{S}^{\mu}\right)
\cong\operatorname*{Hom}\nolimits_{\mathcal{A}}\left(  \mathcal{A}%
\mathbf{E}_{P},\mathcal{A}\mathbf{E}_{Q}\right)  $ (since $\operatorname*{Hom}%
\nolimits_{\mathcal{A}}$ is functorial).

From \cite[Lemma 5.11.13]{sga}, we know that the coefficient of the
permutation $\operatorname*{id}\in S_{n}$ in $\mathbf{E}_{P}$ is $1$. Thus,
the single element $\mathbf{E}_{P}$ is $\mathbf{k}$-linearly independent.
Hence, $\mathbf{kE}_{P}\cong\mathbf{k}$ as $\mathbf{k}$-module.

However, \cite[Theorem 5.11.3]{sga} shows that $\mathbf{E}_{P}^{2}=\dfrac
{n!}{f^{\lambda}}\mathbf{E}_{P}$. This shows that the element
$\widetilde{\mathbf{E}}_{P}:=\dfrac{f^{\lambda}}{n!}\mathbf{E}_{P}%
\in\mathcal{A}$ is idempotent\footnote{\textit{Proof.} From
$\widetilde{\mathbf{E}}_{P}=\dfrac{f^{\lambda}}{n!}\mathbf{E}_{P}$, we obtain
$\widetilde{\mathbf{E}}_{P}^{2}=\left(  \dfrac{f^{\lambda}}{n!}\mathbf{E}%
_{P}\right)  ^{2}=\left(  \dfrac{f^{\lambda}}{n!}\right)  ^{2}%
\underbrace{\mathbf{E}_{P}^{2}}_{=\dfrac{n!}{f^{\lambda}}\mathbf{E}_{P}%
}=\left(  \dfrac{f^{\lambda}}{n!}\right)  ^{2}\dfrac{n!}{f^{\lambda}%
}\mathbf{E}_{P}=\dfrac{f^{\lambda}}{n!}\mathbf{E}_{P}=\widetilde{\mathbf{E}%
}_{P}$. In other words, $\widetilde{\mathbf{E}}_{P}$ is idempotent.}.
Moreover, $\widetilde{\mathbf{E}}_{P}$ is a nonzero scalar multiple of
$\mathbf{E}_{P}$ (since the scalar $\dfrac{f^{\lambda}}{n!}$ is clearly
nonzero). Thus, $\mathcal{A}\mathbf{E}_{P}=\mathcal{A}\widetilde{\mathbf{E}%
}_{P}$.

But \cite[Lemma 5.13.4]{Etingof-et-al} says that if $A$ is a $\mathbf{k}%
$-algebra and $e\in A$ is an idempotent, then $\operatorname*{Hom}%
\nolimits_{A}\left(  Ae,M\right)  \cong eM$ for any left $A$-module $M$.
Applying this to $A=\mathcal{A}$ and $e=\widetilde{\mathbf{E}}_{P}$ and
$M=\mathcal{A}\mathbf{E}_{Q}$, we find%
\[
\operatorname*{Hom}\nolimits_{\mathcal{A}}\left(  \mathcal{A}%
\widetilde{\mathbf{E}}_{P},\mathcal{A}\mathbf{E}_{Q}\right)  \cong%
\widetilde{\mathbf{E}}_{P}\mathcal{A}\mathbf{E}_{Q}=\mathbf{E}_{P}%
\mathcal{A}\mathbf{E}_{Q}%
\]
(since $\widetilde{\mathbf{E}}_{P}$ is a nonzero scalar multiple of
$\mathbf{E}_{P}$). Altogether,%
\begin{align}
\operatorname*{Hom}\nolimits_{\mathcal{A}}\left(  \mathcal{S}^{\lambda
},\mathcal{S}^{\mu}\right)   &  \cong\operatorname*{Hom}\nolimits_{\mathcal{A}%
}\left(  \underbrace{\mathcal{A}\mathbf{E}_{P}}_{=\mathcal{A}%
\widetilde{\mathbf{E}}_{P}},\mathcal{A}\mathbf{E}_{Q}\right)
=\operatorname*{Hom}\nolimits_{\mathcal{A}}\left(  \mathcal{A}%
\widetilde{\mathbf{E}}_{P},\mathcal{A}\mathbf{E}_{Q}\right) \nonumber\\
&  \cong\mathbf{E}_{P}\mathcal{A}\mathbf{E}_{Q}.
\label{pf.prop.specht.tensA.all}%
\end{align}


Now, if $\lambda\neq\mu$, then \cite[Proposition 5.11.15]{sga} (applied to
$S=P$ and $T=Q$) shows that $\mathbf{E}_{P}\mathbf{aE}_{Q}=0$ for all
$\mathbf{a}\in\mathcal{A}$, and thus $\mathbf{E}_{P}\mathcal{A}\mathbf{E}%
_{Q}=0$. Hence, if $\lambda\neq\mu$, then (\ref{pf.prop.specht.tensA.all})
becomes%
\begin{equation}
\operatorname*{Hom}\nolimits_{\mathcal{A}}\left(  \mathcal{S}^{\lambda
},\mathcal{S}^{\mu}\right)  \cong\mathbf{E}_{P}\mathcal{A}\mathbf{E}_{Q}=0.
\label{pf.prop.specht.tensA.distinct}%
\end{equation}


This proves Proposition \ref{prop.specht.tensA} in the case when $\lambda
\neq\mu$.

Thus, we now WLOG assume that $\lambda=\mu$. Hence, $\mu=\lambda$ and thus
$\mathcal{S}^{\mu}=\mathcal{S}^{\lambda}\cong\mathcal{A}\mathbf{E}_{P}$.
Hence, the same argument that we used to prove (\ref{pf.prop.specht.tensA.all}%
) can be applied to $\lambda$ and $P$ instead of $\mu$ and $Q$. This results
in%
\begin{equation}
\operatorname*{Hom}\nolimits_{\mathcal{A}}\left(  \mathcal{S}^{\lambda
},\mathcal{S}^{\lambda}\right)  \cong\mathbf{E}_{P}\mathcal{A}\mathbf{E}_{P}.
\label{pf.prop.specht.tensA.7}%
\end{equation}
However, \cite[Proposition 5.11.5]{sga} (applied to $T=P$) shows that each
$\mathbf{a}\in\mathcal{A}$ satisfies $\mathbf{E}_{P}\mathbf{aE}_{P}%
=\kappa\mathbf{E}_{P}$ for some $\kappa\in\mathbf{k}$. In other words, each
$\mathbf{a}\in\mathcal{A}$ satisfies $\mathbf{E}_{P}\mathbf{aE}_{P}%
\in\mathbf{kE}_{P}$. In other words, $\mathbf{E}_{P}\mathcal{A}\mathbf{E}%
_{P}\subseteq\mathbf{kE}_{P}$. On the other hand, each $\kappa\in\mathbf{k}$
satisfies
\[
\kappa\mathbf{E}_{P}=\dfrac{\kappa f^{\lambda}}{n!}\cdot\underbrace{\dfrac
{n!}{f^{\lambda}}\mathbf{E}_{P}}_{\substack{=\mathbf{E}_{P}^{2}=\mathbf{E}%
_{P}1\mathbf{E}_{P}\\\in\mathbf{E}_{P}\mathcal{A}\mathbf{E}_{P}}}\in
\dfrac{\kappa f^{\lambda}}{n!}\mathbf{E}_{P}\mathcal{A}\mathbf{E}_{P}%
\subseteq\mathbf{E}_{P}\mathcal{A}\mathbf{E}_{P}.
\]
Thus, $\mathbf{kE}_{P}\subseteq\mathbf{E}_{P}\mathcal{A}\mathbf{E}_{P}$.
Combining this with $\mathbf{E}_{P}\mathcal{A}\mathbf{E}_{P}\subseteq
\mathbf{kE}_{P}$, we obtain%
\[
\mathbf{E}_{P}\mathcal{A}\mathbf{E}_{P}=\mathbf{kE}_{P}.
\]
Hence, (\ref{pf.prop.specht.tensA.7}) becomes%
\[
\operatorname*{Hom}\nolimits_{\mathcal{A}}\left(  \mathcal{S}^{\lambda
},\mathcal{S}^{\lambda}\right)  \cong\mathbf{E}_{P}\mathcal{A}\mathbf{E}%
_{P}=\mathbf{kE}_{P}\cong\mathbf{k}.
\]
In other words, $\operatorname*{Hom}\nolimits_{\mathcal{A}}\left(
\mathcal{S}^{\lambda},\mathcal{S}^{\mu}\right)  \cong\mathbf{k}$ if
$\lambda=\mu$. Thus, Proposition \ref{prop.specht.tensA} is proved in the case
$\lambda=\mu$ as well.
\end{proof}

\begin{noncompile}
\begin{proof}
[Old proof of Proposition \ref{prop.specht.tensA}.]Recall (e.g., from
\cite[Theorem 5.12.2 and Corollary 5.12.4]{Etingof-et-al}\footnote{We are
using the fact that the Specht module $\mathcal{S}^{\lambda}$ is isomorphic to
the $V_{\lambda}$ in \cite{Etingof-et-al}. This holds by \cite[Theorem 5.5.13
\textbf{(b)}]{sga}.}) that the Specht modules $\mathcal{S}^{\nu}$ for the
partitions $\nu$ of $n$ are pairwise non-isomorphic split irreducible $S_{n}%
$-representations. Hence, the hom-space $\operatorname*{Hom}%
\nolimits_{\mathcal{A}}\left(  \mathcal{S}^{\lambda},\mathcal{S}^{\mu}\right)
$ is given by%
\[
\operatorname*{Hom}\nolimits_{\mathcal{A}}\left(  \mathcal{S}^{\lambda
},\mathcal{S}^{\mu}\right)  \cong%
\begin{cases}
\mathbf{k}, & \text{if }\lambda=\mu;\\
0, & \text{if }\lambda\neq\mu
\end{cases}
\ \ \ \ \ \ \ \ \ \ \text{as }\mathbf{k}\text{-vector spaces.}%
\]
This proves Proposition \ref{prop.specht.tensA}.
\end{proof}
\end{noncompile}


\newcommand{\etalchar}[1]{$^{#1}$}
\begin{thebibliography}{AFBC{\etalchar{+}}24}
\expandafter\ifx\csname url\endcsname\relax
  \def\url#1{\texttt{#1}}\fi
\expandafter\ifx\csname doi\endcsname\relax
  \def\doi#1{\burlalt{doi:#1}{https://doi.org/#1}}\fi
\expandafter\ifx\csname urlprefix\endcsname\relax\def\urlprefix{URL }\fi
\expandafter\ifx\csname href\endcsname\relax
  \def\href#1#2{#2}\fi
\expandafter\ifx\csname burlalt\endcsname\relax
  \def\burlalt#1#2{\href{#2}{#1}}\fi

\bibitem[AFBC{\etalchar{+}}24]{2407.08644}
I.~Axelrod-Freed, S.~Brauner, J.~H.-H. Chiang, P.~Commins, and V.~Lang.
\newblock Spectrum of random-to-random shuffling in the hecke algebra, 2024,
  \burlalt{2407.08644}{http://arxiv.org/abs/2407.08644}.
\newblock \urlprefix\url{https://arxiv.org/abs/2407.08644}.

\bibitem[Cla91]{Clausen91}
M.~Clausen.
\newblock Multivariate polynomials, standard tableaux, and representations of
  symmetric groups.
\newblock {\em J. Symb. Comput.}, 11(5-6):483--522, 1991.
\newblock \doi{10.1016/S0747-7171(08)80117-4}.

\bibitem[Egg19]{Egge19}
E.~S. Egge.
\newblock {\em An introduction to symmetric functions and their combinatorics},
  volume~91 of {\em Stud. Math. Libr.}
\newblock Providence, RI: American Mathematical Society (AMS), 2019.
\newblock \doi{10.1090/stml/091}.

\bibitem[EGH{\etalchar{+}}11]{Etingof-et-al}
P.~Etingof, O.~Golberg, S.~Hensel, T.~Liu, A.~Schwendner, D.~Vaintrob, and
  E.~Yudovina.
\newblock {\em Introduction to representation theory. {With} historical
  interludes by {Slava} {Gerovitch}.}, volume~59 of {\em Stud. Math. Libr.}
\newblock Providence, RI: American Mathematical Society (AMS), 2011.

\bibitem[Fulton97]{Fulton97}
W.~Fulton.
\newblock {\em {Y}oung Tableaux}, volume~35 of {\em London Mathematical Society
  Student Texts}.
\newblock Cambridge University Press, 1997.
\newblock \doi{10.1017/CBO9780511626241}.
\newblock Errata at \url{https://mathoverflow.net/questions/456463/}.

\bibitem[GL24]{s2b1-pub}
D.~Grinberg and N.~Lafreni\`ere.
\newblock The one-sided cycle shuffles in the symmetric group algebra.
\newblock {\em Algebr. Comb.}, 7(2):275--326, 2024.

\bibitem[GP25]{GriPar25}
D.~Grinberg and J.~Parlett.
\newblock Top to random and reverse: analysis of a new descent algebra shuffle,
  2025, \burlalt{2508.06740v1}{http://arxiv.org/abs/2508.06740v1}.
\newblock \urlprefix\url{https://arxiv.org/abs/2508.06740v1}.

\bibitem[GR20]{GriRei}
D.~Grinberg and V.~Reiner.
\newblock Hopf {Algebras} in {Combinatorics}.
\newblock Preprint, {arXiv}:1409.8356v7 [math.{CO}] (2020), 2020.
\newblock \urlprefix\url{https://arxiv.org/abs/1409.8356v7}.

\bibitem[Gri25]{sga}
D.~Grinberg.
\newblock An introduction to the symmetric group algebra.
\newblock Preprint, {arXiv}:2507.20706v1 [math.{CO}] (2025), 2025.
\newblock \urlprefix\url{https://arxiv.org/abs/2507.20706v1}.

\bibitem[GriLaf22]{s2b1}
D.~Grinberg and N.~Lafreni\`ere.
\newblock The one-sided cycle shuffles in the symmetric group algebra.
\newblock arXiv:2212.06274v4, 2023.

\bibitem[GriLaf24]{fps2024sn}
D.~Grinberg and N.~Lafreni\`ere.
\newblock The somewhere-to-below shuffles in the symmetric group and {H}ecke
  algebras.
\newblock extended abstract at the FPSAC 2024 conference, 2023.

\bibitem[Grinbe16]{Grinberg-Orzech}
D.~Grinberg.
\newblock A constructive proof of {O}rzech's theorem.
\newblock
  \urlprefix\url{https://www.cip.ifi.lmu.de/~grinberg/algebra/orzech.pdf}.

\bibitem[Grinbe23]{s2b2}
D.~Grinberg.
\newblock Commutator nilpotency for somewhere-to-below shuffles.
\newblock arXiv:2309.05340v2, 2023.

\bibitem[Lor18]{Lorenz18}
M.~Lorenz.
\newblock {\em A tour of representation theory}, volume 193 of {\em Grad. Stud.
  Math.}
\newblock Providence, RI: American Mathematical Society (AMS), 2018.
\newblock \doi{10.1090/gsm/193}.

\bibitem[RW84]{RemmelWhitney}
J.~B. Remmel and R.~Whitney.
\newblock Multiplying {Schur} functions.
\newblock {\em J. Algorithms}, 5:471--487, 1984.
\newblock \doi{10.1016/0196-6774(84)90002-6}.

\bibitem[Sag01]{Sagan-SG}
B.~E. Sagan.
\newblock {\em The symmetric group. {Representations}, combinatorial
  algorithms, and symmetric functions.}, volume 203 of {\em Grad. Texts Math.}
\newblock New York, NY: Springer, 2nd ed. edition, 2001.

\bibitem[Sta24]{Stanley-EC2}
R.~P. Stanley.
\newblock {\em Enumerative combinatorics. {Volume} 2. {With} an appendix by
  {Sergey} {Fomin}}, volume 208 of {\em Camb. Stud. Adv. Math.}
\newblock Cambridge: Cambridge University Press, 2nd edition edition, 2024.
\newblock \doi{10.1017/9781009262538}.

\end{thebibliography}

\end{document}
